\documentclass[11pt]{article}

\usepackage{maiacustom}

\begin{document}

\psettitle{Astronomy Question Bank}

These questions were carefully written or selected to prepare students for the astronomy team selection process in Brazil. Some questions are not original; those are marked with () before the statement. The question-bank template is the same one used by Professor Kevin Zhou. His work is invaluable, and many ideas in this set can be found in his handouts.

% \begin{psidea}{Idea title}{}
% Idea
% \end{psidea}

% \begin{psexample}{Example title}{}
% Example
% \end{psexample}

% \begin{pssolution*}{}{}
% Solution
% \end{pssolution*}

% \begin{psremark*}{Remark title}{}
% Remark
% \end{psremark*}

\section{Celestial Mechanics}
\pts{5}    
\begin{pproblem}
    Iaum and Sevla, two renowned physicists, intend to launch a space probe to explore the limits of Newtonian physics and the relativistic corrections required near a black hole. They need your help to study the \textit{effective potential} of such a body.
    \begin{enumerate}[label=\textbf{\alph*)}]
        \item The energy of a body of mass \(m\) orbiting a massive body of mass \(M\) can be written as:
             \[
             \frac{m\dot r ^2}{2} + V_{eff}(r)
             \]
             where \(\dot r\) is the radial velocity of the body. Find \(V_{eff}(r)\) in terms of \(G\), \(M\), \(L\), \(m\), and \(r\).

        \item According to Newtonian physics, what is the smallest radius at which a body can have a circular orbit around a black hole?
    
        \item What is the frequency of small radial oscillations, \(\omega_r\), about this radius?

        The gravitational potential near these bodies must include a relativistic correction and has the form:
        \[
        V(r) = -\frac{GMm}{r} - \frac{GML^2}{mc^2r^3}
        \]

        \item Explain why this correction allows a particle to fall into the center of a black hole, \(r=0\), and why this is impossible in Newtonian physics.
        
        \item What is the smallest value of \(L\) for which a particle can orbit the black hole on a circular orbit? What is the value of the radius in this case?
        
        \item Assuming \( \lim_{L \to \infty}\), find the smallest distance to which a particle on an open orbit can approach a black hole.  
    \end{enumerate}
    \begin{pssolution*}{}{}
        \begin{alternativas}
            \item In the classical form, we have:
            \[E = \frac{mv^2}{2}-\frac{GMm}{r}\]
            Note that we can write \(v^2 = v_\theta^2+ v_r^2\), where \(v_\theta\) and \(v_r\) are, respectively, the tangential and radial velocities of the body. Since \(v_r = \dot{r}\), we need an expression that relates \(v_\theta\) to another quantity. The best way to make this relation is by using the angular momentum, since:
            \[L = mrv_\theta \rightarrow v_\theta = \frac{L}{mr}\]
            Substituting into the energy expression, we have:
            \[E = \frac{m\dot{r}^2}{2}+\frac{L^2}{2mr^2} - \frac{GMm}{r}\]
            Since we want an answer of the form \(E = m\dot{r}^2/2 + V_{eff}(r)\), by analogy we have:
            \[\boxed{V_{eff}(r) = -\frac{GMm}{r}+\frac{L^2}{2mr^2}}\]

            \item The greatest limiting speed of a body is \(c\), disregarding relativistic effects:
            \[\frac{mc^2}{2}-\frac{GMm}{R} = -\frac{GMm}{2R}\]
            Isolating \(R\):
            \[\boxed{R = \frac{GM}{c^2}}\]

            \item For small oscillations, we have the approximation:
            \[\omega_r \approx \sqrt{\frac{V_{eff}''(R)}{m}}\]
            Calculating:
            \[V_{eff}'(R) = \frac{GMm}{r^2}-\frac{L^2}{mr^3}\]
            \[V_{eff}''(R) = -\frac{2GMm}{r^3}+\frac{3L^2}{mr^4}\]
            Using \(L = mvr\) (circular orbits) and \(v\equiv c\), we arrive at:
            \[V_{eff}''(R) = \frac{mc^6}{G^2M^2}\]
            With this, we can conclude that:
            \[\omega_r = \frac{c^3}{GM}\]

            \item When we take the limit \(\lim_{r\rightarrow 0}V_{eff}(r)\), in Newtonian physics the dominant term is \(L^2/2mr^2\), that is, an energy of \(+\infty\) would be required to reach the center of the black hole. However, after the relativistic correction, the dominant term is \(-GML^2/mc^2r^3\), that is, an energy of \(-\infty\) would be required to reach the center of the black hole. That is why it is not possible to reach the center of a black hole in Newtonian physics.
            
            \item On a circular orbit, \(r\) is constant, which implies that the force in the radial direction vanishes, that is:
            \[V_{eff}'(r) = 0\]
            Calculating:
            \[V_{eff}'(r) = \frac{GMm}{r^2}-\frac{L^2}{mr^3}+\frac{3GML^2}{mc^2r^4}=0\]
            Simplifying, we have:
            \[r^2-\frac{L^2}{GMm^2}r+\frac{3L^2}{m^2c^2}=0\]
            Note that the solution of this equation must be real, since \(r\) cannot be imaginary, which imposes the condition:
            \[\frac{L^4}{G^2M^2m^4} - \frac{12L^2}{m^2c^2} \geq 0\]
            In the equality case, the value of \(L\) that satisfies this equation is:
            \[\boxed{L = \frac{\sqrt{12}GMm}{c}}\]
            The value of \(r\) in this case is:
            \[r = \frac{L^2}{2GMm^2} = \boxed{r_{min} = \frac{6GM}{c^2}}\] 
            
            \item The complete expression for \(r\) is:
            \[r = \frac{\frac{L^2}{GMm^2}\pm\sqrt{\left(\frac{L^2}{GMm^2}\right)^2-\frac{12L^2}{m^2c^2}}}{2} = \frac{L^2-\sqrt{L^4-12(GMmL/c)^2}}{2GMm^2}\] 
            In the limit \(L\rightarrow \infty\), we have:
            \[r_{min} = \lim_{L \rightarrow \infty}\frac{L^2-\sqrt{L^4-12(GMmL/c)^2}}{2GMm^2} = \frac{6(GMm/c)^2}{2(GMm^2)} = \boxed{\frac{3GM}{c^2}} \] 
        \end{alternativas}
    \end{pssolution*}
\end{pproblem}

\pts{3}
\begin{pproblem} An important point in the study of orbits is what happens when the potential follows something unusual. The goal of this question is to study that phenomenon. Consider the angular momentum \(L\) and the mass of the body \(m\).
    \begin{alternativas}
        \item For a general orbit that has a potential \(V(r)\) of the form \(\psi r^k\), find the value \(r_0\) of a circular orbit for this potential.
        
        \item Find the frequency of small oscillations, \(\omega_r\), about this radius. 
        
        \item Let \(\omega_\theta\) be the angular frequency of the orbit. Find the value of the ratio \(\frac{\omega_r}{\omega_\theta}\). An orbit can be closed only if this ratio is an integer. Why? Find values of \(k\) for which the orbit is closed.
         
    \end{alternativas} 

    \begin{pssolution*}{}{}
        \begin{alternativas}
            \item Since on a circular orbit \(V_{eff}'(r) = 0\), we have:
            \[V_{eff}'(r) = -\frac{L^2}{mr^3}+k\psi r^{k-1} = 0\]
            \[k\psi r^{k-1} = \frac{L^2}{mr^3} \rightarrow \boxed{r_0 = \left(\frac{L^2}{k\psi m}\right)^{\frac{1}{k+2}}}\]
        
            \item Using \(\omega_r^2=V_{eff}''(r_0)/m\):
            \[V_{eff}''(r) = \frac{3L^2}{mr^4}+k(k-1)\psi r^{k-2}\]
            \[V_{eff}''(r) = r^{-4}\left(\frac{3L^2}{m}+k(k-1)\psi r^{k+2}\right)\]
            \[V_{eff}''(r) = r^{-4}\left(\frac{3L^2}{m}+k(k-1)\psi \frac{L^2}{k\psi m}\right)\]
            \[V_{eff}''(r) = \frac{L^2}{r_0^4m}(k+2)\]
            Thus, we have:
            \[\boxed{\omega_r = \frac{L^2}{r_0^2m}\sqrt{k+2}}\]

            \item By the definition of \(L\), we have:
            \[L = mr_0^2\omega_\theta \rightarrow \omega_\theta = \frac{L^2}{mr_0^2}\]
            Therefore:
            \[\frac{\omega_r}{\omega_\theta} = \sqrt{k+2}\]
            Note that, for the orbit to be a closed figure, this ratio must be an integer, so \(k+2\) must be a perfect square. Any value \(k = n^2-2 \ \forall \ n \in \mathbb{N}\) satisfies this equation.
        \end{alternativas}
    \end{pssolution*}
\end{pproblem}

\pts{4} %week 114
\begin{pproblem}
    The hodograph is one of the lesser-known ways of solving celestial-mechanics problems. Using this device, one can solve complex problems (involving, for example, finding orbital parameters, velocity parameters, and how to minimize the eccentricity of an orbit). Although the hodograph has many uses, the goal of this exercise is to prove the Hodograph Theorem, which states that, on closed orbits, the velocity forms a circle in vector space. We will prove this theorem in two ways.
    \begin{align*}
    \textbf{Method 1)}
    \end{align*}
    \begin{enumerate}[label=\textbf{\Roman*.}]
        \item Starting from Newton's second law and the conservation of angular momentum, find an expression for \(d\mathbf{v}/d\theta\). Leave your answer in terms of the angular momentum per unit mass, \(h\), and \(\mu\) = \(GM\). 
        \item Where is the center of this circle located? Taking the central body at the origin of the system, determine an expression for the distance between the central body and the center of the hodograph. 
    \end{enumerate}
    
    \begin{align*}
        \textbf{Method 2)}
    \end{align*}
    Method 2 uses more advanced notions of calculus and starts from the conservation of the \textit{Laplace-Runge-Lenz vector}, \(\mathbf{A}\).
    \begin{enumerate}[label=\textbf{\Roman*.}]
        \item The vector \(\mathbf{A}\) is given by the expression:
        \[
        \mathbf{A} = \mathbf{p}\times \mathbf{L} - GMm^2\mathbf{\hat{r}}
        \]
        The first step of the proof is to show that \(\mathbf{A}\) is constant. To do so, differentiate it with respect to time and draw your own conclusions.
        \item Since the vector \(\mathbf{L}\) is constant, \(\mathbf{A}\times\mathbf{L}\) is also a constant. Using the vector identity \(\mathbf{A}\times(\mathbf{B}\times\mathbf{C}) = \mathbf{B}(\mathbf{A}\cdot\mathbf{C}) - \mathbf{C}(\mathbf{A}\cdot\mathbf{B}) \), prove that \(\mathbf{v}\) forms a circle in vector space.    
    \end{enumerate}
    \begin{pssolution*}{}{}
        \textbf{Method 1)}

        \begin{enumerate}[label=\textbf{\Roman*.}]
            \item Newton's second law tells us that:
            \[-\frac{GMm}{r^2}\mathbf{\hat{r}} =  m\frac{d\mathbf{v}}{dt}\]
            By the definition of angular momentum:
            \[L = mr^2\frac{d\theta}{dt} \rightarrow \ \  \therefore dt = \frac{mr^2}{L}d\theta\]
            Substituting \(dt\) into Newton's second law:
            \[-\frac{GMm}{r^2}\mathbf{\hat{r}} = m\left(\frac{L}{mr^2}\right)\frac{d\mathbf{v}}{d\theta}\]

            Isolating \(dv/d\theta\), we have:

            \[\boxed{\frac{d\mathbf{v}}{d\theta} = -\frac{GMm}{L}\mathbf{\hat{r}} \equiv -\frac{\mu}{h}\mathbf{\hat{r}}}\]

            Note that this expression is a circle in vector space and has radius \(R = \mu/h\), since \(d\mathbf{v}/d\theta\) is constant and points in the radial direction.

            \item As we saw in the previous item, the hodograph forms a circle of radius \(R = \mu/h\). Use the following figure as support:
            \begin{figure}[H]
                \centering
                \includegraphics[width=0.4\linewidth]{imagens/hodografo1.png}
                \caption{Diagram of the hodograph}
            \end{figure}
            The point \(O\) represents the origin of the system (location of the central body). Note that the smallest distance between \(O\) and the hodograph must correspond to the smallest speed (apastron) and the largest distance must correspond to the largest speed (periastron). Also note that \(v_a + \overline{OC} = v_p - \overline{OC}\). Setting up the equations, we have:
            \[v_a + v_p = 2R \equiv \frac{2\mu}{h}\]
            \[v_p - v_a = 2\overline{OC}\]

            We can write \(v_p\) and \(v_a\) in terms of the angular momentum:
            \[L = ma(1-e)v_p \rightarrow v_p = \frac{h}{a(1-e)}\]
            \[L = ma(1+e)v_a \rightarrow v_a = \frac{h}{a(1+e)}\]
            
            Substituting into the first equation:
            \[\frac{h}{a}\left(\frac{1}{1+e}+\frac{1}{1-e}\right) = \frac{2\mu}{h}\]
            
            Here, we can isolate \(a\):
            \[\frac{h}{a}\left(\frac{2}{(1-e)^2}\right) = \frac{2\mu}{h}\]

            \[a = \frac{h^2}{\mu}\frac{1}{(1-e)^2}\]

            With this in mind, we can move on to the second expression:
            \[\frac{h}{a}\left(\frac{1}{1-e}-\frac{1}{1+e}\right) = 2\overline{OC}\]
            \[h\left(\frac{\mu(1-e)^2}{h^2}\right)\left(\frac{2e}{(1-e)^2}\right) = 2\overline{OC}\]
            
            Finally, solving for \(\overline{OC}\), we have:
            \[\boxed{\overline{OC} = \frac{\mu}{h}e}\]

            For an ellipse (\(e<1\)), the origin lies inside the circle; for a parabola (\(e=1\)), the origin lies on the boundary of the circle; and for a hyperbola (\(e>1\)), the origin lies outside the circle.
        \end{enumerate}

    
        \textbf{Method 2}
        \begin{enumerate}[label=\textbf{\Roman*.}]
            \item Differentiating \(\mathbf{A}\) with respect to time, we have:
            \[\frac{d\mathbf{A}}{dt} = \frac{d}{dt}(\mathbf{p}\times\mathbf{L}) - GMm^2\frac{d}{dt}\mathbf{\hat{r}}\] 
            When we differentiate a \textit{unit vector} with respect to time, we have:
            \[\frac{d\mathbf{\hat{r}}}{dt} = \mathbf{\omega}\times\mathbf{\hat{r}}\]
            where \(\mathbf{\omega}\) is the velocity with which the unit vector translates. Thus:
            \[\frac{d\mathbf{A}}{dt} = \left(\frac{d\mathbf{p}}{dt}\right)\times \mathbf{L} + \mathbf{p}\times\left(\frac{d\mathbf{L}}{dt}\right) - GMm^2(\mathbf{\omega}\times \mathbf{\hat{r}})\]
            Now we have a few considerations. First, note that \(d\mathbf{L}/dt=0\), since there are only central forces on the orbit. Second, \(d\mathbf{p}/dt = \mathbf{F}\). Furthermore, we have \(\mathbf{L} = mr^2\mathbf{\omega}\). Returning to the equation:

            \[\frac{d\mathbf{A}}{dt} = -\frac{GMm}{r^2}\mathbf{\hat{r}}\times(mr^2\mathbf{\omega}) - GMm^2(\mathbf{\omega} \times \mathbf{\hat{r}})\]
            \[\frac{d\mathbf{A}}{dt} = -GMm^2(\mathbf{\hat{r}}\times\mathbf{\omega})- GMm^2(\mathbf{\omega} \times \mathbf{\hat{r}})\]

            For any two vectors, \(\mathbf{i}\) and \(\mathbf{j}\), it holds that:
            \[i\times j = -j\times i\]

            Using this property, we can conclude that:

            \[\boxed{\frac{d\mathbf{A}}{dt} = 0}\]

            Therefore, \(\mathbf{A}\) is constant!

            \item Writing the expression and using \textit{BAC-CAB}, we have:
            
            \[\mathbf{A}\times\mathbf{L} = (\mathbf{p}\times\mathbf{L})\times\mathbf{L} - GMm^2\mathbf{\hat{r}}\times\mathbf{L}\]

            \[(\mathbf{A} +GMm^2\mathbf{\hat{r}})\times\mathbf{L} = - \mathbf{L}\times(\mathbf{p}\times\mathbf{L}) = \mathbf{p}(-\mathbf{L}\cdot\mathbf{L})-\mathbf{L}(-\mathbf{L}\cdot\mathbf{p})\]

            Since \(\mathbf{L}\) and \(\mathbf{p}\) are always perpendicular, \(\mathbf{L}\cdot\mathbf{p}=0\), so:

            \[(\mathbf{A}+GMm^2\mathbf{\hat{r}})\times\mathbf{L} = -mL^2\mathbf{v}\]

            Since both \(\mathbf{A}\) and \(\mathbf{\hat{r}}\) are perpendicular to \(\mathbf{L}\) and constant, \(\mathbf{A}+GMm^2\mathbf{\hat{r}}\) traces a circle. The cross product with \(\mathbf{L}\) merely rotates the circle by \(90^\circ\) and multiplies its radius by \(\mathbf{L}\). Therefore, \(\mathbf{v}\) traces a circle in vector space.
            
        \end{enumerate}
    \end{pssolution*}
\end{pproblem}

\pts{2} %problems of week 114
\begin{pproblem} One of the most fascinating phenomena of the Solar System is neutron stars. These are very small, extremely dense bodies that spin very fast. In this question, we will build a theoretical model for these stars.
    \begin{enumerate}[label=\textbf{\alph*)}]
        \item The Vela Pulsar, a neutron star located in the constellation Vela, has a frequency of \(11 \text{Hz}\), that is, it spins about its own axis \(11\) times per second, with equatorial radius \(R_e = 9.6\text{ km}\) and mass \(M = 1.88M_\odot\). Given this, calculate the ratio of its equatorial radius to its polar radius.
        \item What is the smallest rotation period the Vela Pulsar can have so that it does not break apart?
        \item Currently, the rate of change of the period of the Vela Pulsar is \(\frac{\Delta P}{\Delta t} = 1.25 \cdot 10^{-13}\text{s/s}\). This causes the pulsar to release an enormous amount of energy. The surface temperature of the pulsar is \(T \sim 10^6\text{ K}\). Calculate the order of magnitude of the ratio between the power emitted due to the increase of the period and the power emitted by the Stefan-Boltzmann law.
    \end{enumerate}
    \begin{pssolution*}{}{}
        \begin{alternativas}
            \item A star is always in hydrostatic equilibrium, so, writing the balance for a point on the equator and for one of the poles,
            \[-\frac{GM}{r_e}-\frac{1}{2}\omega^2r_e^2 = -\frac{GM}{r_p}\]

            Using \(\omega = 2\pi f\),

            \[\frac{GM}{r_p} = \frac{GM}{r_e}+2\pi^2f^2r_e^2\]

            \[\boxed{\frac{r_e}{r_p} = 1+\frac{2\pi^2f^2r_e^3}{GM} \approx 1 + 8.4\cdot10^{-6}}\]

            That is, a neutron star, even spinning absurdly fast, is still practically a perfect sphere!

            \item In the limiting situation in which it remains intact, we have:
            \[\frac{GM}{r_e^2} = \omega^2r_e\]
            \[\omega^2 = \frac{GM}{r_e^3} \rightarrow \boxed{T_{min} = \frac{2\pi r_e^3}{GM} = 2.2\cdot10^{-8}\text{ s}}\]

            \item The energy due to rotation has the form:
            \[E = \frac{1}{2}I\omega^2\]
            where \(I\) is the moment of inertia of the star (\(I = \frac{2}{5}Mr^2\)). Therefore, the dissipated power is:
            \[P = \frac{dE}{dt} = I\omega\frac{d\omega}{dt}\]

            Since \(\omega = 2\pi/T\),

            \[\frac{d\omega}{dT} = -\frac{2\pi}{T^2} \rightarrow d\omega = -\frac{2\pi}{T^2}dT\]
            
            The expression for the power as a function of the period becomes:

            \[P = -\frac{8\pi^2Mr_e^2}{5T^3}\frac{dT}{dt} = -9.1\cdot10^{29} \text{ W}\]

            Note that, by conservation of energy, all of this power must be converted into luminosity, so \(L_{rot} = -P = 9.1\cdot10^{29}\text{ W}\). The power coming from thermal radiation is given by:

            \[L_{rad} = 4\pi R^2\sigma T^4 = 6.56 \cdot 10^{25}\text{ W}\]

            With this, we can conclude that \(99.99\%\) of its luminosity comes from the rotation.
        \end{alternativas}
    \end{pssolution*}    
\end{pproblem}

\pts{3}
\begin{pproblem} (Adapted from PPP) Geométrio, Paulinho, and Hirata love chocolate. One day, they were exploring a galaxy when they came across a giant ball of chocolate. The three quickly began to eat the ball, first making a straight line from the point \(P\) to the center of the ball at \(O\) (diagram \(1\)). After that, Hirata falls, colliding at the point \(O\), without braking along the way. Surprisingly, he remains alive and is rescued by Geométrio and Paulinho. After that, still hungry, they continue to eat the giant chocolate sphere and leave a spherical hole of diameter \(\overline{PO}\) (diagram \(2\)). Hirata is very clumsy and ends up falling again, starting from the point \(P\) and going to \(O\) once more.
    \begin{enumerate}[label=\textbf{\alph*)}]
        \item Find the ratio of Hirata's impact speeds in cases \(1\) and \(2\).
        \item Find the ratio of Hirata's fall times in cases \(1\) and \(2\).
    \end{enumerate}
    \begin{paracol}{2}
        \begin{figure}[h]
            \centering
            \includegraphics[width=0.31\textwidth]{imagens/q5(1).png}
            \caption{Diagram \(1\)}
        \end{figure}
        \switchcolumn
        \begin{figure}[h]
            \centering
            \includegraphics[width=0.3\textwidth]{imagens/q5(2).png}
            \caption{Diagram \(2\)}
        \end{figure}
    \end{paracol} 

    \begin{pssolution*}{}{}
        \begin{alternativas}
            \item The gravity inside a planet is given by:
            \[\oint\mathbf{g}\cdot d\mathbf{S} = -4\pi GM_{int}\]
            \[4\pi r^2g = -4\pi G \rho \frac{4\pi}{3}r^3\]

            That is, for case \(1\), \(g \propto -\rho r \rightarrow g = -k\rho r\). Using Newton's second law:
            \[F = mv\frac{dv}{dr} = -mkr\]
            Integrating:
            \[\int_0^{v_1} v'dv' =-k\rho\int_0^R rdr\]

            \[\frac{v_1^2}{2} = k\rho\frac{R^2}{2} \rightarrow \boxed{v_1 = R \sqrt{k\rho}}\]

            Let \(\mathbf{j}\) be the vector from the center of the planet to the center of the cavity. Using superposition:
            \[a = -k\rho\mathbf{r} - k(-\rho)(\mathbf{r}-\mathbf{j}) = -k\rho\mathbf{j}\]
            That is, in case \(2\), the acceleration is constant and has magnitude \(|a| = \frac{k\rho R}{2}\).
            
            Using Torricelli:
            \[v_2^2 = 2aR k\rho R^2 \rightarrow \boxed{v_2 = R\sqrt{k\rho}}\]
            Therefore, \(v_1/v_2 = 1\).

            \item In the first case, note that we have simple harmonic motion, since the force is proportional to \(r\). The time to reach the bottom equals \(1/4\) of the period:
            \[\ddot{r} = -k\rho r \rightarrow \omega^2 = k\rho \rightarrow T = \frac{2\pi}{\sqrt{k\rho}} \rightarrow t_1 = \frac{T}{4} = \frac{\pi}{2\sqrt{k\rho}}\]

            In the second case, we have a constant acceleration, so:
            \[t_2 = \frac{v_2}{a_2} = \frac{2R\sqrt{k\rho}}{k\rho R} = \frac{2}{\sqrt{k\rho}}\]

            Therefore, the ratio \(t_1/t_2\) is given by:
            \[\boxed{\frac{t_1}{t_2} = \frac{\pi}{4}}\]

        \end{alternativas}
    \end{pssolution*}
\end{pproblem}

\pts{5}
\begin{pproblem} Calculating escape speeds in certain situations can be more complicated than it seems. One example is calculating the launch speed a rocket needs in order to reach a given location.
    \\
    \textbf{Hint:} The following problems require you to think carefully about the most appropriate reference frame. Do not underestimate the question.
    \\
    \textbf{Data:} A body can orbit the surface of the Earth with speed \(v_0 = 7.9\text{ km/s}\), the orbital speed of the Earth around the Sun is \(u_0 = 29.7\text{ km/s}\), and the radius of the Earth is negligible compared with the Earth-Sun distance. Also assume that, upon leaving the Earth's gravitational field, the distance between the probe and the Sun is the same as the distance between the Earth and the Sun.
    \begin{enumerate}[label=\textbf{\alph*)}]
        \item What is the smallest launch speed a rocket needs in order to reach the Sun, assuming that it applies only one impulse? (To check, \(v = 31.8\) km/s)
        \item What is the smallest launch speed a rocket needs in order to escape the Solar System? (To check, \(v = 16.7\) km/s)
        \item How does the answer to item \textbf{a)} change if the rocket can apply a second, very small impulse at some point of its orbit?
    \end{enumerate}

    \begin{pssolution*}{}{}
        \begin{alternativas}
            \item For this to happen, after leaving the Earth, in the Sun's frame the probe must be at rest and, consequently, in the Earth's frame it must have velocity \(-u_0\). Using conservation of energy:
            \[\frac{mv^2}{2}-\frac{GM_\oplus m}{R_\oplus} = \frac{mu_0^2}{2}\]
            \[v^2 = u_0^2+\frac{2GM_\oplus}{R_\oplus}\]

            Note that the speed \(v_0\) can also be obtained by equating the centripetal acceleration to the gravitational acceleration:
            \[\frac{mv_0^2}{R_\oplus} = \frac{GM_\oplus m}{R_\oplus^2} \rightarrow v_0^2 = \frac{GM_\oplus}{R_\oplus}\]

            Substituting,

            \[v^2=u_0^2+2v_0^2 \rightarrow \boxed{v = 31.8\text{ km/s}}\]

            \item In this case, upon leaving the Earth's gravitational field, in the Sun's frame the body must have speed \(\sqrt{2}u_0\), so in the Earth's frame it must have \((\sqrt{2}-1)u_0\). Again, by conservation of energy:
            
            \[\frac{mv^2}{2}-\frac{GM_\oplus m}{R_\oplus} = \frac{m(\sqrt{2}-1)^2u_0^2}{2}\]

            \[v^2 = (\sqrt{2}-1)^2u_0^2+2v_0^2 \rightarrow \boxed{v = 16.7 \text{ km/s}}\]

            \item Imagine that we send the probe to infinity, that is, we launch it with speed \(v=16.7\text{ km/s}\), and after that we apply an infinitesimal impulse toward the Sun. This makes the probe travel all the way back and, consequently, reach the Sun.
        \end{alternativas}
    \end{pssolution*}
\end{pproblem}

\pts{3} % problems of week 115
\begin{pproblem} (Kevin Zhou) The rocket equation is given by:
    \[
    v = u\ln\left(\frac{M_0}{M}\right)
    \]
    where \(v\) is the speed of the rocket, \(u\) the relative speed at which the rocket ejects fuel, \(M_0\) the initial mass of the rocket, and \(M\) its current mass.
    \begin{enumerate}[label=\textbf{\alph*)}]
        \item Derive the rocket equation starting from conservation of momentum.
        \item Assuming \(u\) is fixed throughout the rocket's trip, what must the value of \(u\) be so that the rocket goes from \(0\) to \(v\) using the least possible fuel? (You will need to solve numerically for \(v/u\))
        \item How must the rocket equation be corrected for a region of space with a constant gravitational field, \(\mathbf{g}\), pointing opposite the rocket's motion? Take \(\eta = dm/dt = \text{constant}\).
    \end{enumerate}

    \begin{pssolution*}{}{}
        \begin{alternativas}
            \item Consider that the rocket has mass \(M\) and velocity \(v\). It ejects fuel at a relative speed \(u\). Let \(dm\) be the mass of ejected gas:
            \[dp = (v-u)dm\]
            But, by definition,
            \[dp = vdm + mdv\]
            Combining these equations:
            \[mdv = -udm\]

            Integrating,
            \[\ln m = -\frac{v}{u} + C\]

            We know that when \(m = M_0\), we have \(v = 0\), so \(C = \ln M_0\). Thus:

            \[v = u \ln\left(\frac{M_0}{M}\right)\]

            As expected.

            \item The kinetic energy of the gas is given by:
            
            \[K_{g} =\frac{1}{2}M_gu^2 = \frac{(M_0-M)u^2}{2}\]

            We can find an expression for \(M\) using the rocket equation:

            \[\frac{M_0}{M} = e^{v/u} \rightarrow M_0 = Me^{v/u}\]

            Returning to the previous formula:
            
            \[K_g = \frac{M(e^{v/u}-1)u^2}{2}\]

            Here, \(M\) is equivalent to the ``structural'' mass of the rocket.

            To minimize the amount of energy spent on fuel, we have:

            \[\frac{dK_g}{du}=0  \rightarrow -e^{v/u}v +2u(e^{v/u}-1)=0\]

            Defining \(x=v/u\),

            \[e^xx = 2(e^x-1) \rightarrow x = 2(1-e^{-x})\]

            Solving by iteration, we find:

            \[x \approx 1.59 \rightarrow \boxed{u \approx \frac{v}{1.59}} \]

            \item In this case, we have:
            
            \[dp = (v-u)dm - mgdt = mdv +vdm\]
            \[\frac{dm}{m} = -\frac{dv}{u} - \frac{g}{u}dt\]

            Integrating,

            \[\ln\left(\frac{M}{M_0}\right) = -\frac{v}{u} - \frac{gt}{u}\]

            Since \(\eta\) is constant, \(M = M_0 -\eta t \rightarrow t = (M_0-M)/\eta\).

            \[\ln\left(\frac{M}{M_0}\right) = -\frac{v}{u} - \frac{g}{u\eta}(M_0-M)\]

            Solving for \(v\),

            \[\boxed{v = u\ln\left(\frac{M_0}{M}\right)-\frac{g}{\eta}(M_0-M)}\]
        \end{alternativas}
    \end{pssolution*}
\end{pproblem}  

\pts{4}
\begin{pproblem} One of the most common rocket propellants is a mixture of liquid hydrogen and liquid oxygen. When it begins to burn, the following chemical reaction occurs:
    \[
    2H_2 + O_2 \rightarrow 2H_2O
    \] 
    For each mole of hydrogen, this reaction releases \(241.8 \text{ kJ}\) of energy. Throughout the question, assume that all of this energy is used to move the rocket.
    \begin{enumerate}[label=\textbf{\alph*)}]
        \item A space expedition is to be carried out in such a way that a Hohmann transfer is required to launch a rocket from Earth to Mars. Calculate the total velocity change \(\Delta v\) needed to perform this maneuver.
        \item From the \(\Delta v\) calculated above, estimate how many tons of propellant must be used to perform this maneuver for a rocket that ejects propellant at speed \(u = 3.0\) km/s and has structural mass \(M = 140\cdot10^3\) kg.
    \end{enumerate}
    \begin{pssolution*}{}{}
        \begin{alternativas}
            \item For the first impulse, we have:
            \[v_{0,1}^2=\frac{GM}{R_T}\]
            \[v_{f,1}^2 = GM\left(\frac{2}{R_T}-\frac{1}{a}\right)\]
            The transfer orbit has \(a=(R_T+R_M)/2\),
            \[v_{f,1}^2 = \frac{2GM}{R_T}\left(\frac{R_M}{R_T+R_M}\right)\]

            For the second impulse,
            \[v_{0,2}^2 = GM\left(\frac{2}{R_M}-\frac{1}{a}\right) = \frac{2GM}{R_M}\left(\frac{R_T}{R_T+R_M}\right)\]

            \[v_{f,2}^2 = \frac{GM}{R_M}\]

            With \(\Delta v = \Delta v_1 + \Delta v_2 = (v_{f,1}-v_{0,1})+(v_{f,2}-v_{0,1})\), we can write:

            \[\boxed{\Delta v = \sqrt{\frac{2GM}{R_T}\left(\frac{R_M}{R_T+R_M}\right)}-\sqrt{\frac{GM}{R_T}} + \sqrt{\frac{GM}{R_M}} - \sqrt{\frac{2GM}{R_M}\left(\frac{R_T}{R_T+R_M}\right)}}\]
            
            Using \(M = M_\odot = 2\cdot10^{30}\) kg, \(R_M = 1.5\) AU, and \(R_T = 1\) AU, we find the numerical value \(\Delta v \approx 5.42 \) km/s.

            \item Using the rocket equation, we have:
            
            \[v = u\ln\left(\frac{M_0}{M}\right) \rightarrow \frac{M_0}{M}  = 1+ \frac{\mu}{M}= e^{v/u}\]

            where \(\mu\) is the mass of gas.

            \[\boxed{\mu = M(e^{v/u}-1) \approx 7.12 \cdot 10^5 \text{ kg}}\]

            This is consistent with reality, since the fuel mass is \(\approx 83\%\) of the total mass of the rocket.
        \end{alternativas}
    \end{pssolution*}
\end{pproblem}

\pts{3} % problems of week 115
\begin{pproblem} Marisso was studying a binary system with inclination \(i\). He found that the largest redshift coming from star \(1\) was \(z_1\). Given this, find an expression for the mass of star \(2\), \(m_2\). Leave your answer in terms of \(z_1\), the binary period \(P\), the inclination \(i\), the mass ratio \(\lambda  = m_2/m_1\), and fundamental constants.
    \begin{pssolution*}{}{}
        The speed at which \(1\) orbits the CM is:

        \[v_1 = \frac{2\pi a_1}{P}\]

        By geometry, the radial velocity is \(v_{r,1}=v_1\sin i\), so:

        \[z_1 \approx \frac{v_{r,1}}{c} = \frac{2\pi a_1\sin i}{Pc}\]

        From the center-of-mass theorem:

        \[m_1a_1=m_2a_2 \rightarrow a_2 = \frac{a_1}{\lambda}\]

        Using Kepler's third law:

        \[\frac{(a_1+a_2)^3}{P^2} = \frac{G(m_1+m_2)}{4\pi^2}\]

        \[\frac{a_1^3(1+1/\lambda)^3}{P^2} = \frac{Gm_2(1+1/\lambda)}{4\pi^2}\]

        Isolating \(m_2\):

        \[m_2 = \frac{4\pi^2a_1^3(1+1/\lambda)^2}{GP^2}\]

        \(a_1\) can be obtained by evaluating the redshift formula, \(a_1 = Pcz_1/2\pi\sin i\):

        \[m_2 = \frac{4\pi^2}{GP^2}\left(\frac{z_1Pc}{2\pi\sin i}\right)^3(1+1/\lambda)^2\]

        Finally, simplifying:

        \[\boxed{m_2 = \frac{z_1^3c^3P}{2\pi G \sin^3i}(1+1/\lambda)^2}\]
    \end{pssolution*}
\end{pproblem}

\pts{2}
\begin{pproblem} What is the ratio between \(a)\) the gravitational forces caused by the Sun and by the Moon on the surface of the Earth? And \(b)\) the tidal forces caused by the same?
    \begin{pssolution*}{}{}
        \begin{alternativas}
            \item The gravitational force has the form:
            \[F_G = \frac{GMm}{d^2}\]

            Thus,

            \[\frac{F_\odot}{F_L} = \frac{M_\odot}{M_L}\left(\frac{d_{T-L}}{d_{T-\odot}}\right)^2 \boxed{\approx 178.33}\]

            \item For the tidal forces, we have:
            
            \[F_M \propto \frac{m}{d^3}\]

            \[\therefore \frac{F_\odot}{F_L} = \frac{M_\odot}{M_L}\left(\frac{d_{T-L}}{d_{T-\odot}}\right)^3 \boxed{\approx 0.456}\]

        \end{alternativas}
    \end{pssolution*}
\end{pproblem}

\pts{4}
\begin{pproblem} (Morin 7.7) A particle of mass \(m\) travels on a hyperbolic orbit with a mass \(M\) fixed at one of the foci. The speed at infinity is \(v_0\) and the impact parameter is \(b\).
    \begin{alternativas}
    \item Show that the deflection angle of the particle is given by:
    \[
    \phi = \pi-2\tan^{-1}(\gamma b)
    \]
    where \(\gamma \equiv v_0^2/GM\).
    \item Let \(d\sigma\) be the cross section of the particle (measured when it is at infinity) that is deflected into a solid angle of size \(d\Omega\). Show that:
    \[
    \frac{d\sigma}{d\Omega} = \frac{1}{4\gamma^2\sin^4(\phi/2)}
    \] 
    As a point of curiosity, this quantity is called the \textit{differential cross section}.
    \end{alternativas}

    \begin{pssolution*}{}{}
        \begin{alternativas}
            \item The diagram of the situation is the following:
            
            \begin{figure}[H]
                \centering
                \includegraphics[width=0.87\linewidth]{imagens/orbita hiperbolica esquema 1.png}
                \caption{Diagram of the hyperbolic orbit}
                \label{fig.1}
            \end{figure}

            To find the deflection angle \(\phi\), we will use the impulse theorem, considering the horizontal and vertical components of the attractive force:

            \[F_x = -\frac{GMm}{r^2}\cos\theta, \ \ F_y = -\frac{GMm}{r^2}\sin\theta\]

            The impulse theorem tells us that:
            \[\int F_i dt = \Delta p\]

            where \(p_i\) is the momentum of the particle in the direction \(i\).

            On the \(x\)-axis, we have:

            \[-\int \frac{GMm}{r^2}\cos\theta dt = \Delta p_x \]

            Note that integrating with respect to time makes things inconvenient, because finding how \(\theta\) and \(r\) vary in time is complicated. For this reason, we turn to the definition of angular momentum. Using \(L = mr^2d\theta/dt\), we have:

            \[dt = \frac{mr^2}{L}d\theta\]

            Substituting this value into our integral:
            \[-\int \frac{GMm^2}{L}\cos\theta d\theta = \Delta p_x\]

            Notice that our integral goes from \(\theta = 0\), the situation at infinity, to \(\theta = \pi - \phi\), the situation at infinity after the deflection. Thus:

            \[-\frac{GMm^2}{L}\int_0^{\pi-\phi}\cos\theta d\theta = -\frac{GMm^2}{L}\sin\theta\big|_0^{\pi-\phi} = -\frac{GMm^2}{L}\sin(\pi-\phi) = \Delta p_x\]

            Analogously for \(y\):

            \[-\frac{GMm^2}{L}\int_0^{\pi-\phi}\sin\theta d\theta = \Delta p_y\]
            \[=\frac{GMm^2}{L}\cos\theta\big|_0^{\pi-\phi} \equiv -\frac{GMm^2}{L}(1-\cos(\pi-\phi)) = \Delta p_y\]
        
            Also note that \(\Delta p_x = m(v_{f,x}-v_0), \ \Delta p_y = m(v_{f,y}-0)\). We can also find a relation between \(\phi\) and the velocities and, by pure geometry, we obtain:
            \[\tan\phi =-\frac{v_{f,y}}{v_{f,x}}\]
            Since there are only central forces, \(L\) is constant, and from the initial situation it equals \(L = mbv_0\). Now, turning to the calculation:

            \[-\frac{GM}{bv_0}\sin(\pi-\phi) = v_{f,x}-v_0, \ \ -\frac{GM}{bv_0}(1-\cos(\pi-\phi))=v_y\]
            
            With these two equations, we obtain:

            \[\frac{v_{f,y}}{v_{f_x}}= -\frac{GM(1-\cos(\pi-\phi))}{bv_0^2-GM\sin(\pi-\phi)}\]

            That is:

            \[\tan\phi = \frac{GM(1-\cos(\pi-\phi))}{bv_0^2-GM\sin(\pi-\phi)}\]

            Solving for \(\phi\), we obtain the expected result. Another way to carry out the same question is to think purely about the geometry of the situation:
            \begin{figure}[H]
                \centering
                \includegraphics[width=0.7\linewidth]{imagens/morin hiperbole.png}
                \caption{Source: Morin 7.7}
            \end{figure}

            Here, it is evident that:

            \[\phi = \pi - 2\tan^{-1}\left(\frac{b}{a}\right)\]

            where \(b\) is the impact parameter and \(a\) the semiaxis of the hyperbola. From the theory of conics, we have:
            \[\frac{b}{a} = \sqrt{e^2-1} = \sqrt\frac{2EL^2}{m(GMm)^2} = \frac{v_0^2b}{GM} \equiv \gamma b\]

            In both approaches, we arrive at the same result:

            \[\boxed{\phi = \pi - \tan^{-1}(\gamma b)}\]

            as desired.

            \item Consider a ring of thickness \(db\) and radius \(b\). Now consider a very large sphere, centered at \(M\). Any particle that passes through the cross section of the ring of radius \(b\) will hit this sphere making an angle \(\phi\) with the \(x\)-axis, with an angular separation \(d\phi\). Using \(d\cot x/dx = -1/\sin^2x\), we have:
            
            \[\big|\frac{db}{d\phi}\big| = \frac{1}{2\gamma \sin^2(\phi/2)}\]

            The cross-sectional area is given by \(d\sigma = 2\pi bdb\) and the solid angle on a sphere is given by \(d\Omega = 2\pi \sin\phi d\phi\). Therefore:

            \[\frac{d\sigma}{d\Omega} = \frac{2\pi b}{2\pi \sin \phi}\frac{db}{d\phi} = \left(\frac{(1/\gamma)\cot(\phi/2)}{2\sin(\phi/2)\cos(\phi/2)}\right)\left(\frac{1}{2\gamma\sin^2(\phi/2)}\right)\]

            \[\boxed{\frac{d\sigma}{d\Omega} = \frac{1}{4\gamma^2\sin^4(\phi/2)}}\]
        \end{alternativas}            
    \end{pssolution*}
\end{pproblem}

\begin{psidea}{}{}
    In certain questions, the use of vectors is appropriate. One example is the question below. To make things easier, when finding the inclination of the orbit, look at the following figure:

    \begin{figure}[H]
        \centering
        \includegraphics[width=0.7\linewidth]{imagens/vetores ecliptica.png}
        \caption{Vectors}
    \end{figure}

    Here, \(C\) represents a body, \(T\) the Earth, and \(S\) the Sun. The angles \(\lambda\) and \(b\) denote, respectively, the ecliptic longitude and the ecliptic latitude. \(\vec{r}\) is the vector between the Sun and the body, \(\vec{r}_T\) the vector between the Sun and the Earth, and \(\vec{d}\) the vector between the Earth and the body.
    From the figure,

    \[\vec{r} = r(\cos b\cos\lambda \hat{x}+ \cos b \sin\lambda \hat{y} + \sin b \hat{z})\]

    \[\vec{r}_T = r_T(\cos(\omega_\oplus t)\hat{x}+\sin(\omega_\oplus t)\hat{y})\]

    Thus, \(\vec{d}=\vec{r}-\vec{r}_T\):

    \[
    \vec{d} = \big(r \cos b \cos \lambda - r_T \cos (\omega_\oplus t)\big) \hat{x} 
    + \big(r \cos b \sin \lambda - r_T \sin (\omega_\oplus t)\big) \hat{y} 
    + r \sin b \, \hat{z}
    \]

    Calculating the magnitude of the vector \(\vec{d}\), we obtain:

    \[
    d = |\vec{d}| = \sqrt{r_T^2 + r^2 - 2 r_T r \cos b \cos (\lambda - \omega_\oplus t)}
    \]

    Therefore:

    \[
    \boxed{\cos b \cos (\lambda - \omega_\oplus t) = \frac{r_T^2 + r^2 - d^2}{2 r_T r}}
    \]
\end{psidea}

\pts{5}
\begin{pproblem}(Problem set \(1\) - Vinhedo 2024) 
    CR4b-2023 is a probe on a heliocentric orbit that was built and launched in 2023 by the young prodigy Caranguejo. The goal of the probe was to study solar storms and collect data for Caranguejo to analyze at his observatory in Espírito Santo. CR4b-2023, however, during one of its expeditions toward the Sun, was hit by a solar storm and suffered serious damage. Because of this, the probe became disoriented and thus had all of its orbital parameters changed, so that Caranguejo no longer knew its location.

    Aiming to track the position of CR4b-2023 again, Caranguejo used his observatory to collect the values of the angular separations \(\Delta\phi\) between the probe and the Sun and the values of the angular diameter \(\theta_S\) of the probe, all as a function of time \(t\). The table obtained by Caranguejo is shown below.

    \begin{table}[H]
        \centering
        \begin{tabular}{|c|c|c|}
            \hline
            \(\Delta\phi\) (Degrees) & \(\theta_S\) (mas) & \(t\) (Days) \\
            \hline
            0.000 & 99.27 & 0.00 \\
            4.889 & 103.67 & 1.79 \\
            17.354 & 104.72 & 7.43 \\
            29.015 & 93.06 & 20.62 \\
            27.793 & 81.18 & 25.41 \\
            13.460 & 79.77 & 45.98 \\
            11.539 & 77.97 & 51.61 \\
            2.466 & 76.44 & 54.53 \\
            2.656 & 74.84 & 55.31 \\
            7.177 & 73.90 & 57.19 \\
            10.800 & 70.63 & 58.43 \\
            22.464 & 72.92 & 91.67 \\
            13.823 & 80.50 & 130.48 \\
            \hline
        \end{tabular}
        \caption{Values measured by Caranguejo}
    \end{table}

    Assume that, at the initial instant \(t_0 = 0\) when the probe is hit by the solar storm, the Sun was exactly at the point of Libra and the motion of the probe was ascending relative to the ecliptic plane. Also assume that the radius of the probe, supposed spherical, is given by \(R = 30\) km. For this question, an error analysis is not required (although it is important to try to use methods aimed at reducing statistical errors). Based on what was presented:

    \begin{enumerate}[label=(\alph*)]
        \item Calculate the orbital parameters of the orbit of CR4b-2023, that is, its semimajor axis \(a\), eccentricity \(e\), inclination \(i\), longitude of the ascending node \(\Omega\), and argument of perihelion \(\omega\). As in all data-analysis questions, you must provide clearly labeled data tables, clearly labeled graphs, and enough formula derivations to make clear what you measured and how you are deriving your results in order to reduce statistical errors.
        
        \item Taking \((x', y', z')\) as the coordinates of CR4b-2023 in a right-handed Cartesian system in which the \(x'y'\) plane lies in the plane of the probe's orbit, the Sun is at the origin, and the \(x'\) axis points toward the vernal point, write, in a table, the values of at least 5 distinct points \((x', y', z')\). Based on the data found, sketch, on a graph, the orbit of CR4b-2023, indicating the coordinates of its center.
        
        \item Now considering a new right-handed coordinate system \((x, y, z)\), in which the Sun is at the origin, the \(xy\) plane represents the ecliptic plane, and the \(x\) axis points toward the vernal point, write the coordinates of the center of the orbit of CR4b-2023. Finally, calculate the distance \(d_{TC}\) between the center of Earth's orbit and the center of the probe's orbit.
    \end{enumerate}

    \begin{pssolution*}{}{}
        \begin{alternativas}
            \item The idea of this part of the question is to find a relation between the days and the distance from the probe to the Sun. Using basic trigonometry, we can say that the distance from the probe to the Earth is given by:
            \[d_{S-\oplus} = R/\tan^{-1}\left(\frac{\theta_S}{2}\right)\]
            where \(R\) is the radius of the probe.
            Here, we can use the law of cosines to relate this distance to the distance from the probe to the Sun. Considering the Earth's orbit to be circular with radius \(d_{\odot-\oplus}\) and defining \(d\) as the distance between the probe and the Sun, we have:
            \[d^2 = d_{\odot-\oplus}^2+d_{S-\oplus}^2-2d_{\odot-\oplus}d_{S-\oplus}\cos\Delta\phi\]      
            
            With this, we can build the following table, with the values of \(d\) in AU and the values of \(t\) in days.

            \[
                \begin{array}{|c|c|}
                \hline
                \text{t (Days)} & \text{d (UA)} \\
                \hline
                0.00 &    0.17 \\
                1.79 &    0.22 \\
                7.43 &    0.34 \\
                20.62 &    0.50 \\
                25.41 &    0.49 \\
                45.98 &    0.24 \\
                51.61 &    0.22 \\
                54.53 &    0.09 \\
                55.31 &    0.12 \\
                57.19 &    0.18 \\
                58.43 &    0.27 \\
                91.67 &    0.44 \\
                130.48 &    0.25 \\
                \hline
                \end{array}
            \]

            Plotting a graph of \(t\) versus \(d\), we have:
            
            \begin{figure}[H]
                \centering
                \includegraphics[width=0.85\linewidth]{imagens/graficoorbita.png}
            \end{figure}

            Now let us get to work and find the requested parameters: semimajor axis, \(a\), eccentricity, \(e\), inclination, \(i\), longitude of the ascending node \(\Omega\), and argument of perihelion, \(\omega\). The first two can be obtained directly from the data we have analyzed so far.

            Analyzing the data, we can see that the smallest distance between the probe and the Sun is given by \(a_p = 0.09\) AU and the largest distance is given by \(a_a = 0.49\) AU. These distances must correspond to perihelion and aphelion, respectively. Using the definition:
            \[a_p = a(1-e), \ a_a=a(1+e) \rightarrow \frac{a_p}{a_a}=\frac{1-e}{1+e} \approx 0.18\]
            
            Solving, we obtain \(e \approx 0.67\).

            Using \(a = \frac{a_p+a_a}{2}\), we obtain \(a\approx 0.30\) AU.

            To find the angular factors, we have to think a bit more. Starting with \(\Omega\), note that at the initial instant, at \(t=0\), the angular separation between the Sun and the probe is \(\Delta \phi = 0\). The statement tells us that at this moment the Sun is at the point of Libra, so the probe is also at the point of Libra and, since its motion was ascending, we have \(\Omega = 0^\circ\).

            To find \(\omega\) and \(i\), let us look at the following spherical triangle:

            \begin{figure}[H]
                \centering
                \includegraphics[angle=90, width=0.8\textwidth]{imagens/trigesfericoorbita.jpg}
                \caption{Diagram}
            \end{figure}

            From the polar form of the ellipse, \(r=\frac{a(1-e^2)}{1+e\cos\theta}\), where \(\theta\) is the true anomaly of the orbit, we can find \(\omega\). Notice that at the initial instant, \(t=0\), the probe is exactly at the point of Libra, so \(\theta = 2\pi-\omega\) as in the diagram.

            Thus, we have:

            \[d(0) = \frac{a(1-e^2)}{1+e\cos(2\pi-\omega)}\]

            Solving for \(\omega\):

            \[\omega = 2\pi-\cos^{-1}\left(\frac{a(1-e^2)}{d(0)e}-\frac{1}{e}\right) \approx 267.6^\circ \approx 270^\circ\]

            substituting the values found.

            To find the inclination, we will use the law of cosines:

            \[\cos(\theta-\omega)=\cos\lambda\cos b + \sin\lambda\sin b\cos(\pi/2) = \cos\lambda\cos b\]
            
            Note that \(\cos(\theta-\omega) \approx \cos(\theta-270^\circ) = \sin \theta\).

            Rewriting,

            \[\cos b = \frac{\sin\theta}{\cos\lambda}\]

            Using the relation obtained in the previous idea,

            \[\cos b  = \frac{r_T^2 + r^2 - d^2}{2 r_T r\cos (\lambda - \omega_\oplus t)}\]

            Equating these expressions and solving for \(\lambda\),

            \[\lambda = \tan^{-1} \left( 
                \frac{r_T^2 + r^2 - d^2}{2 r_T r \sin \theta \sin (\omega_\oplus t)} - \cot (\omega_\oplus t)
                \right)\]

            where:

            \[
            \theta = \arccos \left( \frac{a (1 - e^2)}{e r} - \frac{1}{e} \right)
            \] 
            
            With this, we can make the following table, with the values of \(b\) and \(\lambda\).

            \[
            \begin{array}{|c|c|}
            \hline
            \lambda \, (\text{Degrees}) & b \, (\text{Degrees}) \\ \hline
            0.000     & 0.00     \\ 
            17.495    & 9.847    \\ 
            45.905    & 22.521   \\ 
            78.492    & 29.499   \\ 
            101.508   & 29.499   \\ 
            134.095   & 22.521   \\ 
            163.504   & 9.847    \\ 
            180.000   & 0.000    \\ 
            197.495   & -9.847   \\ 
            216.005   & -18.747  \\ 
            270       & -30      \\ 
            333.435   & -14.478  \\ 
            17.495    & 9.847    \\ 
            78.492    & 29.499   \\ 
            143.994   & 18.747   \\ \hline
            \end{array}
            \]

            Returning to the spherical triangle and using the four-parts formula, we have:

            \[\cos\lambda\cos90^\circ = \sin\lambda\cot b -\sin90^\circ\cot i\]

            Simplifying:

            \[\tan b = \tan i \sin \lambda\]

            Note that this resembles a function of the type \(y= A + Bx\). Performing a linear regression on a calculator, we find \(B\approx 0.578\). Since \(\tan i = B\), we arrive at:

            \[\boxed{i \approx 30^\circ}\]

        \item Since the Sun is at the origin, the values of \(x'\) and \(y'\) can be calculated by:
            \[
            x' = r \sin \theta
            \]
            and
            \[
            y' = -r \cos \theta
            \]

            Writing five different values of \(x'\) and \(y'\) in a table, we obtain:

            \[
            \text{Table 7: Values of \(x'\) and \(y'\)}
            \]

            \[
            \begin{array}{|c|c|}
            \hline
            x'\text{ AU} & y'\text{ AU} \\ \hline
            0.167 & 0.000 \\ \hline
            0.084 & 0.478 \\ \hline
            -0.219 & 0.261 \\ \hline
            -0.203 & 0.074 \\ \hline
            0.000 & -0.100 \\ \hline
            \end{array}
            \]

            Notice that any values of \(x'\) and \(y'\), provided they are correct, may be used. Furthermore, it is important to increase the spacing between the values of \(x'\) and \(y'\) in order to reduce the errors associated with sketching the orbit. Based on this, we can plot the graph:

            \begin{figure}[H]
                \centering
                \includegraphics[width=0.7\linewidth]{imagens/graforbt.png}  
                \caption{Graph of the orbit} 
            \end{figure}

            Observing the graph, we see that the coordinates \((x', y', z')\) of the center of the ellipse are \(\boxed{(0;\ 0.2; \ 0)}\).

        \item Here, we will use a rotation matrix about the \(x\)-axis, given by:
        
            \[
            \begin{pmatrix}
            x \\
            y \\
            z
            \end{pmatrix}
            =
            \begin{pmatrix}
            1 & 0 & 0 \\
            0 & \cos \alpha & \sin \alpha \\
            0 & -\sin \alpha & \cos \alpha
            \end{pmatrix}
            \begin{pmatrix}
            x' \\
            y' \\
            z'
            \end{pmatrix}
            \] 

            For \(\alpha=i=30^\circ\):

            \[
            \begin{pmatrix}
            x \\
            y \\
            z
            \end{pmatrix}
            =
            \begin{pmatrix}
            1 & 0 & 0 \\
            0 & \cos 30^\circ & \sin 30^\circ \\
            0 & -\sin 30^\circ & \cos 30^\circ
            \end{pmatrix}
            \begin{pmatrix}
            x' \\
            y' \\
            z'
            \end{pmatrix}
            \]

            Then:
            \[
            x = x'
            \]
            \[
            y = y' \cos 30^\circ + z' \sin 30^\circ
            \]
            \[
            z = -y' \sin 30^\circ + z' \cos 30^\circ
            \]

            Substituting the numerical values, we obtain that the coordinates \((x, y, z)\) of the center of CR4b-2023 are given by:
            \[
            (0.000, 0.173, -0.100).
            \]

            Finally, we obtain that the distance between the center of Earth's orbit and the center of the probe's orbit is given by:
            \[
            d = \sqrt{\lvert x^2 + y^2 + z^2 \rvert}
            \]

            that is:
            \[
            \boxed{d = 0.200 \, \text{UA}}
            \]
        \end{alternativas}
    \end{pssolution*}
\end{pproblem}

\pts{5}
\begin{pproblem} (Iran Problem Set) 
    A particle of mass \(m\) is orbiting a massive object of mass \(M\). Show that the impulse required to rotate the orbit by an angle \(\eta\) about one of the foci is given by:
    \[
    \Delta v = \frac{2\mu e}{h}\sin\left(\frac{\eta}{2}\right),
    \]
    where:
    \begin{itemize}
        \item \(h\) is the angular momentum per unit mass;
        \item \(\mu = GM\), where \(G\) is the gravitational constant.
    \end{itemize}
    \begin{pssolution*}{}{}
        Drawing the orbit and the rotated orbit:

        \begin{figure}[H]
            \centering
            \includegraphics[width=0.7\linewidth]{imagens/esquemarotorb.png}
            \caption{Diagram of the rotated orbit.}
        \end{figure}

        Notice that the impulse can be applied only at the intersection points of the two orbits and must be radial, since the angular momentum must remain constant in order to keep the orbit identical. Let us define \(X\) as a point on the orbit before the impulse and \(X'\) as a point on the orbit after the rotation.

        Let us set the following definitions: \(F\) is the focus, \(P\) is the perihelion of the orbit, and \(A\) is the intersection point at which the impulse is applied.

        By symmetry, note that the angles \(AFP\) and \(AFF'\) are equal and therefore have the value \(\eta/2\).

        Since the two orbits have the same angular momentum, the tangential velocity at the point \(A\) must be the same for both orbits. Furthermore, we can see that the point \(A\) has true anomaly \(\theta = \eta/2\), which guarantees that the speed is also the same for both orbits.

        \begin{figure}[H]
            \centering
            \includegraphics[width=0.7\linewidth]{imagens/rotorbt1.png}
            \caption{Velocity vectors.}
        \end{figure}

        The only way for this to be possible is if the radial impulse is given by \(\Delta v = 2v_r\), where \(v_r\) is the radial velocity immediately before the impulse.

        Of course \(v_r^2 = v^2 - v_t^2\), where \(v_t\) is the tangential velocity. Starting from basic definitions:

        \[
        h^2 = \mu a (1 - e^2), \quad r = \frac{a(1 - e^2)}{1 - e\cos\theta}, \quad v_t = \frac{h}{r}, \quad v^2 = \mu \left(\frac{2}{r} - \frac{1}{a}\right).
        \]

        Working with these equations, we arrive at:

        \[
        v^2 = \frac{\mu^2}{h^2}(1 + 2e\cos\theta + e^2), \quad v_t^2 = \frac{\mu^2}{h^2}(1 + e\cos\theta)^2.
        \]

        Finally, substituting into \(v_r^2 = v^2 - v_t^2\), we obtain:

        \[
        v_r^2 = \frac{\mu^2}{h^2}(e^2 - e^2\cos^2\theta) \equiv \frac{\mu^2e^2}{h^2}\sin^2\theta.
        \]

        Substituting \(\Delta v = 2v_r\) and \(\theta = \eta/2\), we arrive at the desired result:

        \[
        \boxed{\Delta v = \frac{2\mu e}{h}\sin\left(\frac{\eta}{2}\right) \blacksquare}
        \]
    \end{pssolution*}
\end{pproblem}


\pts{3}
\begin{pproblem}
    In this question, we will study a simplified model of the effect that confirmed Einstein's theory of general relativity: gravitational lensing. The presence of a massive body curves spacetime, so that stars can serve as lenses in space. Some telescopes, such as the JWST, use this effect to photograph very distant galaxy clusters. One of the photographs taken by the JWST of a gravitational lens is shown below:
    \begin{figure}[H]
        \centering
        \includegraphics[width=0.7\linewidth]{imagens/fotoJWSTeinstenring.png}
        \caption{Source: National Geographic}
    \end{figure} 

    A simplified diagram of gravitational lensing can be seen below. This ``circle'' is known as the Einstein ring.
    \begin{figure}[H]
        \centering
        \includegraphics[width=0.9\linewidth]{imagens/anelE.png}
        \caption{Diagram of the gravitational lens}
    \end{figure}
    In the diagram, the point \(O\) represents the observer, \(L\) the massive body that acts as a lens, \(S\) the true position of the observed object, and \(S'\) its apparent position. \(\alpha\) is the impact parameter, \(\beta\) the angular separation between the observed body and the lens, \(\theta_E\) the angular radius of the Einstein ring, and \(\varphi\) the angular deflection caused by a massive body. The distance \(\overline{OL}\) equals \(D_L\) and the distance \(\overline{OP}\) equals \(D_S\), so that \(D_S-D_L = D_{LS}\).

    All the angles are very small, so that \(\sin x \approx \tan x \approx x\). Furthermore, because they lie at infinity, the lines \(\overline{OP}\), \(\overline{OS}\), and \(\overline{OS'}\) may be approximated as parallel.

    It is a known result of the theory of general relativity that the deflection of light caused by a massive body is given by:
    \[\varphi = \frac{4GM}{\alpha c^2}\]
    \begin{alternativas}
        \item In the limit \(\lim_{\alpha \rightarrow 0}\), find an expression for \(\theta_E\) in terms of the other distances and angles provided.
    
        \item The JWST telescope obtains images in the infrared, of wavelength \(\lambda_{IV}\). How large must its diameter be so that it can resolve an Einstein ring?
    \end{alternativas}

    \begin{pssolution*}{}{}
        \begin{alternativas}
            \item Note that, when \(\alpha \rightarrow 0\), \(\overline{PL}=\overline{AS}=D_{LS}\). Using \(\tan\theta = \sin\theta=\theta\), we have:
            
            \[(1)\ \overline{PS} = D_S\beta = \alpha, \ \ (2)\ \overline{SS'} = D_{SL}\varphi, \ \ (3)\ \overline{PS'} = D_S\theta_E,\] 
            \[(4)\ \overline{OL} = D_L = \overline{LA}/\theta_E = \alpha\theta_E, \ \ (5)\ \overline{LA}=\overline{PS}, \ \ (6)\ \overline{PS'} = \overline{PS}+\overline{SS'}\]

            Substituting \((1), \ (2), \ (3)\) into \((6)\), 

            \[D_S\theta_E=D_S\beta+(D_S-D_L)\varphi\]

            \[\theta_E-\beta = \frac{D_S-D_L}{D_S}\varphi\]

            Substituting the value of \(\varphi\):

            \[\theta_E-\beta = \frac{D_S-D_L}{D_S}\frac{4GM}{\alpha c^2}\]

            Using \((4)\), we have \(\alpha=D_L\theta_E\), and substituting into our equation:

            \[\theta_E-\beta = \frac{D_S-D_L}{D_S}\frac{4GM}{D_L\theta_E c^2}\]

            \[\theta_E^2-\beta\theta_E = \frac{D_S-D_L}{D_S}\frac{4GM}{c^2}\]

            The limit \(\alpha\rightarrow 0\) implies that \(\beta\rightarrow 0\), so:

            \[\boxed{\theta_E=\sqrt{\frac{D_S-D_L}{D_S}\frac{4GM}{c^2}}}\]

            \item Using the Rayleigh criterion \(\theta = R\frac{\lambda}{D}\), where \(R\approx 1.22\). Isolating \(D\) and substituting \(\theta = \theta_E\):
            
            \[\boxed{D =\sqrt{\frac{c^2R^2\lambda_{IV}^2}{4GM}\frac{D_S}{D_S-D_L}}}\]
        \end{alternativas}
    \end{pssolution*}
\end{pproblem}

\pts{4}
\begin{pproblem}
    In this exercise, we will study properties of the eccentricity of orbits and understand how it relates to energy and angular momentum. Denoting \(\mu = GM\) and \(h=L/m\), do what is asked in the following items.
    \begin{center}
        \textbf{Part I: The eccentricity vector}
    \end{center}
    \begin{alternativas}
        \item Write an expression for Newton's second law in vector form for a body of mass \(m\) orbiting a body of mass \(M\) at a distance \(r\).
        
        \item Take the cross product of your expression with \(\mathbf{h}\) and prove that:
        
        \[\frac{d}{dt}(\dot{\mathbf{r}}\times \mathbf{h}) = \frac{\mu}{r}\mathbf{v} - \frac{\mu \dot{r}}{r^2}\mathbf{r}\]

        where \(\dot{\mathbf{r}}\) represents the radial velocity and \(\mathbf{v}\) the velocity.

        \item We can write the right-hand side of the previous equation as:
        \[\frac{d}{dt}\left(A\mathbf{r}\right)\]

        Find the value of \(A\).

        \item Integrate both sides of the equality you found and take the scalar product of both sides with \(\mathbf{r}\). After that, solve for \(r\). Your result should be very similar to what is predicted by geometry alone:
        \[r = \frac{p}{1+e\cos\theta}\]

        Find the values of \(p\) and \(e\). \textbf{Do not forget the integration constant. Your answers may depend on it.}

        \item In the previous item, \(e\) is the eccentricity of the orbit. Returning to the equation initially obtained in item \(c\), after the integration, find a way to express \(\mathbf{e}\) as a vector of magnitude \(e\) that points directly toward the orbited body.
        
        \item Prove that the relation found in the previous item can be written as:
        \[\mu \mathbf{e} = (v^2-\frac{\mu}{r})\mathbf{r} - (\mathbf{r}\cdot\mathbf{v})\mathbf{v}\]

        Where does this vector point?
    \end{alternativas}
    \begin{center}
        \textbf{Part II: Determining orbital elements from the eccentricity vector}     
    \end{center}
    Let us define our coordinate system as follows. Consider a right-handed system, oriented so that the \(xy\) plane coincides with the equatorial plane and \(\mathbf{z}\) points toward the north pole. Define the nodal vector as \(\mathbf{n} = \mathbf{z} \times \mathbf{h}\).
    \begin{alternativas}
        \item Find an expression for \(\mathbf{h}\) in terms of the components of the vector \(\mathbf{r}\) and of the vector \(\mathbf{v}\). After that, find the vector \(\mathbf{n}\) in terms of the components of \(\mathbf{h}\).

        \item Draw a celestial sphere and represent on it: the equatorial plane and an arbitrary orbit (with inclination different from 0). On it, mark the following angles: \(i\), the orbital inclination, \(\Omega\) the longitude of the ascending node, \(\omega\) the argument of periastron, and \(\theta\), the true anomaly. 

        \item Find expressions for all of these angles in terms of \(\mathbf{e}\), \(\mathbf{h}\), \(\mathbf{n}\), \(\mathbf{r}\), and any of the vectors defined in our \(xyz\) coordinate system.
    \end{alternativas}

    \begin{pssolution*}{}{}
        \begin{center}
            \textbf{Part I}
        \end{center}
        \begin{alternativas}
            \item Using the vector form of the gravitational force, we have:
            \[-\frac{GMm}{r^3}\mathbf{r} = m\ddot{\mathbf{r}} \rightarrow \boxed{\ddot{\mathbf{r}}= -\frac{\mu}{r^3}\mathbf{r}}\]

            \item Taking the cross product of this expression with \(\mathbf{h}\), we have:
            
            \[\ddot{\mathbf{r}}\times \mathbf{h} = -\frac{\mu}{r^3}(\mathbf{r}\times\mathbf{h}) = \frac{\mu}{r^3}(\mathbf{h}\times\mathbf{r})\]

            Since \(\mathbf{h}\) is constant, the left-hand side of the equation can be written as \(d(\dot{\mathbf{r}}\times \mathbf{h})/dt\). For the right-hand side, let us write \(\mathbf{h} = \mathbf{r}\times \mathbf{v}\) and use the BAC-CAB rule, which states that \((A \times B) \times C = B(A\cdot C) - C(A \cdot B)\). Thus:

            \[\frac{d}{dt}(\dot{\mathbf{r}}\times \mathbf{h}) = \frac{\mu}{r^3}((\mathbf{r}\times\mathbf{v})\times \mathbf{r}) = \frac{\mu}{r^3}(r^2\mathbf{v} - \mathbf{r}(\mathbf{r}\cdot \mathbf{v}))\]

            Note that \(\mathbf{r}\cdot\mathbf{v}\) is the product of the components of \(\mathbf{r}\) and \(\mathbf{v}\) that lie in the same direction. Thus, \(\mathbf{r}\cdot\mathbf{v} = r\dot{r}\). Substituting into the previous expression:

            \[\boxed{\frac{d}{dt}(\dot{\mathbf{r}}\times \mathbf{h}) = \frac{\mu}{r}\mathbf{v} - \frac{\mu \dot{r}}{r^2}\mathbf{r}}\]

            \item Expanding the expression given in the statement:
            
            \[\frac{d}{dt}(A\mathbf{r}) = \mathbf{r}\frac{d}{dt}A + A \mathbf{v}\]

            Comparing the expressions, it is easy to see that \(\boxed{A = \mu/r}\).

            \item Using the information obtained in the previous items, we have:
            \[\frac{d}{dt}(\dot{\mathbf{r}}\times \mathbf{h}) = \mu\frac{d}{dt}\left(\frac{\mathbf{r}}{r}\right)\]

            Multiplying both sides by \(dt\) and integrating, we have:

            \[\dot{\mathbf{r}}\times\mathbf{h} = \frac{\mu}{r}\mathbf{r} + \mathbf{B}\]

            where \(B\) is a constant vector of integration. Now, taking the scalar product with \(\mathbf{r}\):

            \[\mathbf{r}\cdot(\dot{\mathbf{r}}\times \mathbf{h}) = \mu r + \mathbf{r}\cdot\mathbf{B}\]

            Using \(a\cdot(b\times c) = (a\times b)\cdot c\) and defining \(\nu\) as the angle between \(\mathbf{r}\) and \(\mathbf{B}\), we have:

            \[h^2 = \mu r + rB\cos\nu\]

            Solving for \(r\), we arrive at:

            \[r = \frac{h^2/\mu}{1+(B/\mu)\cos\nu}\]

            Thus:

            \[\boxed{p = \frac{h^2}{\mu}, \ \ \ \ e = \frac{B}{\mu}}\]

            \item Using \(\mathbf{e} = \mu \mathbf{B}\) and returning to item \(d)\), we have:
            \[\dot{\mathbf{r}}\times \mathbf{h} = \frac{\mu}{r}\mathbf{r}+\mu \mathbf{e}\]

            Solving for \(\mathbf{e}\) and noting that \(\dot{\mathbf{r}}\times\mathbf{h} = \mathbf{v}\times\mathbf{h}\), we have:

            \[\boxed{\mathbf{e} = \frac{\mathbf{v}\times\mathbf{h}}{\mu} - \frac{\mathbf{r}}{r}}\]

            \item Writing \(\mathbf{h} = \mathbf{r}\times\mathbf{v}\), we have:
            \[\mathbf{e} = \frac{\mathbf{v}\times(\mathbf{r}\times\mathbf{v})}{\mu} - \frac{\mathbf{r}}{r} \]

            Multiplying both sides of the expression by \(\mu\), and using BAC-CAB, we have:

            \[\mu\mathbf{e} = v^2\mathbf{r} - \mathbf{v}(\mathbf{v}\cdot\mathbf{r}) - \frac{\mu}{r}\mathbf{r}\]

            Rearranging the terms:

            \[\boxed{\mu \mathbf{e} = (v^2-\frac{\mu}{r})\mathbf{r} - (\mathbf{r}\cdot\mathbf{v})\mathbf{v}}\]

            \end{alternativas}
    \end{pssolution*}
\end{pproblem}

\pts{5}
\begin{pproblem}\(\star\star\star\) (Adapted from IPhO 2018) One of the most interesting effects of general relativity is the emission of gravitational waves (GWs) by close binaries. One consequence of this is the loss of energy of the system. In this question, we will study this effect in depth.

\begin{alternativas}
    \item Consider a system formed by stars 1 and 2, with masses \(M_i\) and at distances \(r_i\) from the center of mass. Using Newton's second law, one can show that:
    \[\mathbf{a_1} = -\alpha \frac{\mathbf{r_1}}{r_1^n}\]
    where \(\mathbf{a_1}\) is the acceleration vector of star 1. Find the value of \(\alpha = \alpha (G, M_1, M_2)\) (read as \(\alpha\) in terms of \(G, M_1, M_2\)) and of \(n\).

    \item The total energy of the binary system can be expressed as:
    \[E = A(\omega, \mu, L) - \frac{GM\mu}{L}\]

    where \(\mu\) is the reduced mass of the system, \(M = M_1+M_2\), and \(L = r_1+r_2\). Find the value of \(A\).

    \item The previous expression can be simplified to:
    \[E = \frac{\beta GM\mu}{L}\]

    where \(\beta\) is a dimensionless constant. Find its value.

    The correct theory of gravity, general relativity, was formulated by Einstein in 1915 and predicts that gravity propagates at the speed of light. The messengers that carry information about the interaction are called gravitational waves (GWs). GWs are emitted whenever masses are accelerated, causing the system of masses to lose energy.

    Consider a system of two point particles, isolated from the rest of the Universe. Einstein proved that, for sufficiently small speeds, the emitted waves: 1) have a frequency twice the orbital frequency; 2) can be characterized by a luminosity, that is, an emitted power \(\mathcal{P}\), which is dominated by the formula:

    \[
    \mathcal{P} = \frac{G}{5c^5} \sum_{i=1}^3 \sum_{j=1}^3 \left( \frac{d^3 Q_{ij}}{dt^3} \right) \left( \frac{d^3 Q_{ij}}{dt^3} \right).
    \]

    Here, \( c \) is the speed of light \( c \simeq 3 \times 10^8 \, \mathrm{m/s} \). For a system of two point particles orbiting in the \( x-y \) plane, \( Q_{ij} \) is the following table (\( i, j \) indicate the row/column number):

    The components of \( Q_{ij} \) are given by:

    \[
    Q_{11} = \sum_{A=1}^2 \frac{M_A}{3} \left( 2x_A^2 - y_A^2 \right),
    \]

    \[
    Q_{22} = \sum_{A=1}^2 \frac{M_A}{3} \left( 2y_A^2 - x_A^2 \right),
    \]

    \[
    Q_{33} = -\sum_{A=1}^2 \frac{M_A}{3} \left( x_A^2 + y_A^2 \right),
    \]

    \[
    Q_{12} = Q_{21} = \sum_{A=1}^2 M_A x_A y_A.
    \]

    And \(Q_{i,j} = 0\) for all other possibilities. Here, \(x_A, y_A\) are the positions of star \(A\), with \(A=1,\ 2\), in a Cartesian plane with the center of mass at the origin.

    \item Write the coordinates (\(x_1, y_1\)) and (\(x_2, y_2\)) in terms of \(r_1, r_2, \omega\), and \(t\), with \(t\) the time elapsed since some specific point of the orbit.
    
    \item To calculate the dissipated power, the formula contains a quadruple sum. Solving this expression by brute force is quite complicated and time-consuming. Fortunately, there is a nicer and more elegant way to solve it. To do so, begin by writing \(Q_{i,j}\) as a 3x3 matrix, of the following form:
    
    \[Q = A\left(\begin{matrix}
        a_{11} & a_{12} & a_{13} \\
        a_{21} & a_{22} & a_{23} \\
        a_{31} & a_{23} & a_{33} \\
    \end{matrix}\right)\]

    Find the coefficient \(A = A(\mu, L)\) and complete the matrix \(Q\).

    \item From the previous item, you should obtain:
    \[Q_{ii} = A(b_i + j_i\cos(k t)), \ \ \ Q_{ij}^{i\neq j}= A(p_{ij} \sin(kt)) \]

    Find the values of \(b_i\), \(j_i\), \(p_{ij}\), and \(k\).

    \item You are now able to solve for \(\mathcal{P}\) more smoothly. The expression can be simplified to:
    \[\mathcal{P} = \xi \frac{G}{c^5}\mu^2L^4\omega^6\]
    Find the numerical value of \(\xi\).

    \textbf{HINT: } The double sum:

    \[\sum_{i=1}^N\sum_{j=1}^N (A_{ij})(A_{ij})\]

    where \(A\) is an \(N\times N\) matrix, can be simplified to:

    \[\sum_{i=1}^N\sum_{j=1}^N A_{ij}^2\]

    which represents the sum of the squares of all the elements of the matrix \(A\).

    \item If there were no emission of GWs, the system would remain in equilibrium indefinitely, but because of that emission the energy of the system is not conserved, so that there is a change in the angular velocity of the system, \(\omega\). The formula for the time variation of \(\omega\) has the following form:
    \[\left(\frac{d\omega}{dt}\right)^3 = (3\xi )^3 \frac{\omega^{11}}{c^{15}}(GM_C)^5\]
    
    where \(M_c\) is the so-called \textit{chirp} mass and \(M_c = M_c(M, \mu)\). Find a formula for \(M_c\).

    \item Using the information obtained above, relate the orbital angular velocity \(\omega \) to the gravitational-wave frequency \( f_{\text{GW}} \). Given that, for any smooth function \( F(t) \) and \( a \neq 1 \):

    \[
    \frac{dF(t)}{dt} = \chi F(t)^a \implies F(t)^{1-a} = \chi (1-a)(t_0 - t),
    \]
    
    where \( \chi \) is a constant and \( t_0 \) is an integration constant, show that the equation of the previous item implies that the gravitational-wave frequency is:
    
    \[
    f_{\text{GW}}^{-8/3} = 8\pi^{8/3} \xi \left( \frac{GM_c}{c^3} \right)^{(2/3)+p} (t_0 - t)^{2-p},
    \]
    
    and determine the constant \( p \).

    On September 14, 2015, the event GW150914 was recorded by the LIGO detectors, consisting of two L-shaped arms, each 4 km long. These arms changed their relative length according to the figure below. The detector arms respond linearly to a passing gravitational wave, and the response pattern mimics the wave. This wave was created by two black holes on nearly circular orbits; the loss of energy by gravitational radiation caused the orbit to shrink and, eventually, the black holes to collide. The collision point corresponds, approximately, to the peak of the signal after point D in the figure below.

    \begin{figure}[H]
        \centering
        \includegraphics[width=0.7\linewidth]{imagens/grafico OG 1.png}
        \caption{Strain, that is, the relative change in the length of each arm, in the LIGO H1 detector. The horizontal axis represents time, and the points A, B, C, and D correspond to \(t = 0.000\), \(0.009\), \(0.034\), and \(0.040\) seconds, respectively.}
    \end{figure}

    \item From the figure, estimate \( f_{GW}(t) \) at:

    \[
    t_{AB} = \frac{t_B + t_A}{2} \quad \text{and} \quad t_{CD} = \frac{t_D + t_C}{2}.
    \]
    
    Assuming that the equation for \(f_{GW}\) is valid up to the collision (which, strictly speaking, is not true) and that the two objects have equal masses, estimate the chirp mass, \( M_c \), and the total mass of the system, in terms of solar masses \( M_\odot \simeq 2 \times 10^{30} \, \text{kg}\).

    \item Estimate the minimum orbital separation between the two objects at \( t_{CD} \). 
    Then estimate a maximum size for each object, \( R_{\text{max}} \). 
    Obtain \( \frac{R_\odot}{R_{\text{max}}} \) in order to compare this size with the radius of our Sun, \( R_\odot \simeq 7 \times 10^5 \, \text{km} \). 
    Also estimate their linear orbital speed at the same instant, \( v_{\text{col}} \), comparing it with the speed of light, \( \frac{v_{\text{col}}}{c} \).
\end{alternativas}

\begin{pssolution*}{}{}
    \begin{alternativas}
        \item From Newton's second law,
        \[-\frac{GM_1M_2}{(r_1+r_2)^3}\mathbf{d} = M_1\mathbf{a_1}\]

        where \(\mathbf{d}\) is a vector that leaves star two and goes to star one. Therefore, \(|\mathbf{d}| = r_1+r_2\). Using the center-of-mass theorem, we have

        \[r_2 = \frac{M_1r_1}{M_2}\]

        Substituting and simplifying,

        \[\mathbf{a_1} = -\frac{GM_2}{r_1^3(1+\frac{M_1}{M_2})^3}\mathbf{d}\]

        Now working with the vector \(\mathbf{d}\), note that its direction is the same as that of \(\mathbf{r_1}\), since the line joining 1 and 2 necessarily passes through the CM of the system. Thus,

        \[\mathbf{d} = |\mathbf{d}|\hat{\mathbf{r_1}} = (r_1+r_2)\hat{\mathbf{r_1}} = (1+M_1/M_2)\mathbf{r_1}\]

        Substituting,

        \[\mathbf{a_1} = -\frac{GM_2}{(1+M_1/M_2)^2}\frac{\mathbf{r_1}}{r_1^3}\]

        With this, we can conclude that

        \[\boxed{\alpha = \frac{GM_2}{(1+M_1/M_2)^2}, \ \ \ n = 3 \ }\]

        \item The energy of the system is given by
        \[E = \frac{M_1v_1^2+M_2v_2^2}{2}-\frac{GM_1M_2}{L}\]

        Focusing on the first term and using \(v_A = r_A\omega\), we have

        \[M_1v_1^2+M_2v_2^2 = (M_1r_1^2 + M_2r_2^2)\omega^2 \]

        By the center-of-mass theorem

        \[(M_1r_1-M_2r_2)^2=0 \rightarrow M_1r_1^2 +M_2r^2_2 = \frac{M_1M_2}{M_1+M_2}L^2\equiv \mu L^2\]

        For the second term, it suffices to multiply the numerator and the denominator by \((M_1+M_2)\). Thus,

        \[E = \frac{\mu L^2}{2} - \frac{G}{L}\frac{M_1M_2}{M_1+M_2}(M_1+M_2) \equiv \frac{\mu L^2\omega^2}{2} + \frac{GM\mu }{L}\]

        Thus, \(\boxed{A = \frac{\mu L^2\omega^2}{2}}\).

        \item By Kepler's third law,
        
        \[\omega^2 = \frac{GM}{L^3}\]

        Substituting,

        \[E = \frac{GM\mu}{2L} - \frac{GM\mu}{L} = -\frac{GM\mu}{L}\]

        thus concluding that \(\boxed{\beta = -1/2}\).

        \item Let \(\theta\) be the angle that the vector \(r_1\) makes with the \(x\)-axis. Thus,
        \[(x_1, \ y_1) = r_1(\cos\theta, \ \sin\theta)\]

        The vector \(r_2\) must make an angle of \(\pi + \theta\). Using \(\sin(\pi+\theta) = -\sin(\theta)\) and \(\cos(\pi+\theta) = -\cos\theta\),

        \[(x_2, \ y_2) = -r_2(\cos\theta, \sin\theta)\]

        Defining \(t\) as the time elapsed since the stars crossed the \(x\)-axis, one has \(\theta = \omega t\). Therefore,

        \[\boxed{(x_1, \ y_1) = r_1 (\cos(\omega t), \ \sin(\omega t)) \ \ \ (x_2, \ y_2) = -r_2 (\cos(\omega t), \ \sin(\omega t))}\]

        \item The matrix will have the following form,
        
        \[Q = \left(\begin{matrix}
            Q_{11} & Q_{12} & Q_{13} \\
            Q_{21} & Q_{22} & Q_{23} \\
            Q_{31} & Q_{32} & Q_{33} \\
        \end{matrix}\right)\]

        Invoking the formulas from the statement,
        \[
        Q_{11} = \sum_{A=1}^2 \frac{M_A}{3} \left( 2x_A^2 - y_A^2 \right),
        \]

        \[
        Q_{22} = \sum_{A=1}^2 \frac{M_A}{3} \left( 2y_A^2 - x_A^2 \right),
        \]

        \[
        Q_{33} = -\sum_{A=1}^2 \frac{M_A}{3} \left( x_A^2 + y_A^2 \right),
        \]

        \[
        Q_{12} = Q_{21} = \sum_{A=1}^2 M_A x_A y_A.
        \]

         For \(Q_{11}\), we have

        \[Q_{11} = \frac{M_1}{3}(2x_1^2-y_1^2) + \frac{M_2}{3}(2x_2^2-y_2^2)\]

        Substituting \(x\) and \(y\),

        \[Q_{11} = \frac{M_1}{3}r_1^2(2\cos^2(\omega t) - \sin^2(\omega t)) + \frac{M_2r^2_2}{3}(2\cos^2(\omega t) - \sin^2(\omega t))\]

        \[Q_{11} = \frac{1}{3}(2\cos^2(\omega t) - \sin^2(\omega t))(M_1r_1^2 + M_2r_2^2)\]  
        \[Q_{11} = \frac{\mu L^2}{3}(2\cos^2(\omega t) - \sin^2(\omega t))\]

        For \(Q_{22}\),

        \[Q_{22} = \frac{M_1r_1^2}{3}(2\sin^2(\omega t) - \cos^2(\omega t)) + \frac{M_2r_2^2}{3}(2\sin^2(\omega t) - \cos^(\omega t))\]

        \[Q_{22} = \frac{\mu L^2}{3}(2\sin^2(\omega t) - \cos^2(\omega t))\]


        For \(Q_{33}\)

        \[Q_{33} = -\frac{M_1r_1^2}{3}(\cos^2(\omega t) + \sin^2(\omega t))-\frac{M_2r_2^2}{3}(\cos^2(\omega t) + \sin^2(\omega t))\]

        \[Q_{33} = -\frac{\mu L^2}{3}\]

        Finally, for \(Q_{12} = Q_{21}\),

        \[Q_{12} = Q_{21} = M_1r_1^2\cos(\omega t) \sin(\omega t) + M_2r_2^2\cos(\omega t)\sin(\omega t)\]

        \[Q_{12} = Q_{21} = \mu L^2 \sin(\omega t)\cos(\omega t) = \frac{\mu L^2}{2}\sin(2\omega t)\]

        Since the remaining terms are \(0\), we can choose as \(A\) the value \(\mu L^2/2\), so our matrix \(Q\) takes the form

        \[Q = \frac{\mu L^2}{2}\left(\begin{matrix}
            \frac{4}{3}\cos^2(\omega t) - \frac{2}{3}\sin^2(\omega t)  & \sin(2\omega t) & 0 \\
            \sin(2\omega t) & \frac{4}{3}\sin^2(\omega t) - \frac{2}{3}\cos(\omega t) & 0 \\
            0 & 0 & -\frac{2}{3} \\
            \end{matrix}\right)\]

        For the terms \(a_{11}\) and \(a_{22}\), we will use trigonometric identities, starting with \(a_{11}\),
        \[\frac{4}{3}\cos^2(\omega t) - \frac{2}{3}\sin^2(\omega t) = \frac{2}{3}(2\cos^2(\omega t) - \sin^2(\omega t))\]

        Using \(\cos^2(\omega t) = 1-\sin^2(\omega t)\), we have

        \[ = \frac{2}{3}\left(2-3\sin^{2}(\omega t)\right)\]

        Using the double-angle form, we have

        \[\cos(2x) = \cos^2x -\sin^2x = 1-2\sin^2x \rightarrow \sin^2x = \frac{1-\cos(2x)}{2}\]

        Making this substitution,

        \[=\frac{2}{3}\left(2 - \frac{3}{2} + \frac{3}{2}\cos(2\omega t)\right) = \frac{1}{3}+\cos(2\omega t)\]

        Using the same process for \(a_{22}\), we find
        \[a_{22} = \frac{1}{3}-\cos(2\omega t)\]

        so that our matrix can be simplified to


        \[\boxed{Q = \left(\begin{matrix}
            \frac{1}{3}+\cos(2\omega t) & \sin(2\omega t) & 0 \\
            \sin(2\omega t) & \frac{1}{3} - \cos(2\omega t) & 0 \\
            0 & 0 & -\frac{2}{3} \\
        \end{matrix}\right)}\]

        \item Analyzing the previous matrix, we have
        \[\boxed{A = \frac{\mu L^2}{2}}\]
        and
        \[\boxed{b_1 = b_2 = \frac{1}{3}, \ \ b_3 = -\frac{2}{3}, \ \ j_1 = 1, \ \ j_2 = -1, \ \ j_3 = 0, \ \ p_{12}=p_{21} = 2, \ \ k= 2}\]

        \item The formula provided for the power was
        
        \[
        \mathcal{P} = \frac{G}{5c^5} \sum_{i=1}^3 \sum_{j=1}^3 \left( \frac{d^3 Q_{ij}}{dt^3} \right) \left( \frac{d^3 Q_{ij}}{dt^3} \right).
        \]

        To begin, let us differentiate the matrix \(Q\) with respect to time,

        \[\frac{d^3 Q}{dt^3} = \frac{\mu L^2}{2}\frac{d^3}{dt^3}\left(\begin{matrix}
            \frac{1}{3}+\cos(2\omega t) & \sin(2\omega t) & 0 \\
            \sin(2\omega t) & \frac{1}{3} - \cos(2\omega t) & 0 \\
            0 & 0 & -\frac{2}{3} \\
        \end{matrix}\right)\]

        To differentiate a matrix, it suffices to differentiate all of its elements. Doing this, we arrive at

        \[\frac{d^3Q}{dt^3} = \frac{\mu L^2}{2}\left(\begin{matrix}
            8\omega^3 \sin(2\omega t) & -8\omega ^3\cos(2\omega t) & 0 \\
            -8\omega^3\cos(\omega t) & -8\omega^3\sin(\omega t) & 0 \\
            0 & 0 & 0 \\ 
        \end{matrix}\right)\]

        Taking that factor outside the matrix,

        \[\frac{d^3Q}{dt^3} = 4\mu\omega^3L^2\left(\begin{matrix}
            \sin(2\omega t) & - \cos(2\omega t) & 0 \\
            -\cos(2\omega t) & - \sin(\omega t) & 0 \\
            0 & 0 & 0 \\
        \end{matrix}\right)\]

        It was given as a hint that the term

        \[\sum_{i=1}^3 \sum_{j=1}^3 \left( \frac{d^3 Q_{ij}}{dt^3} \right) \left( \frac{d^3 Q_{ij}}{dt^3} \right)\]

        corresponds to the sum of the squares of the entries of the matrix, so

        \[\sum_{i=1}^3 \sum_{j=1}^3 \left( \frac{d^3 Q_{ij}}{dt^3} \right) \left( \frac{d^3 Q_{ij}}{dt^3} \right) = (4\mu \omega^3 L^2)^2(\sin^2(2\omega t) + 2 \cos^2(\omega t))\]

        Thus,

        \[\mathcal{P} = \frac{G}{5c^5} \sum_{i=1}^3 \sum_{j=1}^3 \left( \frac{d^3 Q_{ij}}{dt^3} \right) \left( \frac{d^3 Q_{ij}}{dt^3} \right) = \frac{G}{5c^5}(32\mu^2\omega^6 L^4)\]

        Rearranging,

        \[\mathcal{P} = \frac{32}{5}\frac{G\mu^2\omega^6L^4}{c^5}\]

        Comparing this expression with the one given in the statement, we have

        \[\boxed{\xi = \frac{32}{5}}\]

        \item \(\mathcal{P}\) represents the power being dissipated by the GWs, so
        \[\mathcal{P} = -\frac{dE}{dt}\]

        Equating \(E\) to the orbital energy,

        \[\mathcal{P} = \frac{GM\mu}{2} \frac{d}{dt}\frac{1}{L}\]

        We can write \(\omega\) as a function of \(L\). To do so, it suffices to use Kepler's third law,


        \[\omega^2 = \frac{GM}{L^3} \rightarrow \frac{1}{L }= (GM)^{-1/3}\omega^{2/3}\]

        Substituting,

        \[\mathcal{P} = \frac{(GM)^{2/3}\mu}{2}\frac{d}{dt}\omega^{2/3} = \frac{(GM)^{2/3}\mu}{2}\frac{d}{d\omega}\omega^{2/3}\frac{d\omega}{dt}\]

        \[= \frac{(GM)^{2/3}\mu}{3\omega^{1/3}}\frac{d\omega}{dt}\]

        Equating this to the expression obtained in the previous item,

        \[\frac{\xi G\mu^2\omega^6L^4}{c^5} = \frac{(GM)^{2/3}\mu}{3\omega^{1/3}}\frac{d\omega}{dt}\]

        Solving for \(d\omega/dt\),
        
        \[\frac{d\omega }{dt} = (3\xi)\frac{G\mu \omega^{19/3}L^4}{(GM)^{2/3}c^5}\]

        Substituting \(L = \left(\frac{GM}{\omega^2}\right)^{1/3}\), we have

        \[\frac{d\omega}{dt} = (3\xi)\frac{G\mu \omega^{19/3}}{(GM)^{2/3}c^5}\left(\frac{GM}{\omega^2}\right)^{4/3}\]

        Cubing,

        \[\left(\frac{d\omega}{dt}\right)^{3} = (3\xi)^3\frac{G^3\mu^3\omega^{19}}{G^2M^2c^{15}}\frac{G^4M^4}{\omega^8}\]

        Simplifying,

        \[\boxed{\left(\frac{d\omega}{dt}\right)^3 = (3\xi)^3\frac{\omega^{11}}{c^{{15}}}(G^5\mu^3M^2)}\]

        Analyzing the expression in the statement, we have

        \[\boxed{M_c = (\mu^3 M^2)^{1/5}}\]

        \item Using the fact given in the statement,
        
        \[\frac{d F(t)}{dt} = \chi F(t)^a \rightarrow F(t)^{1-a} = \chi (1-a)(t_0-t)\]

        we can find an expression for \(\omega(t)\).

        \[\frac{d\omega}{dt} = \frac{3\xi (GM_c)^{3/5}}{c^5}\omega^{11/3} \rightarrow \omega^{1-11/3} = \frac{3\xi (GM_c)^{3/5}}{c^5}(1-11/3)(t_0-t)\]

        Simplifying,

        \[\omega ^{-8/3} = -\frac{3\xi (GM_c)^{3/5}}{c^5}\frac{8}{3}(t_0-t)\]

        \[\omega^{-8/3} = \frac{8\xi (GM_c)^{3/5}(t-t_0)}{c^5}\]

        To find an expression for \(f_{GW}\), we must use

        \[f_{GW} = 2f_{orbt} = \frac{\omega}{\pi}\rightarrow \omega = \pi f_{GW}\]

        Making this substitution, we arrive at


        \[f_{GW}^{-8/3} = \frac{8\pi^{8/3}(GM_c)^{3/5}(t-t_0)}{c^5}\]

        Comparing with the form in the statement, we obtain

        \[\boxed{p = 1}\]

        \item From the figure, we take the two $\Delta t$ as half-periods. Thus, the cyclic gravitational-wave frequency is $f_{GW} = 1/(2 \Delta t)$. The four given points then allow us to calculate the frequency at the mean time of the two intervals as:

        \[
        \begin{array}{|c|c|}
        \hline
        t \, (\text{s}) & f_{GW} \, (\text{Hz}) \\
        \hline
        t_{AB} = 0.0045 & f_{GW}(t_{AB}) = (2 \times 0.009)^{-1} \\
        t_{CD} = 0.037 & f_{GW}(t_{CD}) = (2 \times 0.006)^{-1} \\
        \hline
        \end{array}
        \]
        
        Now we have two pairs of values $(f_{GW}, t)$ for two unknowns $(t_0, M_c)$. In terms of $t_{AB}$ and $t_{CD}$, and dividing the two equations, we obtain:
        
        \[
        t_0 = \frac{A t_{CD} - t_{AB}}{A - 1}, \quad A = \left( \frac{f_{GW}(t_{AB})}{f_{GW}(t_{CD})} \right)^{-8/3}.
        \]
        
        Substituting the numerical values, $A \simeq 2.95$ and $t_0 \simeq 0.054 \, \text{s}$. Now we can use the equation for \(f_{GW}\) at either $t_{AB}$ or $t_{CD}$ and determine $M_c$. One obtains the chirp mass:
        
        \[
        M_c \simeq 6 \times 10^{31} \, \text{kg} \simeq 30 \, M_\odot\]
        
        Thus, the total mass $M$ is:
        
        \[
        M = 4^{3/5} M_c \simeq 69 \, M_\odot
        \]
        
        This result is, in fact, remarkably close to the best estimates that use the full theory of general relativity! [Even though the real objects do not have exactly equal masses and the theory we have just used is not valid very close to the collision.]
    
        \item Kepler's law gives $L = (GM / \Omega^2)^{1/3}$. The second pair of points highlighted on the graph corresponds to the cycle before the merger. Thus, using \(f_{GW} = \omega/\pi\),

        \[
        \omega_{t_{CD}} \sim 2.6 \times 10^2 \, \text{rad/s}.
        \]
        
        Next, using the total mass obtained previously:
        
        \[
        L \sim 5 \times 10^2 \, \text{km}
        \]
        
        Thus, these objects have a maximum radius of $R_{\text{max}} \sim 250 \, \text{km}$. Therefore, they have more than 30 times as much mass, and:
        
        \[
        \frac{R_\odot}{R_{\text{max}}} \sim 3 \times 10^3,
        \]
        
        they are 3000 times smaller than the Sun!
        
        The linear speed is:
        
        \[
        v_{\text{col}} = \frac{L}{2} \omega \simeq 7 \times 10^4 \, \text{km/s}. 
        \]
        
        They are moving at more than 20\% of the speed of light!
        
        \[
        \boxed{
        \begin{aligned}
        \quad & L_{\text{col}} \sim 5 \times 10^2 \, \text{km}, \\
        & \frac{R_\odot}{R_{\text{max}}} \sim 3 \times 10^3, \\
        & \frac{v_{\text{col}}}{c} \sim 0.2.
        \end{aligned}
        }
        \]
    \end{alternativas}
    
\end{pssolution*}
\end{pproblem}

\pts{3}
\begin{pproblem}(Barra 2024) During his stay at Hotel Fazenda Ribeirão, Hugo was observing the star Plo I with his telescope. Comparing its spectrum with those of other stars, he discovers that Plo I is a neutron star and estimates its mass to be \( M_1 = 2M_{\odot} \). He also observes that Plo I has sinusoidal variations in its radial velocity over time and concludes that Plo I must be part of a binary system with another celestial body, Plo II. Observing the velocity curve of Plo I, Hugo also determines that Plo I has a circular orbit of period \( P \) and maximum radial velocity \( K \). However, he does not know the mass \( M_2 \) of Plo II, nor the inclination \( i \) of the binary's orbital plane with respect to the plane of the sky, shown in the following figure.

    \begin{figure}[H]
        \centering
        \includegraphics[width=0.8\linewidth]{imagens/figurabarra.png}
        \caption{Representation of the orbit}
    \end{figure}

    \begin{alternativas}
        \item From Kepler's third law and Newton's second law, find an expression for the \textbf{mass function} of the binary system formed by Plo I and Plo II in terms of \( P \), \( K \), and universal constants. You may use that the mass function is given by:
        \[
            f = \frac{M_2^3 \sin^3 i}{(M_1 + M_2)^2}
        \]
    
        \item Using his measurements, Hugo calculates that the mass function of the binary is \( f = 46.2M_{\odot} \). We wish to find the minimum value \( M_{2,\text{min}} \) of the mass of Plo II, knowing that it corresponds to the case \( i = 90^\circ \). Assume \( M_{2,\text{min}} > M_1 \). From this, determine \( M_{2,\text{min}} \). Is your result consistent with the hypothesis? Answer YES or NO, justifying with calculations.
        
        \item Based on your answer to the previous item, conclude: is it more likely that Plo II is an exoplanet, a white dwarf, a neutron star, or a black hole?
    \end{alternativas}
    \begin{pssolution*}{}{}
        \begin{alternativas}
            \item First, note that \(v_{r,\ \mathrm{apparent}} = |\mathbf{\omega}\times\mathbf{r}| = \omega r \sin i = v_r \sin i\), so \(K = v_r\sin i = \frac{2\pi a_1}{T}\sin i\). Defining \(a=a_1+a_2\) and \(M = M_1+M_2\), where \(a_i\) and \(M_i\) represent the distance between star \(i\) and the CM and the mass of star \(i\), we have Kepler's third law as
            \[\frac{a^3}{T^2} = \frac{GM}{4\pi^2}\]
        \end{alternativas}
    \end{pssolution*}
\end{pproblem}

\newpage

\section{Positional Astronomy}

\pts{5} % problem of the week 114
\begin{pproblem} After getting lost on a voyage, you and Professor Klafke are shipwrecked on an island and find the following sundial:
    \begin{figure}[H]
        \centering
        \includegraphics[width=0.5\linewidth]{imagens/q6.jpg}
        \caption{Sundial}
    \end{figure}
    It is now your task to discover properties of the sundial.
    \begin{enumerate}[label=\textbf{\alph*)}]
        \item In which hemisphere did you find yourselves? How did you reach that conclusion?
        \item After several weeks on the island, you notice that the sundial's shadow seems to follow a straight line. Letting \(\phi\) be the latitude at which you found yourselves, determine the Sun's declination on that day and explain why this happens on only \(2\) days of the year.
        \item Prove that, when the Sun has the declination found above, the shadow follows a straight line. This can be done by the following steps:
        \begin{enumerate}[label=\roman*)]
            \item Determine the orientation of the line and explain why it must have that orientation;
            \item Determine the length of the shadow for a given position of the Sun in alt-azimuth coordinates;
            \item Derive a relation between altitude and azimuth, given that the tip of the shadow lies on the line;
            \item Determine a constant quantity and show that it is constant for every position of the Sun on that day.
        \end{enumerate}
    \end{enumerate}


\begin{pssolution*}{}{}
    \begin{alternativas}
        \item From the figure, we can see that the markings on the left correspond to the morning hours. The Sun is therefore rising on the right, so that its shadow falls on the left. From these arguments, the cardinal point \textit{East} lies on the right and \textit{West} on the right. A vertical sundial always has its gnomon pointed toward the non-elevated celestial pole. Since this one points \textit{south}, we conclude that the sundial is in the \(\boxed{\textit{Northern Hemisphere}}\).

        \item This happens only at the equinoxes, when \(\delta_\odot = 0\).

        \item Following the steps in the statement, 
        \begin{enumerate}[label=\roman*)]
            \item Since \(\delta_\odot=0\), the sky moves exactly from east to west, 
            \item If the gnomon has length \(L\) and the Sun is at an altitude \(h\), geometry gives the shadow length \(l\) as \(l=\frac{L}{\tan h}\).
            \item Consider the following figure, 
            \begin{figure}[H]
                \centering
                \includegraphics[width=0.66\linewidth]{imagens/retarelogiosolvert.png}
                \caption{Diagram of the imaginary line}
            \end{figure}
        \end{enumerate}

        The distance from the center of the sundial to the shadow is

        \[x = l\sin (A-90^{\circ})\equiv l\sin a = L\frac{\sin a}{\tan h}\]

        We now want to prove that this quantity is constant.

        \item Consider the following triangle: 
        
        \begin{figure}[H]
            \centering
            \includegraphics[width=0.5\linewidth]{imagens/trigesf1.jpg}
            \caption{Sun on the equator}
        \end{figure}

        Using the four-parts formula, 

        \[\cot(90^\circ -\phi)\sin(90^{\circ})+\cos a \cos(90) = \cot h \sin a\]

        Simplifying, we have 

        \[\frac{\sin a}{\tan h } = \tan {\phi}\]

        which is constant. Our condition is therefore proved.

    \end{alternativas}
\end{pssolution*}
\end{pproblem}

\pts{5}
\begin{pproblem} (Adapted from T1 - 2024)
    The equation of time is the difference between the right ascension of the mean Sun and the right ascension of the true Sun. With that in mind, let us calculate some of its properties.
    \begin{alternativas}
        \item Find an expression for the equation of time, neglecting Earth's eccentricity, i.e.: the Sun moving at constant speed along the ecliptic. Leave your answer in terms of the time since the March equinox \(T_M\), Earth's period \(T\), and the obliquity of the ecliptic \(\epsilon\).
        \item Now, taking Earth's eccentricity into account, the equation of time has the form:
        \[E.T. = -2e \cdot \sin(k_1 \cdot M) + \tan^2 \left( \frac{\epsilon}{2} \right) \cdot \sin\left( k_2 \cdot (M + \lambda_P) \right)\]
        where \(M\) is the mean anomaly, \(\lambda_P\) the ecliptic longitude of periapsis, and \(e\) Earth's eccentricity. Determine the constants \(k_1\) and \(k_2\). Think about the symmetries involving each term of the equation.
        \\
        The aim of the question is now to find the Sun's declination at the node of the analemma. See the following figure:
        \begin{figure}[H]
            \centering
            \includegraphics[width=0.7\textwidth]{imagens/q16.png}
            \caption{Analemma diagram}
        \end{figure}
        The node is the point at which the path crosses itself.
        \item For small values of the Sun's declination, obtain a formula for it in terms of \(M\), \(\lambda_P\), and \(\epsilon\). Also assume that the eccentricity and the orbital inclination are sufficiently small.
        \item Find the two mean anomalies that have the same declination and the same \(E.T.\).
        \item Find the value of the solar declination at the node.
    \end{alternativas}

\begin{pssolution*}{}{ }
    \begin{alternativas}
        \item The equation of time is given by
        
        \[E.T. = \alpha_M - \alpha_V\]

        where \(\alpha_M\) is the right ascension of the mean Sun, i.e.: a ``fictitious Sun'' that moves along the celestial equator at constant speed, and \(\alpha_V\) is the right ascension of the true Sun.

        Consider the following spherical triangle

        INSERT image

        First, both the true Sun and the mean Sun have constant and equal angular speeds \(\omega_\odot = 2\pi/T\), so \(\alpha_M=\omega_\odot T_M\).

        Using the four-parts formula

        \[\cot \pi/2 \sin\epsilon + \cos\epsilon\cos\alpha_V=\cot(\omega_\odot T_M)\sin\alpha_V\]

        Since \(\cot\pi/2 =0\), we obtain

        \[\alpha_V = \tan^{-1}(\cos\epsilon\tan(\omega_\odot T_M))\]

        Note that when \(\epsilon = 0\), \(\alpha_V = \alpha_M\), as expected.

        Finally, returning to the definition of the equation of time

        \[\boxed{E.T. = \tan^{-1}\left(\cos\epsilon\tan\left(\frac{2\pi}{T} T_M\right)\right) - \frac{2\pi}{T}T_M}\]

        \item To solve this problem, consider the symmetries. For an elliptical orbit, the only symmetry is about the major axis. Once the obliquity of the ecliptic is included, there are two possible symmetries, about the line of the solstices and about the line of the equinoxes. It is therefore reasonable to expect that the sinusoid tied to the orbital eccentricity has one period per year, while the one tied to the obliquity of the ecliptic has two periods per year. Thus \(\boxed{k_1=1, \ k_2 = 2}\)
        
        \item Returning to the previous spherical triangle and using the law of sines, we have
        
        \[\frac{\sin\delta_\odot}{\sin\epsilon} = \sin\lambda\]

        where \(\lambda \) is the ecliptic longitude of the Sun. Using \(\sin\theta \approx \theta\) for small angles, we have

        \[\delta_\odot = \epsilon\sin\lambda\]

        Neglecting the eccentricity, \(\lambda = \lambda_P + M\), so

        \[\boxed{\delta_\odot = \epsilon\sin(\lambda_P + M)}\]

        \item The condition for two mean anomalies to have the same equation of time is 
        \[\sin(M_1+\lambda_P) = \sin(M_2+\lambda_P)\]

        Using \(\sin\theta = \sin(\pi-\theta)\), we have

        \[\boxed{M_2 =\pi -M_1-2\lambda_P}\]

        \item Equating the equations of time
        
        \[-2e\sin(M_1) + \tan^2 \left( \frac{\epsilon}{2} \right)\sin\left(2(M_1 + \lambda_P) \right) = -2e\sin(M_2) + \tan^2 \left( \frac{\epsilon}{2} \right)\sin\left(2(M_2 + \lambda_P) \right)\]

        Substituting \(M_2\)  

        \[-2e\sin(M_1) + \tan^2 \left( \frac{\epsilon}{2} \right)\sin\left(2(M_1 + \lambda_P) \right) = -2e\sin(M_1+2\lambda_P) - \tan^2 \left( \frac{\epsilon}{2} \right)\sin\left(2M_1 + 2\lambda_P \right)\]

        \[\tan^2 \left( \frac{\epsilon}{2} \right)\sin(2M_1+2\lambda_P) = e(\sin M_1 - \sin(M_1+\lambda_P))\]

        Using the double-angle formula for sine and the prosthaphaeresis formulas:

        \[\tan^2\left(\frac{\epsilon}{2}\right)\sin(M_1+\lambda_P)\cos(M_1+\lambda_P) = -e\sin\lambda_P\cos(M_1+\lambda_P)\]

        \[\sin(M_1+\lambda_P) = -\frac{e\sin\lambda_P}{\tan^2\left(\frac{\epsilon}{2}\right)}\approx -\frac{4e\sin\lambda_P}{\epsilon^2}\]

        Finally, substituting into the declination formula:

        \[\boxed{\delta_\odot \approx -\frac{4e\sin\lambda_P}{\epsilon}}\]
    \end{alternativas}
\end{pssolution*}
\end{pproblem}

    
\pts{5}
    \begin{pproblem}
        Porílio, after a long day of classes at ITA, on the summer solstice (\(\phi = 23^\circ 11' S, \ \lambda = 45^\circ 53' W\)) wants to watch the sunset. On the day in question, the equation of time is \(E.T. = -3\) minutes. Porílio does not want to miss the Sun under any circumstances and starts running so that the (center of the) Sun stays on the horizon.
    
        \begin{alternativas}
            \item São José dos Campos follows Brasília time (UTC-3h). At what civil time should Porílio start his run?
            \item If Porílio runs toward one of the geographic poles, what should his initial speed be?
            \item If Porílio decides to climb a mountain of inclination \(i = 30^\circ\) at a constant speed \(v = 1\) m/s, what is the longest time for which Porílio can keep the Sun above the horizon?
            \textbf{Hint: } Small-angle approximations may be useful.
        \end{alternativas}

\begin{pssolution*}{}{ }
    \begin{alternativas}
        \item To find the rising or setting time of a star, we need its hour angle at the moment it rises. Consider the following spherical triangle, 
        
        \begin{figure}[H]
            \centering
            \includegraphics[width=0.8\linewidth]{imagens/trigangulohorario.png}
            \caption{Diagram of the Sun on the horizon}
        \end{figure}

        When the Sun is on the horizon, \(h = 0\). Because this is the summer solstice (in the Southern Hemisphere), \(\delta = -23.47^\circ\). Using the law of cosines together with \(\sin(90 - \theta) = \cos(\theta)\) and \(\cos(90 - \theta) = \sin(\theta)\), we have 

        \[\sin h = \sin\phi\sin\delta + \cos\phi\cos\delta\cos H\]

        Solving for \(H\), and using \(\sin h = \sin 0^\circ = 0\), 

        \[\cos H = -\tan\phi\tan\delta\]

        Solving, we obtain \(H = 100.72^\circ = 6^h 42^m 52^s\)

        The time of sunset is given by \(TSL = 12 + H = 18^h 42^m 52^s\).

        Using \(TC = TSL + (\lambda_{\text{zone}} - \lambda_{\text{local}})/15 - E.T.\), we obtain 

        \[\boxed{TC = 18^h 49^m 20^s}\]

        \item Porílio must run toward the South Celestial Pole, since the Sun's declination is negative. To find his initial speed, differentiate the law of cosines with respect to time at \(h = 0\).
        \[\frac{d\cos H}{dt} = -\tan\delta \frac{d\tan\phi}{dt}\]

        Using the chain rule, 

        \[-\sin H \dot{H} = -\frac{\tan\delta}{\cos^2\phi}\dot{\phi}\]

        where \(\dot{a} = \frac{da}{dt}\). Isolating \(\dot{\phi}\), 

        \[\dot{\phi} = \frac{\cos^2\phi\sin H}{\tan\delta}\dot{H}\]

        The hour angle changes at \(360^\circ / 24^h = 15^\circ / 1^h\). Solving for \(\dot{\phi}\), 

        \[\dot{\phi} = \frac{-28.68^\circ}{h} = -1.39\cdot 10^{-4}\text{ rad/s}\]

        The sign of \(\dot{\phi}\) only means that he is moving toward the South Pole. Using 

        \[\dot{\phi} = \frac{v}{R_\oplus}\]

        we arrive at \(v = R_\oplus|\dot{\phi}|\),
        
        \[\boxed{v \approx 885 \text{ m/s}}\]

        \item We now use the horizon angle. Consider the following figure, 
            \begin{figure}[H]
                \centering
                \includegraphics[width=0.90\linewidth]{imagens/angulohorizonte.png}
                \caption{Source: Magna notes}
            \end{figure}

            The angle \(\theta\) is given by

            \[\cos\theta = \frac{R_T}{R_T + h}\]

            Since \(h<<R_T\), set \(x=h/R_T\) and use the approximations for \(x<<1\).

            \[\cos\theta = \left(\frac{R_T+h}{R_t}\right)^{-1} = \left(1+\frac{h}{R_T}\right)^{-1} \approx 1-\frac{h}{R_T}\]

            Since \(\theta\) is small, use the second-order Taylor expansion 

            \[f(x) \approx \sum \frac{f^{(n)}(a) x^n}{n!}\]

            where \(f^{(n)}(x)\) is the \(n\)th derivative of \(f(x)\). Expanding about \(a = 0\), 

            \[\cos\theta \approx 1 - \frac{\theta^2}{2}\]

            Equating the two expressions, 

            \[1 - \frac{\theta^2}{2} = 1 - \frac{h}{R_T} \rightarrow \theta = \sqrt{\frac{2h}{R_T}}\]

            To find how long the Sun stays above the horizon, compute the time over which the Sun's apparent altitude is offset by the angle \(\theta\), so that the Sun remains on the horizon. The law of cosines then becomes 

            \[\sin\theta = \sin\phi\sin\delta + \cos\phi\cos\delta\cos H\]

            Using \(\sin\theta \approx \theta\), 

            \[\theta = \sin\phi\sin\delta + \cos\phi\cos\delta\cos H\]

            Differentiate both sides with respect to time. (Here \(\dot{\phi}\) is neglected, because Porílio's speed is very small.)

            \[\frac{d\theta}{dt} = -\cos\phi\cos\delta\sin H \dot{H}\]

            To differentiate \(\theta\) with respect to time, note that \(h(t) = v t \sin i\).

            \[\frac{d\theta}{dt} = \sqrt{\frac{2v\sin i }{R_T}}\frac{d}{dt}\sqrt{t} = \sqrt{\frac{v\sin i}{2 R_T t}}\]

            Equating the two expressions, 

            \[\frac{v\sin i}{2R_T t} = (\cos\phi\cos\delta\sin H \dot{H})^2\]

            \[t = \frac{v\sin i}{2R_T(\cos\phi\cos\delta\sin H \dot{H})^2}\]

            \[\boxed{t \approx 10.82 \text{ s}}\]

    \end{alternativas}
    
\end{pssolution*}
\end{pproblem}



\pts{3}
\begin{pproblem}
    Mauí, a renowned astronomer, wants to spend his year-end vacation in Sergipe \((10^\circ \ 54'\ 33'' S, \ 37^\circ 4'29'' W)\) and wants to know the rising time of one of his favorite stars, All Kaf al Dij Ma \Romannum{3} \((\delta = +10^\circ \ 56'\ 58'', \ \alpha = 2h \ 58m \ 43s\)), and needs your help to calculate its rising time in the following situations:
    
    \begin{alternativas}
        \item Calculate the rising time of All Kaf al Dij Ma \Romannum{3} at the December solstice.
        \item How long will the star spend above the horizon?
        \item Estimate the length of time during which the star is visible.
    \end{alternativas}

\begin{pssolution*}{}{ }
    \begin{alternativas}
        \item As seen above, the hour angle at which a star rises is given by 

        \[\cos H = - \tan \delta \tan \phi \rightarrow H = \]

        The time follows from 

        \[TSL = \alpha + H\]

        which gives \(\boxed{TSL = }\)

        \item The time the star spends above the horizon is 
        
        \[\Delta t = 2 H = \]

        \item At the summer solstice, \(\delta_\odot = -23.47\), so the Sun's hour angle is \(H_\odot = -\tan\delta_\odot \tan\phi = \). That is, the time for which the Sun is visible is \(\delta t_\odot = \). A valid estimate of the time for which the star is visible is 
        \[\Delta t = 2 (H - H_\odot)\]

        These are the intervals during which the star is above the horizon while the Sun is not: 
       
        \[\boxed{\Delta t = }\]
    \end{alternativas}
\end{pssolution*}
\end{pproblem}


\pts{4}
\begin{pproblem}(Lucas Cavalcante - \href{https://noic.com.br/olimpiadas/astronomia/problemas-da-semana/astronomia-semana-111/}{Week 111})
    Consider a horizontal coordinate system in which azimuth \(0\) corresponds to the north cardinal point. One asteroid was observed at altitudes \(h_1\) and \(h_2\) and azimuths \(A_1\) and \(A_2\), while another asteroid was observed at the coordinates \(h_3\), \(h_4\), \(A_3\), and \(A_4\). Find the coordinates of the points where the asteroids meet.
    
    \begin{psidea}{Einstein notation}{}
    While solving this problem, I will use Einstein notation. It is a tool from physics that makes vector calculations easier. In ordinary notation, a vector is written as
    \[\mathbf{a} = (a_1 \hat{\mathbf{e}}_1, \ a_2 \hat{\mathbf{e}}_2, \ a_3 \hat{\mathbf{e}}_3)\]
    In Einstein notation, this is replaced by:
    \[\mathbf{a} = a_i \hat{\mathbf{e}}_i\]
    Repeated indices indicate a sum (\(a_i \hat{\mathbf{e}}_i = a_1 \hat{\mathbf{e}}_1 + a_2 \hat{\mathbf{e}}_2 + a_3 \hat{\mathbf{e}}_3 + ...\))
    
    The scalar product of two vectors, in Einstein notation, is defined by:
    \[\mathbf{a}\cdot \mathbf{b} = a_ib_j\delta{ij}\]
    where \(\delta_{ij}\) is the \textit{Kronecker delta}, defined by:
    \[\delta_{ij} = \left \{ \begin{matrix} 0, & \mbox{if } i \ne j \\ 
                                            1, & \mbox{if } i = j\end{matrix} \right.\]
    
    Thus, \(\mathbf{a}\cdot \mathbf{b} = a_i b_i \equiv a_1 b_1 + a_2 b_2 + a_3 b_3\)

    For the cross product, we have:
        \[
        \mathbf{a} \times \mathbf{b} = a_i b_j \Epsilon_{ijk} \hat{\mathbf{e}}_k
        \]
        Here, \(\Epsilon_{ijk}\) is called the \textit{Levi-Civita tensor}, or Levi-Civita symbol, and is defined as:
    
        \[
        \Epsilon_{ijk} = 
        \begin{cases} 
        +1 & \text{if } (i, j, k) \text{ is an even permutation of } (1, 2, 3), \\
        -1 & \text{if } (i, j, k) \text{ is an odd permutation of } (1, 2, 3), \\
        0 & \text{if two or more indices are equal}.
        \end{cases}
        \]

        For example, \(\Epsilon_{123} = \Epsilon_{312} = \Epsilon_{231} = 1\) and \(\Epsilon_{132} = \Epsilon_{213} = \Epsilon_{321} = -1\).
        Explicitly 
        \[\mathbf{a} \times \mathbf{b} = (a_2 b_3 - a_3 b_2) \hat{\mathbf{e}}_1 + (a_3 b_1 - a_1 b_3) \hat{\mathbf{e}}_2 + (a_1 b_2 - a_2 b_1) \hat{\mathbf{e}}_3\]
    \end{psidea}

\end{pproblem}


\pts{5}
\begin{pproblem} (Problem set \(1\) - Vinhedo 2024) Consider a sundial consisting of a vertical dial
    and a gnomon, as shown in the figure.
    \begin{figure}[H]
        \centering
        \includegraphics{imagens/q20.png}
        \caption{Sundial diagram.}
    \end{figure}
    \begin{alternativas}
        \item Suppose a vertical sundial whose gnomon is correctly pointed south.
        Let \(\phi\) be the observer's latitude, \(H\) the Sun's hour angle, and \(\delta\) the solar declination. Determine
        the relation between \(\theta\) and \(\gamma\) in terms of these variables.
        
        \item From the previous item, what must \(\theta\) be for the sundial to work throughout the year?

        \item Using the correct value of \(\theta\), determine the relation between \(\gamma\) and \(H\), in terms of \(\phi\) only.
    \end{alternativas}
\end{pproblem}

\pts{2} 
\begin{pproblem}
    What is the smallest height a giant (who lives a VERY long time), located at the South Pole, must have in order to see every star in the sky?

    \begin{pssolution*}
        The angle below the horizon that the giant can see is given by 
    
        \[\cos\theta = \frac{R_\oplus}{R_\oplus + H}\]
    
        where \(H\) is the giant's height. Because of the precession of the equinoxes, the south pole will lie \(23.47^\circ\) above the ecliptic. The horizon angle must therefore be \(\theta = 90^\circ - 23.47^\circ = 66.53^\circ\). Solving for \(H\), 
    
        \[\boxed{H = R_\oplus\left(\frac{1}{\cos\theta}-1\right) \approx 9.62\cdot 10^6\text{ m}}\]
    \end{pssolution*}
\end{pproblem}    


\pts{3}
\begin{pproblem}
    Toleduardo wants to watch the sunset, but he spent too long at SorveterITA. Toleduardo is a \textit{Just in time} man and wants to know how long the sunset will last, so that he can be late at his leisure. Given that Toleduardo only goes to SorveterITA on March 21 or September 22, calculate the duration of the sunset at SorveterITA.
    \\
    \textbf{Data: } Latitude of SorveterITA \(\phi = 23^\circ \ 10' \ 45''S, \ \ \lambda = 45^\circ 53' 14''W\).
\end{pproblem}



\pts{3}
\begin{pproblem} (Problem set 2 - Vinhedo 2021)
    In February 2015, the Moon began a cycle of monthly occultations of Aldebaran ($\alpha$ Tau). That is, every month the Moon passed in front of Aldebaran for an observer on Earth. Note that these occultations did not necessarily occur for observers at the same location every month. 

    Calculate the date (month and year) when this cycle of monthly occultations ends. Assume that the Moon's orbit is circular.

    \subsection*{Data:}

    \begin{itemize}
        \item Ecliptic latitude of Aldebaran = $5.47^\circ$
        \item Nodal precession period of the Moon = 18.6 years
    \end{itemize}
\end{pproblem}





\pts{4}
\begin{pproblem} (Problem set 2 - Vinhedo 2021)
    Miguel lives on an isolated island in the South Pacific Ocean, at longitude
    $\lambda = 176^{\circ}09^{\prime}137.7^{\prime\prime}\text{W}$. Over the year, the place where the Sun rises, as seen by Miguel, varies by $\Delta A = 67^{\circ}03^{\prime}81^{\prime\prime}$ along the horizon. For the first two items, neglect atmospheric refraction. With this information, find:

    \begin{alternativas}
        \item The latitude $\phi$ and the name of the island. (Consult Google Earth or similar software)
        \item The range of times at which the Sun rises on the island, given the peculiar time zone $UT + 12\frac{3}{4}$.
        \item Taking atmospheric refraction into account, would the length of the interval in the previous item decrease, increase, or stay constant?
    \end{alternativas}


\end{pproblem}


\pts{3}
\begin{pproblem}(Problem set 2 - Vinhedo 2022) 
    Bruno decided to rent a house for his vacation in Cuiabá ($15.3^{\circ} \text{S}$, $56.1^{\circ} \text{W}$). As a good astronomer, Bruno spent his nights sitting in a chair, watching the stars through a giant door facing the south cardinal point.

    The door was 4.00 meters high and 1.50 meters wide. Bruno used to sit
    1.00 meter from the door, perfectly aligned with its center horizontally. That is, the line
    segment between Bruno's eyes and the point exactly halfway across the door horizontally and
    at eye level forms an angle of $90^{\circ}$ with the plane of the door. Bruno's eyes are
    1.20 meters above the floor when he is sitting in the chair.
    
    To make his observations easier, Bruno created a coordinate system based on the position of the
    door, from which he saw the stars at the place where he was sitting, using meters as
    the reference unit. The origin of the system is at the lower-left corner. The coordinates
    in $x$ increase to the right and the coordinates in $y$ increase upward. Thus, a
    star seen 1 meter in from the left side of the door and 2 meters above the floor would be represented
    by the coordinates $(1, 2)$.
    
    Bruno was very interested in Shaula ($\lambda \ \text{Sco}, \delta = 37.1^{\circ} \text{S}$). Determine the coordinates of
    Shaula at the instant the star became visible to Bruno when observed through the door. Assume that Shaula was below the horizon when Bruno began watching the sky.
    

\end{pproblem}

\newpage

\section{Photometry and Modern Physics}

\pts{3} 
\begin{pproblem}
    The goal of this question is to derive an expression for the \textit{optical depth}. Imagine that a beam of light passes through a region of space containing a certain number of particles. If \(I_0\) is the intensity of the light before it passes through that region of space and \(I\) is the intensity afterward, the optical depth is defined as:
    \[I = I_0e^{-\tau}\]
    where \(\tau\) is the optical depth.
    \\
    Consider a region of space with particle number density \(n\).
    \begin{alternativas}
        \item What is the number of particles in an area \(A\) and thickness \(dz\)?
        \item Assuming that each particle has cross section \(\sigma\), what is the area blocked by the particles?
        \item Find the formula for \(\tau\). 
    \end{alternativas}

\begin{pssolution*}{}{}
    \begin{alternativas}
        \item The volume number density of particles is \(n\), so the number \(dN\) of particles in a volume \(dV\) is
        \[\boxed{dN = ndV \equiv nAdz}\]

        \item If each particle occupies an area \(\sigma\), the area occupied by \(dN\) particles is
        \[\boxed{dS = \sigma dN = \sigma nAdz}\]

        \item The area occupied by the particles is expected to obstruct the passage of light. Since the intensity is proportional to the area available for the light to pass through, we have
        \[\frac{dI}{I} = -\frac{dS}{A} = -n\sigma dz\]
        Integrating, we obtain

        \[I = I_0 e^{-n\sigma z}\]

        That is,

        \[\boxed{\tau = n\sigma z}\]
    \end{alternativas}
\end{pssolution*}
\end{pproblem}

\pts{2}
\begin{pproblem} The Triangulum Galaxy, \(M33\), is the third-largest galaxy in the Local Group. It lies at a distance \(d=970 \text{ kpc}\) from us and has an apparent magnitude of \(5.72\). Given that it contains approximately \(40\) billion stars, find the average luminosity of the stars in \(M33\). Does your estimate seem consistent with reality? Why?
\begin{pssolution*}{}{}
    Using Pogson's equation

    \[m-M_\odot = -2.5\log \left(\frac{L_{g}}{L_\odot}\left(\frac{10 \text{ pc}}{d}\right)^2\right)\]

    Solving for \(L_g\),

    \[L_g = L_\odot \left(\frac{d}{10\text{ pc}}\right)^2 10^{-0.4(m-M_\odot)}\]

    We also have \(L_g = \overline{L_\star} N\), and thus

    \[\overline{L_\star} = \frac{L_\odot}{N} \left(\frac{d}{10\text{ pc}}\right)^2 10^{-0.4(m-M_\odot)}\]

    Using the values from the table of constants, we obtain

    \[\overline{L_\star} \approx 0.1L_\odot\]

    This estimate does not match reality, since factors such as interstellar extinction increase the apparent magnitude of \(M33\). A consistent estimate would be of the same order of magnitude as the luminosity of the Sun.

\end{pssolution*}
\end{pproblem}

\pts{4}
\begin{pproblem}
    Consider a universe with a constant stellar number density, that is, a number of stars per unit volume, of value \(n\). Assume that every star has \(L = L_\odot\) and that the universe contains a total of \(N_0\) stars.
    
    \begin{alternativas}
    \item What is the probability that the absolute magnitude of a star, seen from the center of the universe, lies between \(m\) and \(m+dm\), where \(dm\) is an infinitesimal interval of magnitude? Leave your answer in terms of \(m\) and the largest possible magnitude, \(m_{lim}\), of a star at the "edge" of the universe.

    \item What is the probability that a star can be observed with the naked eye?
    
    \textbf{Data:} \[\int_{-\infty}^{6}10^{0.6(x-A)}dx \approx 2881.6 e^{-1.381 A}\]
    \end{alternativas}
\begin{pssolution*}{}{}
    \begin{alternativas}
    \item The apparent magnitude of a star is given by

    \[m = -2.5\log F + C\]

    where \(C\) is a constant. Writing \(F = L/4\pi r^2\), we have

    \[m = -2.5\log\left(\frac{L}{4\pi r^2}\right) + C = 5\log r -2.5 \log \left(\frac{L}{4\pi}\right)+ C\]

    Since \(L\) is the same for every star, we can absorb the second term into the constant. Thus

    \[m=5\log r + C \rightarrow r = 10^{0.2(m-C)}\]

    Differentiating the expression for \(m\), we have

    \[\frac{dm}{dr} = \frac{5}{r\ln 10}\]

    Now we find the probability that a star lies at a distance \(r\).

    The number of stars between \(r\) and \(r+dr\) is

    \[dN(r) = 4\pi n r^2dr\]

    Dividing by \(N_0\), the probability that a star lies at that distance is

    \[dP(r) = \frac{dN(r)}{N_0} = \frac{4\pi n r^2}{N_0}dr \rightarrow dr = \frac{N_0}{4\pi n r^2} dP(r)\]

    Substituting the expression for \(dr\),

    \[\frac{r\ln 10}{5}dm = \frac{N_0}{4\pi n r^2}dP(r)\]

    Isolating \(dP(r)\),

    \[dP(r) = \frac{4\pi \ln 10 n r^3}{5 N_0}dm\]

    Substituting the expression for \(r\), we change variables \(r\rightarrow m\) in \(dP\):

    \[dP(m) = \frac{4\pi n \ln 10  }{5 N_0}10^{0.6(m-C)}dm\]

    To find \(C\), compare with the limiting magnitude at the "edge" of the universe. Since the universe contains \(N_0\) stars, its radius is given by

    \[N_0 = \frac{4\pi R^3}{3} \rightarrow R = \left(\frac{3N_0}{4\pi}\right)^{1/3}\]

    Writing the equation for \(m_{lim}\),
    
    \[m_{lim} = 5\log R + C\]

    we can isolate the constant \(C = m_{lim}-5\log R\). Substituting into the expression for \(dP(m)\),

    \[dP(m) = \frac{4\pi n \ln 10  }{5 N_0}10^{0.6(m-m_{lim}+5\log R)}dm\]

    Working through this expression,

    \[dP(m) = \frac{4\pi n \ln 10  }{5 N_0}10^{0.6(m-m_{lim})}10^{3\log R}dm\]

    Using the logarithm identities \(a\log b = \log b^a\) and \(10^{\log x} = x\), we have

    \[dP(m) = \frac{4\pi n R^3\ln 10}{5N_0}10^{0.6(m-m_{lim})}dm\]

    Substituting \(R\),

    \[\boxed{dP(m) = \frac{3\ln 10}{5}10^{0.6(m-m_{lim})}dm}\]

    The probability therefore depends on neither \(n\) nor \(N_0\).

    \item For a star to be visible, its magnitude must satisfy \(m\leq 6\). Thus
    
    \[P(m\leq 6) = \int_{-\infty}^6 dP(m) = \frac{3\ln 10}{5}\int_{-\infty}^610^{0.6(m-m_{lim})}dm\]

    Using the integral given in the statement,

    \[P(m\leq 6) = \frac{3\ln 10}{5}2881.6 e^{-1.381 m_{lim}}\]

    Simplifying the numerical factors,

    \[\boxed{P(m\leq 6) \approx 3981.08 e^{-1.381 m_{lim}}} \]

    Out of curiosity, taking \(m_{lim}\approx 40\) as an estimate, we would have

    \[P(m\leq 6)\approx 10^{-20}\]

    Even in the simplest model of the universe, our capability and our insignificance prevail.
    \end{alternativas}

\end{pssolution*}
\end{pproblem}

\pts{4}
\begin{pproblem}
    The \textit{surface brightness}, flux per solid angle per frequency, \(B_\nu\), is given by Planck's law:
    \[B_\nu = \frac{dF}{d\nu d\Omega} = \frac{2h\nu^3}{c^2(e^{\frac{h\nu}{k_BT}}-1)}\] 
    \begin{alternativas}
        \item Find an expression for:
        \[B_\lambda = \frac{dF}{d\lambda d\Omega}\]
        \item Starting from \(B_\lambda\), find the peak wavelength \(\lambda_{max}\) emitted by a star of temperature \(T\) (you will have to solve something numerically).
        \item At low frequencies we have the Rayleigh-Jeans approximation. Obtain an expression for \(B_\nu\) at low frequencies.
        \end{alternativas}

\begin{pssolution*}{}{}
    \begin{alternativas}
        \item We have
        \[B_\lambda = \frac{dF}{d\lambda d\Omega} = \frac{dF}{d\nu d\Omega}\frac{d\nu}{d\lambda}\]

        Using the fundamental wave relation \(c=\nu\lambda \rightarrow \nu = c/\lambda\), we arrive at

        \[B_\lambda = \frac{2h(c/\lambda)^3}{c^2(e^{\frac{hc}{\lambda k_BT}}-1)}\frac{d}{d\lambda}\left(\frac{c}{\lambda}\right) = \frac{2hc^2}{\lambda^5}\frac{1}{(e^{\frac{hc}{\lambda k_BT}}-1)}\]

        \item To solve this part we need the maximum of \(B_\lambda\). Differentiating \(B_\lambda\) directly is extremely tedious. A more efficient approach is to take the natural logarithm of \(B_\lambda\) and differentiate that. At a maximum, both procedures give the same result.
        
        \[\ln B_\lambda = -5\ln \lambda - \ln (e^{hc/\lambda k_BT}-1) + C\]

        where \(C\) is a constant that absorbs the \(\ln\) of the other constants in \(B_\lambda\). Differentiating further,

        \[\frac{d \ln B_\lambda}{d\lambda} = -\frac{5}{\lambda} + \frac{e^{hc/\lambda k_BT}}{e^{hc/\lambda k_BT}-1}\frac{hc}{k_BT\lambda^2}=0\]

        Defining \(x = \frac{hc}{k_BT\lambda}\),

        \[\frac{xe^x}{e^x-1}=5\]

        \[x=5e^{-x}(e^x-1) = 5(1-e^{-x})\]

        Solving by iteration for \(x\) gives \(x\approx 4.965\). Returning to the definition of \(x\),

        \[4.965 = \frac{hc}{k_B T\lambda_{max}}\rightarrow \boxed{\lambda_{max} \approx \frac{2.989 \cdot 10^{-3}}{T}}\]

        This is Wien's law, commonly written \(\lambda = b/T\), where \(b\approx 2.989 \cdot 10^{-3}\).

    \end{alternativas}
\end{pssolution*}
\end{pproblem}

\pts{3}
\begin{pproblem}
    The luminosity of a secondary body depends on the visible illuminated area of the object. In this question we briefly study that phenomenon.
    \begin{alternativas}
        \item  Consider the following situation:
        \begin{figure}[H]
            \centering
            \includegraphics[width=0.8\linewidth]{imagens/fotometria 1.png}
            \caption{Sun-Earth-Moon diagram}
        \end{figure}

        Find an expression for the ratio \(\Phi\) of the illuminated area, as a function of \(\alpha\), to the total area of the planet.
        
        \item Find the angles \(\alpha\) at which the Moon is in the phases new, waxing, full, and waning, respectively.
    \end{alternativas}

\begin{pssolution*}{}{}
    To solve the problem, follow the diagram below:
    \begin{figure}[H]
       \centering
       \includegraphics[width=0.8\linewidth]{imagens/fotometriaplanetaria.png}
       \caption{Source: Introduction to Planetary Photometry} 
    \end{figure}

    On the left is a sketch of how the planet appears in the sky, and on the right a view "from above". The area we see is half the area of the circle plus half the area of an ellipse with semi-major axis \(R\) and semi-minor axis \(R\cos\alpha\). Thus

    \[\Phi = \frac{\pi R^2/2 + \pi R^2\cos\alpha/2}{\pi R^2} \ \boxed{\therefore \Phi = \frac{1+\cos\alpha}{2}}\]

    At new moon, \(\Phi = 0\); at waxing, \(\Phi = 1/2\); at full moon, \(\Phi = 1\); and at waning, \(\Phi = 1/2\). Thus

    \[\boxed{\alpha_{nova} = 180^\circ, \ \alpha_{crescente} = 90^\circ, \ \alpha_{cheia} = 0^\circ, \ \alpha_{minguante} = 90^\circ}\]
\end{pssolution*}
\end{pproblem}

\pts{5}
\begin{pproblem} (Magna notes)
    In this problem we model the effect of the atmosphere on Earth. Suppose that the Sun
is a blackbody of temperature \(T_1\) and radius \(R_1\). Earth is a sphere located
at a distance \(R\) from the Sun, with radius \(R_3\). Earth's emissivity is \(\epsilon_3\).
\begin{alternativas}
    \item If Earth had no atmosphere, determine its equilibrium temperature, \(T_3\).
    \item  Now we include the atmosphere. Model it as a spherical shell of gas,
    with emissivity \(\epsilon_2\) and outer radius \(R_2>R_3\), concentric with Earth. In thermal equilibrium,
    its absorptivity at ultraviolet and infrared wavelengths is \(\epsilon_2\). The
    atmosphere transmits a fraction \(t\) of the ultraviolet radiation but is completely opaque to
    infrared radiation. Assuming that the Sun emits ultraviolet light while Earth emits and re-emits
    in the infrared, determine the temperature \(T_2\) of the atmosphere and the temperature \(T_3\) of Earth in thermodynamic
    equilibrium. Assume that the atmosphere is a perfect thermal conductor, so that
    all radiation incident on it is distributed uniformly over its surface.
\end{alternativas}

\begin{pssolution*}{}{}
    \begin{alternativas}
        \item In thermodynamic equilibrium, the power absorbed equals the power emitted. The emissivity and the absorptivity are also equal. Thus
        \[\frac{4\pi R_1^2\sigma T_1^4}{4\pi R^2}\pi R_3^2\epsilon_3 = 4\pi R_3^2\sigma T_3^4\epsilon_3\]

        Solving for \(T_3\),

        \[\boxed{T_3 = T_1\left(\frac{R_1^2}{4\pi R^2}\right)^{1/4}}\]

         \item The temperature of the atmosphere is
         
         \[T_2 = T_1\left(\frac{R_1^2}{4\pi R^2 }\right)^{1/4}\]

         The reasoning is analogous to the previous part. The luminosity transmitted by Earth is

         \[L_2 = {4\pi R_2^2\sigma T_2^4\epsilon_2} t\]

         Earth reflects all of this radiation, and the atmosphere reflects only a fraction t back, so the total luminosity is

         \[L_{2,T} = 4\pi R_2^2\sigma T_2^4\epsilon_2 (t+t^2+t^3+...)\]

         The expression in parentheses is an infinite geometric series, with first term \(q_1 = t\) and common ratio \(r = t\). Since \(t<1\), the sum is

         \[(t+t^2+t^3+...) = \sum_{n=1}^{\infty}t^n = \frac{t}{1-t}\]

         Equating the two luminosities,

         \[L_{2,T} = L_3 \rightarrow \frac{4\pi R_2^2\sigma T_2^4 \epsilon_2 t}{1-t} = 4\pi R_3^2\sigma T_3^4\]

         Solving for \(T_3\),

         \[T_3 = T_2\left(\frac{R_2^2\epsilon_2}{R_3^2}\frac{t}{1-t}\right)^{1/4}\]

         Substituting \(T_2\),

         \[\boxed{T_3 =T_1\left(\frac{R_1^2R_2^2}{4\pi R^2R_3^2}\frac{\epsilon_2t}{1-t}\right)^{1/4}}\]


    \end{alternativas}
    
\end{pssolution*}
\end{pproblem}


\pts{2}
\begin{pproblem}(Problem set 2 - 2021)
    The Ring Nebula (M57) has an apparent magnitude of 9 and an
angular diameter of 2$^{\prime}$ for an observer on Earth. What would be the apparent magnitude of the night
sky of a planet orbiting a star exactly at the center of M57?
\begin{pssolution*}{}{}
    The night sky covers a solid angle of \(2\pi\) sr, while the solid angle we see is \(\Omega \approx \pi \theta^2\), where \(\theta\) is the angular radius. Since the flux is proportional to \(\Omega\),

    \[m_{ceu}-m = -2.5\log\left(\frac{\Omega_{ceu}}{\Omega}\right)\]

    To convert \(\Omega\) from arcmin\(^2\) to sr, multiply by \(\left(\frac{\pi}{180\cdot60}\right)^2\). Thus
    
    \[m_{ceu} = 9 - \log\left(\frac{2\pi}{\pi(\frac{\pi}{180\cdot60})^2}\right)\]

    \[\boxed{m_{ceu} = -9.43}\]
\end{pssolution*}
\end{pproblem}


\pts{4}
\begin{pproblem} (Adapted from Problem set 4 - 2021)
    The Cherenkov effect was first detected by the Soviet scientist
    Pavel Cherenkov in 1937. Later, together with his colleagues I. E. Tamm and
    I. M. Frank, he gave the phenomenon a physical interpretation and was awarded the Nobel Prize in Physics
    in 1958. Before a mathematical treatment, we first need to understand the principle a little better.

    When high-energy charged particles travel through a dielectric medium, atoms can be excited if
    the particle speed exceeds the phase velocity ($\frac{c}{n}$). Those atoms, in
    turn, emit electromagnetic radiation when they return to the ground state. The emitted waves
    spread spherically and, when superposed, form a cone of opening angle $2\alpha$,
    as the figure below shows.
    \begin{figure}[H]
      \centering
      \includegraphics[width=0.9\linewidth]{imagens/cherenkov.png}  
      \caption{Radiation mechanism of the Cherenkov effect}
    \end{figure}

    This effect is similar to a jet moving at supersonic speed: it follows the same
    principle as the Mach cone, but with light. We now develop the mathematical model of the Cherenkov effect.

    \begin{center}
        \textbf{Part A - Theoretical Model}    
    \end{center}
    
    Consider a particle moving at relativistic speeds through a medium of refractive index
    $n$. Its rest mass is $m_0$, and it has linear momentum $p$ and speed $v$. At a
    given moment, a photon is emitted at an angle $\alpha$, as Figure 1 shows.

    \begin{alternativas}
        \item Letting $\mu$ be the frequency of the emitted photon, determine its linear
        momentum, $p_\mu$, and its energy, $E_\mu$. Your answer must be in terms of $n$, $\mu$, and physical constants.
        \item Find an expression for the linear momentum of the particle after it emits
    the photon, in terms of $p_\mu$, $p$, and $\alpha$.
        \item Letting $\beta_n = \frac{c}{vn}$, prove that the relation below holds:
        \begin{equation}
            \cos\alpha = \frac{1}{\beta_n}
        \end{equation}
        \item Assuming that linear momentum and energy are conserved, determine the
        minimum speed at which the Cherenkov effect can occur.
        \textbf{Hint:} Compared with the other parameters, the factor $(n^2-1)h\mu$ may be neglected.

    \begin{center}
        \textbf{Part B - Nuclear Reactions}    
    \end{center}
    
    The proton-proton chain is a sequence of fusion reactions that convert hydrogen into helium.
    One possible branch of the proton-proton chain is $pp \text{ IV}$, in which, in principle, a
    helium-3 atom reacts directly with a proton, according to the reaction below:

    \begin{equation}
        ^3\text{He} + ^1\text{H} \rightarrow ^4\text{He} + \nu + \_\_
    \end{equation}

        \item  State the nuclear conservation law you use, and identify the missing particle
        in the blank of the reaction above.

        \item  State the nuclear conservation law you use, and identify the missing particle
        in the blank of the reaction below:

        \begin{equation}
            \pi^- \rightarrow \mu^- + \_
        \end{equation}
        
        \textbf{Data:} Pion mass: $140 \ \text{MeV}/c^2$, muon mass: $106 \ \text{MeV}/c^2$.
    \end{alternativas}

\begin{pssolution*}{}{}
    \begin{center}
        \textbf{Part A - Theoretical Model}    
    \end{center}

    \begin{alternativas}
        \item For a photon, the de Broglie relations give
        \[E = h\mu, \ \ p = \frac{h}{\lambda}\]
        where \(h\) is Planck's constant.

        Using the fundamental wave relation, \(\lambda = v/\mu = \frac{c}{n\mu}\), so

        \[\boxed{E = h\mu, \ \ p = \frac{nh\mu}{c}}\]

        \item Since the total momentum is conserved, suppose the particle moves with momentum \(p\) along the \(x\) axis before emitting a photon. Then
        
        \[p = p_\mu\cos\alpha + p'_x\]

        \[p_\mu\sin\alpha = p'_y\]

        \[p' = \sqrt{p_x^2 + p_y^2}\]

        Solving for \(p'\),

        \[p'^2 =(p - p_\mu\cos\alpha)^2 +p_\mu^2\sin^2\alpha \]

        \[p'^2 = p^2 - 2p_\mu p \cos\alpha + p_\mu^2\]

        Finally,

        \[p' = \sqrt{p^2-2pp_\mu\cos\alpha+p^2_\mu}\]

        \item Consider the following figure.
        
        \begin{figure}[H]
            \centering
            \includegraphics[width=0.7\linewidth]{imagens/cherenkov.1.png}
            \caption{Cherenkov diagram}
        \end{figure}

        The radius of the circle is the distance traveled by the photon. The distance from the center of the circle to the end of the cone is \(vt\), so

        \[\boxed{\cos\alpha = \frac{ct}{nvt} = \frac{c}{nv} \equiv \frac{1}{\beta_n}}\]

        \item Setting \(c=1\), energy conservation gives
        \[\sqrt{p^2+m_0^2} = \sqrt{p'^2+m_0^2}+h\mu\]

        Working through this expression,

        \[\sqrt{p'^2 + m_0^2} = \sqrt{p^2+m_0^2} - h\mu\]

        \[p'^2 + m_0^2 = p^2+m_0^2 + h^2\mu^2 - 2h\mu\sqrt{p^2+m_0^2}\]

        Substituting \(p'\),

        \[p^2-2pp_\mu\cos\alpha+p^2_\mu = p^2+h^2\mu^2 - 2h\mu\sqrt{p^2+m_0^2}\]

        Solving for \(\cos\alpha\),

        \[\cos\alpha = \frac{p_\mu^2 - h^2\mu^2+2h\mu\sqrt{p^2+m_0^2}}{2pp\mu}\]

        Simplifying,

        \[\cos\alpha = \frac{p_\mu}{2p} - \frac{h^2\mu^2}{2pp_\mu} + \frac{h\mu\sqrt{p^2+m_0^2}}{pp_\mu}\]
        
        Substituting \(p_\mu = nh\mu\),

        \[\cos\alpha = \frac{nh\mu}{2p} + \frac{h\mu}{2np}+ \frac{\sqrt{p^2+m_0^2}}{np}\]

        \[\cos\alpha = \frac{h\mu(n^2-1)+2\sqrt{p^2+m_0^2}}{2np}\]

        Neglecting the term \(h\mu(n^2-1)\) and substituting \(\cos\alpha\),

        \[\frac{1}{nv} = \frac{\sqrt{p^2+m_0^2}}{np}\]

        Solving for \(v\),

        \[v = \frac{p}{\sqrt{p^2+m_0^2}}\]

        Finally, restoring the factors of \(c\),

        \[\boxed{v = \frac{pc^2}{\sqrt{p^2c^2+m_0^2c^4}}}\]

        \begin{center}
            \textbf{Part B - Nuclear Reactions}    
        \end{center}

        \item Using charge conservation, the left-hand side has \(2+1=3\) protons, but the right-hand side has only \(2\), so the missing particle must contain one proton. The mass number is \(3+1=4\) on the left and \(4\) on the right. The missing particle must therefore have the charge of a proton and a mass much smaller than a proton's. The only particle that fits this description is the \(\boxed{\text{positron, }e^+}\).
            
        \item Lepton number is conserved. The pion is not a lepton, \(n_L=0\), but the muon is a lepton with \(n_L=+1\), so the missing particle must have \(n_L=-1\). The only neutral antilepton associated with the muon and carrying \(n_L=-1\) is the muon antineutrino, \(\bar{\nu}_\mu\). The complete reaction is therefore
        \[\boxed{\pi^{-}\rightarrow \mu^- + \bar{\nu}_\mu}{}\]

    \end{alternativas}

\end{pssolution*}
\end{pproblem}


\pts{3} 
\begin{pproblem}(Problem set 3 - Vinhedo 2022)
Juvelino, observing directly from his observatory in Paris, France, monitors the star Polaris ($\alpha\text{UMi}$). His goal is to find the color temperature $T_c$ of the star. Some of the data he has on his target are:

\begin{itemize}
    \item Apparent magnitude in the $V$ band: $V = 1.98$;
    \item Absolute magnitude in the $V$ band: $M_V = -3.60$;
    \item Absolute magnitude in the $B$ band: $M_B = -3.19$.
\end{itemize}

Using the information provided, help Juvelino!

\begin{alternativas}
    \item From many observations, Juvelino determined that the interstellar extinction in the $V$ band toward Polaris is $a_V = 5.8 \, \text{mag/kpc}$. Find the distance, in pc, from $\alpha\text{UMi}$ to Earth.
    \item Using the empirical relation
    \begin{equation}
        \frac{A_V}{E_{B-V}} = 3.0
    \end{equation}
    where $A_V$ is the total interstellar extinction in the $V$ band and $E_{B-V}$ is the $B-V$ color excess, find the $B-V$ color index of the observed star.
    \item Show the relation
    \begin{equation}
        T_c = \frac{7009}{(B-V) + 0.47}
    \end{equation}
    in which the color temperature is given in kelvin. To do so, use the fact that stars of spectral class $A0$ have $(B-V) = 0$ and $T_c = 15000 \, \text{K}$. Also use the wavelengths of the $B$ and $V$ bands, $\lambda_B = 440 \, \text{nm}$ and $\lambda_V = 548 \, \text{nm}$, respectively. Justify any approximations you make.

    \textbf{HINT: Planck's law may be useful}
    \item Find the color temperature of Polaris.
\end{alternativas}
\begin{pssolution*}{}{}
    \begin{alternativas}
        \item The expression that correctly relates the magnitudes is
        \[V - M_V = 5 \log d - 5 + a_V d\]

        This equation can be solved only by iteration:

        \[d = \frac{V-M_V - 5\log d + 5}{a_v}\]

        Iterating, we obtain \(\boxed{d\approx 100\text{ pc}}\).

        \item By definition, \(E_{B-V} = A_B - A_V\). Thus
        
        \[\frac{A_V}{A_B-A_V} = 3.00\]

        \[A_B = \frac{A_V}{3}+A_V = 0.77\]

        Thus

        \[B-M_B = 5\log d - 5 + A_B\rightarrow B = 2.58\]

        By definition, the \(B-V\) index is \(\boxed{U_{B-V} = B-V = 0.60}\)

        \item Using Planck's law,
        
        \[\frac{B_B}{B_V} = \frac{F_B}{F_V} \left(\frac{\lambda_V}{\lambda_B}\right)^5 \frac{e^{\frac{\beta hc}{\lambda_V}-1}}{e^{\frac{\beta hc}{\lambda_B}-1}}\]

        Using the Wien approximation,

        \[\frac{F_B}{F_V} = \left(\frac{\lambda_V}{\lambda_B}\right)^5 \frac{e^{\frac{\beta hc}{\lambda_V}}}{e^{\frac{\beta hc}{\lambda_B}}}\]

        Using Pogson's equation to compare the magnitudes,

        \[B-V = -2.5\log\left(\frac{F_B}{F_V}\right)+C\]

        A constant is included because the two bands have different wavelengths. Working through the expression,

        \[B-V = -2.5\log\left(\left(\frac{\lambda_V}{\lambda_B}\right)^5 \frac{e^{\frac{\beta hc}{\lambda_V}}}{e^{\frac{\beta hc}{\lambda_B}}}\right)+C\]

        Working through it further,

        \[B-V = -2.5\log\left(\frac{\lambda_V}{\lambda_B}\right)^5 + 2.5 \beta hc\left(\frac{1}{\lambda_B}-\frac{1}{\lambda_V}\right)\log e + C\]

        Using \(\beta = \frac{1}{k_B T_c}\) and substituting the values,

        \[B-V = -1.19 + \frac{7009}{T_c}+C\]

        For type \(A_0\), \((B_V)=0\) and \(T_c = 15000\)K. Thus

        \[-1.19 + \frac{7009}{15000} + C = 0 \rightarrow C = 0.72\]

        Thus

        \[B-V = -0.47 + \frac{7009}{T_c} \rightarrow \ \ \boxed{T_c = \frac{7009}{(B-V) + 0.47}}\]

        \item From the previous equation,
        
        \[\boxed{T_c = \frac{7009}{0.60 + 0.47} \approx 6600 \text{ K}}\]
    \end{alternativas}    

\end{pssolution*}
\end{pproblem}

\newpage

\section{Thermodynamics}
\pts{3}
\begin{pproblem}
    A galaxy contains on the order of \(10^{10}\) stars. Because of this immense number, we can model a galaxy as a cloud of ideal gas, in which each star is equivalent to a particle of the gas.
    
    The goal of this question is to use this theoretical model to study some properties of galaxies. To that end, we make the following assumptions:

    \begin{enumerate}[label=\roman*)]
        \item The galaxy is spherical and in hydrostatic equilibrium.
        \item The mass density of the galaxy is constant and equal to \(\rho\).
        \item The stellar masses are small enough that interstellar interactions can be neglected.
    \end{enumerate}
    
    \begin{alternativas}
        \item Considering an ideal-gas system, find an expression for the pressure in terms of the density \(\rho\), the temperature \(T\), the mass of each particle \(\mu\), and physical constants.
        
        \item In our theoretical model it does not make sense to think in terms of temperature, so we need a substitute for it. Using the equipartition theorem, find an expression for \(T(r)\) and \(P(r)\).
    \end{alternativas}
\begin{pssolution*}{}{}
    \begin{alternativas}
        \item Starting from the Clapeyron equation,
        \[PV = Nk_BT\]

        Multiply both sides by \(\mu\),

        \[P\mu V = N\mu k_B T\]

        \[P\mu = \frac{N\mu}{V}k_BT\]

        Note that \(N\mu\) equals the total mass, so

        \[\rho = \frac{N\mu}{V}\]
        
        With this we obtain

        \[\boxed{P\mu = \rho k_B T}\]

        This expression will be used quite often in the thermodynamics part of this set.

        \item By the equipartition theorem,
        
        \[\frac{\mu v^2}{2} = \frac{3k_B T}{2}\]

        Hence
    
        \[T = \frac{\mu v^2}{3k_B}\]

        Since the stars have mass and speed, this is an acceptable substitution. We can compute \(v\) from the vis-viva equation, assuming circular orbits,

        \[v^2 = \frac{GM}{r^2} = \frac{4\pi G\rho r}{3}\]

        Thus

        \[\boxed{T(r) = \frac{4\pi G \rho \mu r}{9k_B}}\]

        For \(P(r)\), substitute into the formula:

        \[\boxed{P(r) = \frac{\rho k_B T(r)}{\mu} = \frac{4\pi G \rho^2 r}{9}}\]
    \end{alternativas}
    
\end{pssolution*}
\end{pproblem}

\pts{2}
\begin{pproblem}
    In this question we study the adiabatic index of stars. Assume that the exterior of a star is a vacuum at temperature \(T=0\).
    \begin{alternativas}
        \item All stars are bodies in hydrostatic equilibrium. Given that, what is the pressure at the surface of a star of mass \(M\) and radius \(R\)?
        \item Now suppose that the star expands by \(\delta R\). How does the pressure change? If needed, use \((1+x)^n\approx 1+nx\).
        \item Now conclude what minimum value of \(\gamma\), \(\gamma_{min}\), the star must have in order to remain gravitationally bound (assume that it expands adiabatically).
    \end{alternativas}

\begin{pssolution*}{}{}
    \begin{alternativas}
        \item Using \(P = F/A\), we have
        \[P(R) = \frac{GM^2}{4\pi R^4}\]

        \item Thus
        \[P(R+\delta R)  = \frac{GM^2}{4\pi}(R+\delta R)^{-4} = \frac{GM^2}{4\pi R^4}\left(1+\frac{\delta R}{R}\right)^{-4}\]

        Using the approximation given in the statement,

        \[\boxed{P(R+\delta R) = \frac{GM^2}{4\pi R^4}\left(1-\frac{4\delta R}{R}\right) = P\left(1-\frac{4\delta R}{R}\right)}\]
    
        \item In an adiabatic expansion, \(PV^\gamma\) is constant. Since \(V\propto R^3\), it follows that \(PR^{3\gamma}\) is constant. Thus
        
        \[PR^{3\gamma} = (P+\delta P)(R + \delta R)^{3\gamma}\]

        Using the same reasoning as in the previous item,
        \[PR^{3\gamma} = (P+\delta P)R^{3\gamma}\left(1 + \frac{\delta R}{R}\right)^{3\gamma}\]

        \[P = (P+ \delta P)\left(1+\frac{3\gamma\delta R}{R}\right)\]

        Substituting \(P+\delta P\),

        \[1 = \left(1-\frac{4\delta R}{R}\right)\left(1+\frac{3\gamma\delta R}{R}\right)\]

        Hence

        \[\left(1-\frac{4\delta R}{R}\right)^{-1} = 1+\frac{3\gamma\delta R}{R}\]

        Using the binomial approximation again on the left-hand side and solving for \(\gamma\), we obtain

        \[\boxed{\gamma = \frac{4}{3}}\]
    \end{alternativas}
    
\end{pssolution*}
\end{pproblem}

\pts{4}
\begin{pproblem}
    There are several models for the atmosphere of our planet. Let us explore them and find the physical effects of each one.
    \begin{alternativas}
        \item First, consider the isothermal model (\(T=\) const). Given that each air particle has mass \(\mu\) and that the atmosphere has \(P(0) = P_0\), find a formula for the pressure as a function of height, \(P(h)\).
        \item A more realistic model of the atmosphere is in fact adiabatic, since air is a very poor conductor of heat. Given that the air has Poisson coefficient \(\gamma\), find an expression for \(P(h)\) in the adiabatic model. Assume that at sea level the pressure and the temperature are \(P(0)\) and \(T(0)\).
        \item Find an expression for \(dT/dh\) for the previous model and estimate its value. Is the result consistent with reality?
    \end{alternativas}

\begin{pssolution*}{}{}
    \begin{alternativas}
        \item From the relation
        \[P\mu = \rho k_BT\]

        Thus, since \(T = \mathrm{const}\),

        \[\frac{P(h)}{P(0)} = \frac{\rho(h)}{\rho(0)}\]

        The pressure gradient is given by Stevin's law,

        \[\frac{dP(h)}{dh} = -g\rho(h)\]

        Substituting \(\rho(h)\),

        \[\frac{dP(h)}{dh} = \frac{-gP(h)\rho(0)}{P(0)}\]

        Separating the terms,

        \[\frac{dP(h)}{P(h)} = \frac{-g\rho(0)}{P(0)}dh\]

        Integrating both sides,

        \[\int_0^h\frac{dP(h)}{P(h)} = \ln\left(\frac{P(h)}{P(0)}\right)\]

        \[\int_0^hdh = h\]

        Thus

        \[P(h) = P(0)e^{-\frac{g\rho(0)}{P(0)}h}\]

        Finally, we eliminate the term \(\rho(0)\) using the Clapeyron relation,

        \[P(0)\mu = \rho(0)k_B T \rightarrow \rho(0) = \frac{P(0)\mu}{k_BT}\]

        Substituting,

        \[\boxed{P(h) = P(0)e^{-\frac{\mu g h}{k_BT}}}\]

        \item In the adiabatic model we have
        \[PV^\gamma = \mathrm{const}\]

        Since \(V\propto \rho^{-1}\),

        \[P\rho^{-\gamma} = \mathrm{const}\]

        That is,

        \[\rho(h) = \left(\frac{P(0)}{P(h)}\right)^{-1/\gamma}\rho(0) = \left(\frac{P(h)}{P(0)}\right)^{1/\gamma}\rho(0)\]

        Substituting into the expression for the pressure gradient,

        \[\frac{dP(h)}{dh} = -g\rho(h) = -\left(\frac{P(h)}{P(0)}\right)^{1/\gamma}\rho(0)g\]

        Separating the terms and integrating,

        \[\int_0^h P(h)^{-1/\gamma}dP(h) = -P(0)^{-1/\gamma}\rho(0)g\int_0^h dh\]

        \[\frac{\gamma}{\gamma-1}(P(h)^{\frac{\gamma-1}{\gamma}}-P(0)^{\frac{\gamma-1}{\gamma}}) = -\frac{\mu g z}{k_B T(0)}P(0)^{\frac{\gamma-1}{\gamma}}\]

        Here we use

        \[P(0)\mu = \rho(0)k_B T(0)\]

        Solving for \(P(h)\),
    
        \[\boxed{P(h) = P(0)\left(1-\frac{\gamma-1}{\gamma}\frac{\mu g h}{k_B T(0)}\right)^{\frac{\gamma}{\gamma-1}}}\]
       
        \item For adiabatic processes,
        
        \[PV^\gamma = \mathrm{const}\]

        \[V\propto \frac{T}{P}\]

        Hence

        \[P^{1-\gamma} T^\gamma = \mathrm{const}\]

        \[T P^\frac{1-\gamma}{\gamma} = \mathrm{const}\]
        
        Thus

        \[T = P(0)^\frac{1-\gamma}{\gamma}T(0)P^{\frac{\gamma-1}{\gamma}}\]

        Differentiating,

        \[\frac{dT}{dh} = P(0)^\frac{1-\gamma}{\gamma}T(0) \frac{\gamma-1}{\gamma P^{1/\gamma}}\frac{dP}{dh}\]

        From the previous item we have the expression for \(P\) as a function of \(h\). Carrying out the derivative,

        \[\frac{dP}{dh} = P(0)\frac{d}{dh}\left(1-\frac{\gamma-1}{\gamma}\frac{\mu g h}{k_B T(0)}\right)^{\frac{\gamma}{\gamma-1}}\]

        Using the chain rule,

        \[\frac{dP}{dh} = \frac{\gamma}{\gamma-1}P(0)\left(1-\frac{\gamma-1}{\gamma}\frac{\mu g h}{k_B T(0)}\right)^{\frac{1}{\gamma-1}}\left(-\frac{\gamma-1}{\gamma}\frac{\mu g}{k_B T(0)}\right)\]

        Simplifying,

        \[\frac{dP}{dh} = -P(0)\left(1-\frac{\gamma-1}{\gamma}\frac{\mu g h}{k_B T(0)}\right)^{\frac{1}{\gamma-1}}\frac{\mu g}{k_B T(0)}\]

        Substituting into the expression for \(dT/dh\),

        \[\frac{dT}{dh} = -P(0)^\frac{1-\gamma}{\gamma}T(0) \frac{\gamma-1}{\gamma P^{1/\gamma}}P(0)\left(1-\frac{\gamma-1}{\gamma}\frac{\mu g h}{k_B T(0)}\right)^{\frac{1}{\gamma-1}}\frac{\mu g}{k_B T(0)}\]

        \[\frac{dT}{dh} = -\frac{\gamma-1}{\gamma}P(0)^\frac{1}{\gamma} P^{-\frac{1}{\gamma}}\left(1-\frac{\gamma-1}{\gamma}\frac{\mu g h}{k_B T(0)}\right)^{\frac{1}{\gamma-1}}\frac{\mu g}{k_B}\]

        Substituting \(P\),

        \[\frac{dT}{dh} = -\frac{\gamma-1}{\gamma}P(0)^\frac{1}{\gamma} \left(P(0)\left(1-\frac{\gamma-1}{\gamma}\frac{\mu g h}{k_B T(0)}\right)^{\frac{\gamma}{\gamma-1}}\right)^{-\frac{1}{\gamma}}\left(1-\frac{\gamma-1}{\gamma}\frac{\mu g h}{k_B T(0)}\right)^{\frac{1}{\gamma-1}}\frac{\mu g}{k_B}\]

        Continuing to simplify, we obtain a neat expression,

        \[\frac{dT}{dh} = -\frac{\gamma-1}{\gamma}\frac{\mu g}{k_B}\]

        We can approximate air as a diatomic gas, \(\gamma = \frac{7}{5}\), of mass \(\mu\approx 30\)u. Carrying out the calculation, we obtain

        \[\boxed{\frac{dT}{dh} \approx - 10.03 \text{ K/km}}\]

        which is a value quite consistent with reality.
    \end{alternativas}
\end{pssolution*}
\end{pproblem}

\pts{5}
\begin{pproblem}
    Black holes are among the most fascinating bodies in the universe. In this question we study some of the thermodynamics associated with these objects. For this question we use natural units, i.e.\ \(c = G = \hbar = k_B = 1\). In this convention, the mass of the black hole is described by the equation

    \[dM = \frac{\kappa}{8\pi}dA + \Omega dL + \Phi dQ\]

    Here the quantities refer to the event horizon: \(\kappa\) is the gravitational acceleration at the event horizon, \(A\) its area, \(\Omega\) the angular velocity of the black hole, \(L\) its moment of inertia, \(\Phi\) the electric potential, and \(Q\) its charge.

    \begin{center}
    \textbf{Part 1: Basic Thermodynamics}   
    \end{center}
    
    \begin{alternativas}
        \item One of the fundamental concepts of thermodynamics is entropy. Using what you know about it, briefly explain the inequality
        \[\oint dS \geq 0\]

        \item Of the 3 factors that govern the mass of a black hole (\(dA, \ dL, \ dQ\)), only \(dA\) obeys the same inequality as the entropy. Why is this always true?
        
        For the next items, consider a black hole with no spin and no charge.

        \item Bekenstein and Hawking proved the so-called \textit{Bekenstein-Hawking entropy}, which relates the entropy to the area of the black hole (remember that we are using natural units, so some dimensions may not make sense). Bekenstein and Hawking found that for a black hole \(S \equiv \frac{A}{4}\). The equation we then have is
        
        \[dM = \frac{\kappa}{2\pi}d\left(\frac{A}{4}\right)\]

        By analogy between \(dM\) and some state function, find the temperature of the black hole as a function of \(\kappa\).

        An important term is still missing from the previous formula: the pressure associated with the dark-energy density, \(\Lambda\).

        \item The pressure due to dark energy is given by
        \[P = -\frac{\Lambda}{8\pi}\]

        where \(\Lambda\) is constant. This shows that \(VdP\) vanishes, so it may be added freely to the previous expression. We can therefore conclude that the mass of the black hole is in fact related to another thermodynamic potential. Which one is it?

        \item Note that the volume, \(V=V\), and the entropy, \(S = A/4\), are no longer independent for black holes. Assuming that the event horizon of the black hole is a sphere, find a relation between \(S\) and \(V\). This is further evidence that the internal energy, \(U = U(S, V)\), is not the best thermodynamic potential to work with.
    \end{alternativas}

    \centering
    \textbf{Part 2: Carnot Cycle for Black Holes}
    
    \begin{alternativas}
        \item The goal of this part is to build a theoretical model for a Carnot cycle inside black holes. First, however, prove an important result: for black holes, adiabatic and isochoric processes must be equivalent.
        
        \item Calculate the heat capacity at constant pressure of a black hole. Your result should be something bizarre.
        
        \item Use the fact that \(Q = T \Delta S\) along the isotherms, together with the results of the previous parts, to calculate the efficiency of a black-hole Carnot engine and confirm that you
        obtain the Carnot efficiency. Marvel at how much faster this calculation is than the
        typical derivation of the Carnot efficiency, and note that you have also inadvertently
        calculated the efficiency of the Stirling cycle.

        If you are interested in the subject, there is an interesting paper that discusses cycles in black holes specifically, available \href{https://arxiv.org/pdf/1404.5982}{here}.
    \end{alternativas}

    \centering
    \textbf{Part 3: Lifetime and Evaporation of Black Holes}
    
    \begin{alternativas}
        \item As calculated in part 1, black holes have a temperature. As a result, they emit radiation as described by the Stefan-Boltzmann law. Given that, find a relation between the lifetime of a black hole and its mass \(M\).
    \end{alternativas}

    \begin{pssolution*}{}{}
        \begin{center}
            \textbf{Part 1: Basic Thermodynamics}
        \end{center}
        \begin{alternativas}
            \item The inequality represents the second law of thermodynamics. Briefly, it shows that the change in entropy is always positive, that is, everything tends toward disorder.
            
            \item Note that no law of physics prevents a black hole from losing angular velocity, \(dL < 0\), or charge, \(dQ<0\). When we examine the area of the black hole, however, something interesting appears. Consider the event horizon, given by a sphere of radius \(R_{sch}\). Since this radius is defined as the distance at which an object moving at the speed of light can no longer escape the gravitational attraction of the black hole, we can obtain it from energy conservation,
            
            \[-\frac{GMm}{R_{sch}} + \frac{mc^2}{2} = 0 \rightarrow R_{sch} = \frac{2GM}{c^2}\]
            
            Hence the area of the event horizon is given by

            \[A = 4\pi R_{sch}^2 = \frac{8\pi G^2M^2}{c^4}\]

            In natural units,

            \[A = 8\pi M^2\]

            Thus

            \[dA = 16\pi MdM\]

            A black hole cannot lose mass, since nothing can escape the event horizon. Thus \(dM>0\) and consequently \(dA>0\). We can therefore conclude that it is also valid that
            
            \[\oint dA \geq 0\]

            for black holes.

            \item From the mass-energy relation,
            
            \[U = M \rightarrow dU = dM\]

            (Recall that we are using natural units, so \(Mc^2\) = \(M\).) Thus, using the first law of thermodynamics,

            \[dU = TdS - PdV\]

            Hence

            \[T = \left.\frac{dU}{dS}\right\vert_V = \left.\frac{dM}{dS}\right\vert_V\]

            Using the entropy formula given in the statement (analogous to the one obtained in the previous item),

            \[dM = \frac{\kappa}{2\pi}\left(\frac{A}{4}\right) = \frac{\kappa}{2\pi}dS\]

            Using the equation from the previous item, we see that \(\kappa = 1/4M\). Substituting into the formula for \(T\),

            \[\boxed{T = \frac{1}{8\pi M}}\]

            \item Using the formula obtained for \(T\),
            
            \[dM = TdS\]

            Adding the vanishing term \(VdP\), we have

            \[dM = TdS + VdP\]

            Note that this is identical to the enthalpy, \(H\):

            \[H = U + PV \rightarrow dH = (TdS - PdV) +  (PdV - VdP) = TdS - VdP\]

            We can therefore say that the mass is equivalent to the enthalpy.

            \item The volume of a sphere is given by
            
            \[V = \frac{4\pi R^3}{3}\]

            while the area is given by \(A = 4\pi R^2\), so \(R = \sqrt\frac{A}{4\pi}\equiv \sqrt\frac{S}{\pi}\). Substituting into the volume formula,
            
            \[V = \frac{4\pi}{3}\left(\frac{S}{\pi}\right)^{3/2}\]

            Simplifying,

            \[\boxed{V = \frac{4}{3\sqrt{\pi}}S^{3/2}}\]

            That is, the volume and the entropy depend on each other. Since the internal energy is a function of the volume \textbf{and} the entropy, it does not make sense to say that it equals the mass of the black hole.
        \end{alternativas}
    
        \begin{center}
            \textbf{Part 2: Carnot Cycle for Black Holes}
        \end{center}

        \begin{alternativas}
            \item Using the relation from the previous item,
            
            \[dV = \frac{2}{\pi}S dS\]

            That is, an adiabatic cycle (\(dS = 0\)) is equivalent to an isochoric cycle (\(dV = 0\)).

            \item To calculate the heat capacity at constant pressure,
            
            \[C_P = \left.\frac{\partial Q}{\partial T}\right\vert_P = T\left.\frac{\partial S}{\partial T}\right\vert_P = \frac{dM}{dT}\] 
            
            Using the formula for the temperature,

            \[T = \frac{1}{8\pi M} \rightarrow M = \frac{1}{8\pi T}\]
            \[\boxed{C_P = \frac{dM}{dT} = -\frac{1}{8\pi T^2}}\]

            \item The efficiency of a cycle is given by

            \[\eta = \frac{W}{Q_H} = \frac{Q_C-Q_H}{Q_H} = 1-\frac{Q_C}{Q_H}\]

            where \(Q_H\) is the heat that enters the cycle during the hot phase and \(Q_C\) is the amount that enters it during the cold phase. Define \(T_H\) and \(T_C\) as the hot and cold temperatures, respectively. Suppose the cycle operates with \(1-2\) at the hot temperature and \(3-4\) at the cold temperature. Then

            \[Q_H = T_H\Delta S_{1\rightarrow 2}\]

            \[Q_C = T_C\Delta S_{3\rightarrow 4}\]

            Using the expression from item \(1.e\), we have

            \[S = \frac{3\sqrt{\pi}}{4}V^{2/3}\]

            Thus

            \[\Delta S_{1\rightarrow 2} = \frac{3\sqrt\pi}{4}\left(V^{2/3}_2 - V^{2/3}_1\right)\]
            \[\Delta S_{3\rightarrow 4} = \frac{3\sqrt\pi}{4}\left(V^{2/3}_4 - V^{2/3}_3\right)\]

            Hence the efficiency can be written as

            \[\eta = 1 - \frac{T_C\left(V^{2/3}_4 - V^{2/3}_3\right)}{T_H \left(V^{2/3}_2 - V^{2/3}_1\right)}\]

            But since, for black holes, adiabatic and isochoric processes are identical, \(V_1=V_3\) and \(V_2=V_4\), so

            \[\boxed{\eta = 1 -\frac{T_C}{T_H}}\]

            which is the Carnot efficiency!

            \begin{center}
                \textbf{Part 3: Lifetime and Evaporation of Black Holes}
            \end{center}
        \end{alternativas}
        \begin{alternativas}
            \item The Stefan-Boltzmann law says that the luminosity (power) is given by
            \[L = -\frac{dE}{dt} = - \frac{dM}{dt} = A\sigma T^4 \]

            Note that \(A\propto M^2\) and \(T\propto 1/M\), so

            \[\frac{dM}{dt} \propto M^2\left(\frac{1}{M}\right)^4 = \frac{1}{M^2}\]

            Separating the variables,

            \[M^2dM \propto dt\]

            Finally, integrating, we obtain

            \[\boxed{t \propto M^3}\]
        \end{alternativas}
    \end{pssolution*}


\end{pproblem}

\pts{3}
\begin{pproblem}
    Consider a rocket that uses an ideal diatomic gas as fuel. Its mechanism is quite simple: the gas starts in a chamber at temperature \(T_1\) with cross-sectional area \(A_1\); the gas then flows adiabatically and is expelled through an opening of area \(A_2\), at pressure \(P_2\) and temperature \(T_2<T_1\). Assuming that the flow is continuous, determine the thrust felt by the rocket.
\begin{pssolution*}{}{}
    Since the process is adiabatic, we have

    \[PV^\gamma \propto P\left(\frac{T}{P}\right)^\gamma = \mathrm{const}\]

    Since the gas is diatomic, \(C_P = 7/2 R\) and \(C_V = 5/2 R\), and by definition \(\gamma = C_P/C_V = 7/5\).

    Thus we have

    \[P_1\left(\frac{T_1}{P_1}\right)^\gamma = P_2\left(\frac{T_2}{P_2}\right)^\gamma\]

    Substituting the value of \(\gamma\),

    \[\frac{P_1}{P_2} = \left(\frac{T_1}{T_2}\right)^{7/2}\]

    Let \(v_1\) and \(v_2\) be the speed of the gas at the given stages. Since the flow is continuous, we must have

    \[\rho_1v_1A_1 = \rho_2v_2A_2\]

    And by the ideal gas law,

    \[P\mu = \rho R T \rightarrow \rho \propto \frac{P}{T}\]

    Substituting into the previous expression,

    \[\frac{P_1v_1A_1}{T_1}=\frac{P_2A_2A_2}{T_2}\]

    Combining this with the first equation, we have
    
    \[\frac{v_1}{v_2} = \frac{A_2}{A_1}\left(\frac{T_2}{T_1}\right)^{5/2}\]

    Using Bernoulli's equation for gases,

    \[\frac{1}{2}\mu v_1^2 + C_P RT_1 = \frac{1}{2}\mu v_2^2 + C_P T_2\]

    Solving for \(v_2\),

    \[v_2^2 = \frac{7R(T_1-T_2)}{\mu(1-(A_1/A_2)^2(T_2/T_1)^5)}\]

    Finally, using Newton's second law,

    \[F = \frac{dp}{dt} = \frac{dm}{dt}v_2 = \rho_2A_2 v_2^2 = \boxed{\frac{7P_2A_2(T_1-T_2)}{T_2(1-(A_1/A_2)^2(T_2/T_1)^5)}} \]

\end{pssolution*}

\end{pproblem}


\pts{5} 
\begin{pproblem} (Adapted from Iran Problem Set)
    This problem aims to calculate the boiling point of liquids.

Particles in a liquid move at different speeds depending on the temperature, and some of these particles can escape the intermolecular forces and Earth's gravity (which will be neglected in this problem), leaving the surface of the liquid. These particles transfer their momentum, creating pressure when they collide with the surrounding environment. This pressure is known as the vapor pressure of the liquid. The boiling point is the temperature at which the vapor pressure equals the atmospheric pressure around the liquid.

\begin{alternativas}
\item Using the Maxwell-Boltzmann distribution, find a relation for the root-mean-square speed $v_{\text{rms}}$.

The Maxwell-Boltzmann distribution is
\begin{equation}
    n(v) \, dv = n \left(\frac{m}{2\pi k_B T}\right)^{3/2} e^{-\frac{mv^2}{2k_BT}} \, 4\pi v^2 dv
\end{equation}

The root-mean-square speed is given by
\begin{equation}
    v_{\text{rms}}^2 = \frac{1}{N} \sum_{i=1}^N v_i^2 = \frac{1}{n} \int_0^\infty v^2 n(v) dv
\end{equation}

\item Suppose that the speed of the particle under study equals $v_{\text{rms}}$. Also suppose that Earth's atmosphere is 80\% nitrogen and 20\% oxygen, and that the liquid under study is water.

Calculate the distance the escaped particle travels in the atmosphere before colliding, known as the mean free path. Express this distance in terms of the number density of the atmosphere and the geometric cross section of the particles.

\item Calculate the rate of change of the momentum of an escaping particle. Divide the change in momentum by the time interval of the process. Assume that the particles of the liquid lose all of their momentum when they collide with air molecules.

The mean time until the next collision is the mean free path divided by the speed of the particle.

Given that, over this time interval, a momentum equal to the momentum of the liquid particle is transferred to the air molecule, use Newton's second law to calculate the force exerted by the liquid particle on the air molecule.

\item Pressure is the force exerted on a surface. Consider that the cross section of the force exerted on the air molecules equals the geometric cross section of the liquid molecules. Find an expression for the vapor pressure of a liquid.

Hint:
\begin{equation}
    P = n \frac{S_a}{S_i} \, 3k_BT
\end{equation}

where $S_a$ and $S_i$ are the geometric cross sections of the air molecules and of the water molecules, respectively. $n$ is the number density of the air molecules near the surface of the liquid, and $T$ is the temperature of the liquid.

\item Using the hydrostatic-equilibrium relation and assuming constant gravitational acceleration, constant air density, and zero pressure in the upper layers of the atmosphere, find a relation for the pressure near Earth's surface. Express this relation in terms of the air density, the gravitational acceleration, and the thickness of the atmosphere.

\item Impose the condition required for boiling by equating the atmospheric pressure near Earth's surface (obtained above) to the vapor pressure. Simplify the result until you obtain
\begin{equation}
    T = \frac{mhg S_i}{3k_B S_a}
\end{equation}

where $m$ is the mean mass of the air molecules, $g$ is the gravitational acceleration, $h$ is the thickness of the atmosphere, and the other parameters were described in the previous parts.

\item Determine the mean mass of the air molecules for the composition mentioned at the beginning of the problem.

\item Use the Bohr radius to estimate the ratio of the cross sections. Suppose an electron in a circular orbit around a proton, with the Coulomb force as the dominant force. Using Niels Bohr's assumption $L = n\hbar$, find the distance from the electron to the nucleus in terms of physical constants, $n$, and the atomic number $Z$.

\item Calculate the cross section of the air atoms and of the liquid atoms. Assume that each liquid molecule consists of two hydrogen atoms and one oxygen atom, and that the air particles consist of two oxygen atoms and two nitrogen atoms. Find the ratio of the cross sections $S_i/S_a$.

\item Take $g = 9.8 \, \text{m/s}^2$ and the height of the atmosphere to be $100 \, \text{km}$. Determine the boiling point of water.

\end{alternativas}

\begin{pssolution*}{}{}
    \begin{alternativas}
    \item Defining the variables,
    \[C = 4\pi n \left(\frac{m}{2\pi k_B T}\right)^{3/2}, \ \ \beta = \frac{1}{k_BT}\]

    we can write

    \[n(v)dv = C e^{-\frac{\beta mv^2}{2}}v^2dv\]

    Using the expression for \(v_{rms}\) given in the question,

    \[v_{rms}^2 = \frac{C}{n}\int_0^\infty  e^{-\frac{\beta mv^2}{2}}v^4 dv\]

    To evaluate the integral, we have

    \[\int_0^{\infty} e^{-av^2}v^4 dv = \frac{d^2}{da^2}\int_0^{\infty} e^{-av^2}dv\] 

    where

    \[a = \frac{\beta m}{2}\]
    The integral on the right-hand side is well known and equals \(\frac{1}{2}\sqrt{\frac{\pi}{a}}\). That is,

    \[\int_0^{\infty} e^{-av^2}v^4 dv = \frac{\sqrt{\pi}}{2}\frac{d^2}{da^2}a^{-1/2} = \frac{3}{8}\sqrt{\frac{\pi}{a^5}}\]

    Thus

    \[v_{rms}^2 = \frac{3C}{8n}\sqrt{\frac{\pi}{a^5}}\]

    Substituting \(C\) and \(a\), we have

    \[v_{rms}^2 = \frac{3}{2}\pi  \left(\frac{\beta m}{2\pi}\right)^{3/2}\sqrt{\frac{2^5\pi}{\beta^5m^5}}\]
    
    Finally, simplifying,

    \[\boxed{v_{rms}^2 = \frac{3}{m\beta} = \frac{3k_B T}{m}}\]

    \item Consider \(P(t)\) the probability that the particle does \textbf{not} collide during a time \(t\). The probability that the particle does not collide during a time \(t\) and does not collide during the following time \(dt\) is \(P(t+dt) = P(t)P(dt)\) (multiplicative property), but we also have
    
    \[P(t+dt) \approx P + \frac{dP}{dt}dt\]

    In a time \(dt\) the particle sweeps out a volume \(dV = \sigma v dt\), where \(\sigma\) is the impact parameter. The probability that the particle collides in a time \(dt\) is then \(P'(dt) = ndV = n\sigma v dt\), so (recalling that \(P(dt)\) is the probability that the particle does \textbf{not} collide) \(P(dt) = 1-n\sigma vdt\).

    Equating the expressions for \(P+dt\), we have

    \[P + \frac{dP}{dt}dt = P(1-n\sigma v dt)\]

    Working with this expression,

    \[dP = -Pn\sigma vdv \rightarrow \frac{dP}{P} = -n\sigma v dt\]

    Integrating both sides and using \(P(0)=1\), we have

    \[P(t) = e^{-n\sigma vt}\]

    Then the probability that a particle survives for a time \(t\) and then collides in the following time \(dt\) is

    \[J(t)dt = P(t)(1-P(dt)) = e^{-n\sigma v t} n\sigma v dt\]

    Thus the mean time between collisions is

    \[\tau = \int_0^\infty t J(t)dt = \int_0^\infty e^{-n\sigma v t} n\sigma v dt = \frac{1}{n\sigma v}\]

    The mean free path is then given by

    \[\lambda = v_{radial}\tau = \frac{1}{\sqrt{2}n\sigma}\]

    To find \(\sigma\), define \(r_H\) as the radius of the water molecule, and similarly \(r_O\) and \(r_N\). Looking at the following diagram,

    \begin{figure}[H]
        \centering
        \includegraphics[width= 0.7\linewidth]{imagens/secaotransversal.png}
        \caption{Source: Concepts in Thermal Physics, Blundell and Blundell}
    \end{figure}

    For nitrogen, \(\sigma_N = \pi(r_H + r_N)^2\), and for oxygen, \(\sigma_O = \pi (r_H + r_O)^2\). Since the atmosphere consists of \(0.8 N\) and \(0.2 O\), it holds that

    \[\sigma = 0.2\sigma_O + 0.8\sigma N\]

    Thus we have

    \[\boxed{\lambda = \frac{1}{\sqrt{2}n(0.8\sigma_0 + 0.2\sigma_0)}}\]

    \item Assuming that the particle loses all of its momentum when it collides with an air molecule, we have \(\delta p = mv_{rms}\). Each collision occurs in a time \(\delta t =\tau\), so
    \[F = \frac{\delta p }{\delta t} = \frac{mv_{rms}}{\tau} = n\sigma v_{rms}^2 \  \boxed{ = 3n\sigma k_BT}\]

    \item Using the definition of pressure,
    
    \[P = \frac{F}{\sigma} = \ \boxed{3nk_B T} \]

    \item Using the hydrostatic-equilibrium equation,
    
    \[\frac{dP}{dr} = -\rho g\]

    Since \(\rho\) and \(g\) are constant,

    \[P = -\rho g \int_h^0dh= \ \boxed{\rho g h}\]

    where \(h\) is the thickness of the atmosphere.

    \item Equating the pressures,

    \[\rho g h = 3\frac{S_a}{S_w}nk_BT\]

    Solving for \(T\), we have

    \[T = \frac{\rho g h}{3nk_B}\frac{S_w}{S_a}\]

    Since \(n\) is the number density of particles and \(\rho\) is the mass density, we have \(\rho/n = m\), so

    \[\boxed{T= \frac{m g h}{3k_B}\frac{S_w}{S_a}}\]

    \item Taking a weighted average,
    
    \[\boxed{m = 0.8 m_N + 0.2 m_O}\]

    \item The force felt by the electron is given by
    
    \[F = \frac{Ze^2}{4\pi \epsilon_0r^2}\]

    Equating this to the centripetal force,

    \[\frac{Ze^2}{4\pi \epsilon_0r^2} = m_e \omega^2 r\]

    By the definition of angular momentum, we have

    \[L = m_er^2\omega \rightarrow \omega = \frac{L}{m_er^2} = \frac{n\hbar}{m_er^2}\]

    Substituting into the previous expression,

    \[\frac{Ze^2}{4\pi\epsilon_0 r^2} = m_e \left(\frac{n\hbar}{m_er^2}\right)^2 r\]

    Solving for \(r\),

    \[r = \frac{4\pi\epsilon_0 n^2\hbar ^2}{Z e^2 m_e}\]

    \item For water we have \(r_w = 2r_H + r_O\), and for air \(r_a = 2(r_N + r_O)\). Since \(S_i \propto r_i^2\), we have
    
    \[\frac{S_w}{S_a} = \frac{(2r_H + r_O)^2}{4(r_N + r_O)^2}\]

    Assuming that all the electrons are in the lowest level (\(n=1\)), that \(Z_H = 1, \ Z_N = 7, \ Z_O = 8\), and that \(r\propto 1/Z\), we have
    
    \[\boxed{\frac{S_w}{S_a}=\frac{(2/Z_H + 1/Z_O)^2}{4(1/Z_N + 1/Z_O)^2} \approx 15.73}\]

    \item Substituting into the formula,

    \[T = \frac{(2\cdot 14 + 2 \cdot 16) \cdot 1.67 \cdot 10^{-27} \cdot 10^{5} \cdot 15.73}{3 \cdot 1.38\cdot 10^{-23}} \ \boxed{\approx 380\text{ K}}\]

    which is an estimate quite consistent with reality, \(T = 373\) K.
\end{alternativas}
    

    

    
    
\end{pssolution*}
\end{pproblem}

\pts{3}
\begin{pproblem} (Adapted from Iran Problem Set)
    Dr. Shahram Abbassi is one of the most renowned Iranian scientists in the field of accretion disks. In one of his recent studies of the giant molecular cloud B32, he found a Sun-like star at the center of that cloud. According to his research, the cloud has a mass of $10^6 M_\odot$, a radius of $30 \, \text{pc}$, and a very high viscosity, so high that if the cloud were to collapse, the whole system would collapse with spherical symmetry.

    What matters most is to calculate the specific heat at constant volume ($C_v$) of this cloud. Using the data given and reasonable approximations, determine a bound on $C_v$ such that accretion is possible. Do these values depend on the mass and the radius of the cloud? What can we conclude from the result?
\begin{pssolution*}{}{}
    This bound is given by the virial theorem, when

    \[\langle U \rangle = -2 \langle K \rangle\]

    We can model the cloud as a gas, so that its kinetic energy is given by \(\frac{3}{2} N k_B T\). Its potential energy is the potential energy of a sphere, given by \(U = -3GM^2/5R\). Thus

    \[\frac{3GM^2}{5R} = 3\frac{M}{m_p}k_B T\]

    where \(m\) is the mass of a particle in the cloud.

    Thus

    \[T = \frac{GMm_p}{5k_B R}\]

    In terms of the specific heat at constant volume, the energy is given by

    \[M C_v T\]

    Returning to the equality,

    \[C_v = \frac{3GM}{5R T}\]

    Substituting \(T\),

    \[\boxed{C_v = \frac{k_B }{m} \approx 4.1 \cdot 10^3 \text{ J kg\(^{-1}\) K\(^{-1}\)}}\]

    Note that this value depends only on the mass of the particles that make up the gas cloud. For a hydrogen composition, the value is surprisingly close to that of water, \(C_{v, \text{water}} = 4.2\cdot 10^{3}\) J/kg K.
\end{pssolution*}
\end{pproblem}

\newpage

\newpage

\section{Optics and Telescopes}

\pts{3}
    \begin{pproblem} A Keplerian telescope has two converging lenses. The first has a radius of curvature of both faces equal to \(R = 2\text{ m}\), and the second has both radii of curvature equal to \(r = 0.5 \text{ m}\). Considering that both lenses have negligible thickness and are made of a material with refractive index \(n=1.4\), calculate the diameter and the magnification of the telescope, knowing that this telescope is an f/12.
    
    \end{pproblem}


\pts{5}
    \begin{pproblem}
        In this problem, we will become familiar with one of the most powerful tools in optics, \textit{geometric optics}. The goal of this tool is to model lenses and mirrors in the form of matrices. This is very useful for solving problems involving combinations of several lenses, and it is a method for solving geometric-optics problems (practically) without using geometry. For that, pay attention to the following definitions in the figure below.

        \begin{figure}[H]
            \centering
            \includegraphics[width=0.9\linewidth]{imagens/lentes1.png}
            \caption{Diagram 1}
        \end{figure}

        All dashed lines are parallel to the optical axis, represented by the horizontal arrow. \(P_k\) denotes the point where the light passes from one medium to another, and \(y_k\) is the vertical coordinate of the point \(P_k\). The subscripts \(i\) and \(t\) are used to refer to incident and transmitted, respectively. Thus, for example, \(\alpha_{i,k}\) is the angle that the incident light ray makes with the horizontal at the point \(P_k\), while the angle \(\alpha_k\) is the angle between the point \(P_k\) and the center of lens \(k\). For example, in the figure, the angle \(\alpha_1\) is represented by \(V_1\hat{C}P_1\).
        

        For this problem, we will assume that all angles of interest are small, so that \(\tan\theta \approx \sin\theta \approx \theta\).
        \begin{alternativas}
            \item Starting from Snell's law, find a relation between \(n_{i,1}\), \(n_{t,1}\), \(\alpha_1\), \(\alpha_{i,1}\), and \(\alpha_{t,1}\).
        
            \item Find a relation for \((1) \ n_{t,1}\alpha_{t,1}\) and \((2) \ y_{t,1}\). Leave your answers in terms of \(n_{i,1}\), \(\alpha_{i,1}\), \(y_{i,1}\), and \(\mathcal{D}_1\), where
            \[\mathcal{D}_1 = \frac{n_{t,1}-n_{i,1}}{R_1}\]

        Now for one more definition. Let the vector \(\mathbf{r}_{t,1}= (n_{t,1}\alpha_{t,1}, \ y_{t,1})\) and \(\mathbf{r}_{i,1}= (n_{i,1}\alpha_{i,1}, \ y_{i,1})\).
        
            \item The two equations found in the previous item can be written in matrix form, as
            \[\mathbf{r}_{t,1}=\mathcal{R}_1\mathbf{r}_{i,1}\]
            Find the \(2\times 2\) matrix equivalent to \(\mathcal{R}_1\).

        Our interest now is to find the relations between the points \(P_2\) and \(P_1\).

            \item Find an equation for \(n_1\alpha_{i,2}\) and \(y_{i,2}\). Using the reasoning of the previous item, leave your answer in the form
            \[\mathbf{r}_{i,2}=\Gamma_{2,1}\mathbf{r}_{t,1}\]
            where \(\Gamma_{2,1}\) is also a \(2\times 2\) matrix.

            \item We define the lens matrix, \(\mathcal{A}_{2,1}\), by the following equation:
            \[\mathbf{r}_{i,2} = \mathcal{A}_{2,1}\mathbf{r}_{i,1}\]
            so that

            \[\mathcal{A}_{2,1} = \begin{pmatrix}
                a_{11} & a_{12} \\
                a_{21} & a_{22} \\
            \end{pmatrix}\]

            Find \(\mathcal{A}_{2,1}\) explicitly.

            \item Among all the properties of the lens matrix, the most curious is that the term \(a_{12}\) equals \(-1/f\), where \(f\) is the focal length of the lens. Prove this result. (Hint: you will be able to write \(-a_{12}\) with a well-known equation from optics.)
        
        Now let us get to work and make practical use of matrix optics.
        
            \item Consider the following diagram:
                \begin{figure}[H]
                    \centering
                    \includegraphics[width=0.9\linewidth]{imagens/lentes2.png}
                    \caption{Diagram 2}
                \end{figure}
                 Find \(y_i\) in terms of \(y_0\) and the data in the figure.
                 
            \item If the lens is between two different media, \(n_{i,1} \ne n_{t,2}\), we have the following relation:
            \[a_{12} = -\frac{n_{i,1}}{f_0}=-\frac{n_{t,2}}{f_I}\]
                Using this, redo the previous item, considering that there is water with \(n_w = 4/3\) between the two lenses.
        
            \item A very common lens arrangement is the Tessar objective, found in many cameras because of its efficiency in reducing aberration and astigmatism. The arrangement has the following form:
                \begin{figure}[H]
                    \centering
                    \includegraphics[width=0.7\linewidth]{imagens/lentes3.png}
                    \caption{Tessar objective diagram}
                \end{figure}
                Find the equivalent lens matrix, \(\mathcal{A}_{1,7}\), in terms of \(\mathcal{R}_i\) and \(\Gamma_{i,j}\).
        \end{alternativas}
    \begin{pssolution*}{}{}
        \begin{alternativas}
            \item Snell's law tells us that \(n_{i,1}\theta_{i,1} = n_{t,1}\theta_{t,1}\), where \(\theta\) are the angles that the light makes with the normal. Note that we can write both angles \(\theta\) as
            \[\theta_{i,1} = \alpha_1 + \alpha_{i,1}, \ \ \theta _{t,1} = \alpha_1 + \alpha_{t,1}\]

            Therefore, the relation we are looking for is

            \[\boxed{n_{i,1}(\alpha_1+\alpha_{i,1}) = n_{t,1}(\alpha_1+\alpha_{t,1})}\]

            \item Substituting \(\alpha_1 = y_1/R_1\), we have
            
            \[n_{i,1}(y_1/R_1+\alpha_{i,1})=n_{t,1}(y_1/R_1 +\alpha_{t,1})\]

            Isolating \(n_{t,1}\alpha_{t,1}\),

            \[n_{t,1}\alpha_{t,1} = n_{i,1}\alpha_{i,1}+ \frac{y_1}{R_1}(n_{i,1}-n_{t,1})\]

            In terms of \(\mathcal{D}_1\),

            \[n_{t,1}\alpha_{t,1} = n_{i,1}\alpha_{i,1} - y_1\mathcal{D}_1\]

            Note that \(y\) does not change immediately after the light is transmitted. Therefore \(y_1\equiv y_{i,1}\equiv y_{t,1}\).

            \item In matrix form,
            
            \[\left(\begin{matrix}
                n_{t,1}\alpha_{t,1} \\
                y_{t,1}\\
            \end{matrix}\right) = \left(\begin{matrix}
                1 & - \mathcal{D}_1 \\
                0 & 1
            \end{matrix}\right)\left(\begin{matrix}
                n_{i,1}\alpha_{i,1} \\
                y_{i,1}\\
            \end{matrix}\right)\]

            Thus, our matrix \(\mathcal{R}_1\), also known as the refraction matrix, is given by

            \[\boxed{\mathcal{R}_1 = \left(\begin{matrix}
            1 & -\mathcal{D}_1 \\
            0 & 1 \\
            \end{matrix}\right)}\]

        \item First, let us find the height of point \(2\). Let \(d_{21}\) be the thickness of the lens. Using trigonometry, we obtain
        
        \[y_{i,2} = y_{t,1} + \alpha_{t,1}d_{21}\]

        We have that \(\alpha_{t,1}\) and \(\alpha_{i,2}\) are alternate interior angles. That is, \(\alpha_{t,1}=\alpha({i,2}\). Since the refractive index inside the lens is constant, we have
        \[\] 
        \end{alternativas}
\end{pssolution*}

\end{pproblem}

\pts{2}
\begin{pproblem}
    The wave theory of light shows that a perfect focus is not possible, because of diffraction effects associated with the finite aperture of the lens. This lack of a perfect focus prevents objects that are very close together from being distinguished. This problem can be studied from two different points of view:

    The wave theory of light predicts that a lens of diameter $D$ cannot focus a parallel beam of light of wavelength $\lambda$ into an angle smaller than the diffraction limit:
    \begin{equation}
        \theta_m \approx 1.22 \frac{\lambda}{D}
    \end{equation}

    Now consider a quantum approach: the photons that are focused by the lens. These photons are known to have passed somewhere within a radius of the center of the lens. The uncertainty in the position $x$ is associated with an uncertainty in the $x$ component of the photon's momentum. Consequently, a photon that, in the absence of this uncertainty, would have been brought onto the optical axis of the focal plane may now be deflected by an angle $\theta \ll 1$.

    Consider the de Broglie wavelength $\lambda = \frac{h}{p}$. Find a limit for $\theta$.

    \textbf{Hint:} Heisenberg's uncertainty principle relates the imprecision between measurements of the momentum and the position of a particle by
    \[\Delta p_x \Delta x = \frac{\hbar}{2}\]

\begin{pssolution*}{}{}
    The momentum of a photon is given by \(p = h/\lambda\), and we know that the photon has \(\Delta x = D\). Using Heisenberg's uncertainty principle,

    \[\Delta p_x = \frac{\hbar}{2D}\]

    Using \(\theta \approx \frac{\Delta p_x}{p}\), we obtain
    \[\boxed{\theta \approx\frac{\lambda\hbar}{2D h} = \frac{1}{4\pi}\frac{\lambda}{D}}\]
\end{pssolution*}
\end{pproblem}

\pts{3}
\begin{pproblem} (Magna notes)
    Prove that the quantity \(n_iR_i\sin \theta_i\), for spherical shells adjacent to one another, is an invariant, where \(n_i\) is the refractive index of the \(i\)-th spherical shell, \(\theta_i\) is the angle that the light ray makes with the normal to the \(i\)-th spherical shell, and \(R_i\) is the radius of the \(i\)-th spherical shell.
    
\end{pproblem}

\pts{4}
\begin{pproblem}
    The refractive index of a planet's atmosphere is given by

    \[n(h) = \frac{n_0}{1+\epsilon h}\]

    where \(n_0\) and \(\epsilon\) are constants.

    \begin{alternativas}
        \item A light ray reaches the atmosphere parallel to the surface, at a height \(h'\ll R\). What is the trajectory?
        \item Given the trajectory of the light ray, calculate the radius of the planet.
    
        \textbf{Hint:} You may find the following relation useful,

        \[\int \frac{t}{\sqrt{a-bt^2}}dt = -\frac{1}{b}\sqrt{a-bt^2}+C\]
    \end{alternativas}
\begin{pssolution*}{}{}
    \begin{alternativas}
        \item Defining \(ds = \sqrt{dx^2+dh^2}\) as the infinitesimal arc-length element, we have
        \[\sin\theta = \frac{dx}{ds}, \ \ \cos\theta = \frac{dh}{ds}\]

        Since the medium is continuous, we can write

        \[n(h)\sin\theta(h)=cte\]

        Thus,

        \[n\frac{dx}{ds} = K\]

        Substituting \(n\),

        \[\frac{dx}{ds} = \frac{K(1+\epsilon h)}{n_0}\]

        However, by the definition of \(s\), we have

        \[\left(\frac{dx}{ds}\right)^2 + \left(\frac{dh}{ds}\right)^2 = 1\]

        Substituting \(dx/ds\),

        \[\frac{dh}{ds} = \sqrt{1- \left(\frac{K(1+\epsilon h)}{n_0}\right)^2}\]

        Using the chain rule,

        \[\frac{dh}{dx} = \frac{dh/ds}{dx/ds} = \frac{\sqrt{1- \left(\frac{K(1+\epsilon h)}{n_0}\right)^2}}{\frac{K(1+\epsilon h)}{n_0}}\]

        Simplifying the expression,

        \[\frac{dh}{dx} = \frac{\sqrt{n_0^2 - K^2(1+\epsilon h)^2}}{K(1+\epsilon h)}\]

        To solve this ODE, let us introduce the variable \(u = 1+\epsilon h\), so that

        \[dh = \frac{du}{\epsilon}\]

        Thus,

        \[\frac{du}{dx} = \epsilon\frac{\sqrt{n_0^2 - K^2u^2}}{Ku}\]

        Separating the terms and integrating,

        \[\int \frac{u}{\sqrt{n_0^2 - K^2u^2}} du = \frac{\epsilon}{K}\int dx\]
    
        The integral on the left-hand side is the same one given in the problem, so

        \[-\frac{1}{K^2}\sqrt{n_0^2-K^2u^2} = \frac{\epsilon}{K}x + C\]

        Working with this expression,

        \[\sqrt{n_0^2-K^2u^2} = -\epsilon K x - K^2C\]

        \[n_0^2-K^2u^2 = (\epsilon K x + K^2C)^2\]

        Substituting \(u\),

        \[n_0^2-K^2(1+\epsilon h)^2 = \epsilon^2K^2x^2 + 2\epsilon K^3 C x + K^4C^2\]

        Rearranging the terms,

        \[\epsilon^2K^2x^2 + 2\epsilon K^3 C x + K^2(1+\epsilon h)^2 = n_0^2 - K^4C^2 \]

        \[\epsilon^2x^2 + 2\epsilon K C x + (1+\epsilon h)^2 = n_0^2/K^2 - K^2C^2 \]

        Dividing by \(\epsilon^2\),

        \[x^2 + \frac{2KC}{\epsilon}x + \frac{K^2C^2}{\epsilon^2} + \left(h+\frac{1}{\epsilon}\right)^2 = \left(\frac{n_0}{\epsilon K}\right)^2\]

        Collecting the perfect square,

        \[\boxed{\left(x + \frac{KC}{\epsilon}\right)^2+\left(h+\frac{1}{\epsilon}\right)^2 = \left(\frac{n_0}{\epsilon K}\right)^2}\]
   
        which is precisely the equation of a circle of radius \(\frac{n_0}{\epsilon K}\) and center \(C = (-\frac{KC}{\epsilon}, \ - \frac{1}{\epsilon})\).

        \item In the case of a ray that is tangent to the planet, that ray has \(\theta_0 = \pi/2\) and \(h_0=0\). Thus, using
        
        \[n(h)\sin\theta(h) = n_0\sin\theta_0 = K\]

        we obtain \(K = n_0\). Thus, the radius of the planet is given by the radius of the trajectory (the light ray is tangent to the surface) and, therefore,

        \[\boxed{R = \frac{1}{\epsilon}}\]
    \end{alternativas}
    
\end{pssolution*}
\end{pproblem}

\newpage

\section{Cosmology}

\pts{5}
\begin{pproblem}
    The Friedmann equation is one of the most important equations for the study of cosmology and of the Universe. The goal of this problem is to derive the fundamental equations of cosmology. First, let us begin with some definitions:
    \begin{enumerate}[label=\roman*)]
        \item Because of the expansion of the Universe, the distance between two points is given by
        \[r(t) = r_0 a(t)\]
        where \(a(t)\) is known as the scale factor and \(r_0\) is the distance measured at \(t = 0\). 
        
        \item The Universe follows the Robertson-Walker metric, which can be simplified to
        \[-c^2 dt^2 + dr^2 = 0\]
        \begin{alternativas}
            \item Find an expression for the comoving distance, \(r_0\), in integral form, from the definitions given above.
            \\
            \item The first Friedmann equation has the form
            \[H(t)^2 \equiv \left(\frac{\dot{a}}{a}\right)^2 = k_1 \varepsilon(t) + \frac{C}{r_0^2 a^2}\]
            Find the value of the constant \(k_1\) for a spherical universe, in terms of fundamental constants. In the equation above, \(\varepsilon(t)\) is the energy density of the universe, \(C\) is a constant related to its energy, and \(a\) is the scale factor.

            \textbf{Hint:} You can obtain this equation either from energy conservation or by using Newton's second law.

            \item Repeat the previous item for a cylindrical universe expanding radially. The equation you find should have the form

            \[\left(\frac{\dot{a}}{a}\right)^2 = f(a) \varepsilon(t) + \frac{C'}{r_0^2 a^2}\]

            Find the function \(f(a)\).

            \item Returning now to a ``normal universe''. In cosmology, we define the critical energy density as the energy density of a universe in which \(C = 0\). Find an expression for the critical density.
            \item Rewrite the first Friedmann equation in terms of
            \[\Omega(t) = \frac{\varepsilon(t)}{\varepsilon_c(t)}\]
            where \(\varepsilon_c(t)\) is the critical energy density.
        \end{alternativas}
    \end{enumerate}
\begin{pssolution*}{}{ }
    \begin{alternativas}
        \item Isolating \(d_r\),
        \[d_r = c dt\]

        But, by definition, \(r(t) = r_0 a(t) \rightarrow d_r = r_0 da\). Substituting,

        \[dr_0 = c\frac{dt}{a(t)}\]

        Integrating both sides,

        \[r_0 = c\int \frac{dt}{a(t)}\]

        \item \textbf{Method 1: Energy conservation}
        Consider a particle of mass \(m\) located at a point in the universe. Its energy relative to the center of the universe is given by

        \[E = \frac{mv^2}{2} - \frac{GM(r)m}{r}\]

        where \(v = \dot{r} = r_0 \dot{a}\) and \(M(r) = \frac{4}{3} \pi r^3 \rho(t)\). Dividing the equation by \(m/2\) and applying the substitutions, we have

        \[\frac{2E}{m} = r_0^2 \dot{a}^2 - \frac{G}{r_0 a} \frac{8\pi (r_0 a)^3 \rho(t)}{3}\]

        Simplifying,

        \[\frac{2E}{m} = r_0^2 \dot{a}^2 - \frac{8\pi G r_0^2 a^2 \rho(t)}{3}\]

        Multiplying both sides by \(\frac{1}{r_0^2 a^2}\), we obtain

        \[\frac{2E}{mr_0^2 a^2} = \left(\frac{\dot{a}}{a}\right)^2 - \frac{8\pi G \rho(t)}{3}\]

        Rearranging the equation and using the mass-energy equivalence, \(m = E/c^2 \rightarrow \rho(t) = \varepsilon(t)/c^2\), we have

        \[\left(\frac{\dot{a}}{a}\right)^2 = \frac{8\pi G}{3 c^2} \varepsilon(t) + \frac{2E}{mr_0^2 a^2}\]

        From which we can conclude that

        \[\boxed{k_1 = \frac{8\pi G}{3 c^2}}\]


        \textbf{Method 2: Newton's second law}

        Writing Newton's second law for a body relative to the center of the universe,

        \[-\frac{GM(r)m}{r^2} = m \frac{d^2r}{dt^2}\]

        Note that, as it expands, the universe does not create mass, so the mass contained in a shell \(r\) is constant over time. Using this property, we have

        \[\frac{d^2r}{dt^2} = -\frac{GM(r)}{r^2}\]

        Multiplying both sides by \(dr/dt\),

        \[\frac{dr}{dt} \frac{d^2r}{dt^2} = -\frac{GM(r)}{r^2} \frac{dr}{dt}\]

        Rewriting,

        \[\dot{r} \frac{d\dot{r}}{dt} = -\frac{GM(r)}{r^2} \frac{dr}{dt}\]

        Multiplying both sides by \(dt\),

        \[\dot{r} d\dot{r} = -GM(r) \frac{dr}{r^2}\]

        Substituting \(r = r_0 a\),

        \[r_0^2 \dot{a} d\dot{a} = -\frac{GM(r)}{r_0} \frac{da}{a^2}\]

        Integrating both sides,

        \[r_0^2 \int \dot{a} da = -\frac{GM(r)}{r_0^2} \int \frac{da}{a^2}\]

        \[\frac{r_0^2 \dot{a}^2}{2} = \frac{GM(r)}{r_0 a} + A\]

        where \(A\) is a constant of integration. Substituting \(M(r)\),

        \[\frac{r_0^2 \dot{a}^2}{2} = \frac{4\pi G r_0^2 a^2}{3c^2} \varepsilon(t) + A\]

        Multiplying the expression by \(2/a^2 r_0^2\), we have

        \[\left(\frac{\dot{a}}{a}\right)^2 = \frac{8\pi G}{3c^2} \varepsilon(t) + \frac{2A}{r_0^2 a^2}\]

        From which we can likewise conclude that

        \[\boxed{k_1 = \frac{8\pi G}{3c^2}}\]

        \item Using Gauss's law,

        \[\oint \mathbf{g} \cdot d\mathbf{S} = -4\pi G M(r)\]

        where

        \[\oint \mathbf{g} \cdot d\mathbf{S} = -g(2\pi r L)\]

        for cylindrical symmetry. Substituting into the equation,

        \[g = \frac{4\pi G M(r)}{2\pi r L} = \frac{2G M(r)}{r L}\]

        Using the energy-conservation method, we first have to calculate the potential energy,

        \[U = -\int F dr = -\int \frac{2GM(r)m}{rL} dr = -\int \frac{2GM(r)m}{L} \frac{da}{a} - \frac{2GM(r)m}{L} \ln(a)\]

        By energy conservation,

        \[E = \frac{m r_0^2 \dot{a}^2}{2} - \frac{2GM(r)m}{L} \ln(a)\]

        Substituting \(M(r) = \pi r_0^2 a^2 L \varepsilon(t)/c^2\),

        \[E = \frac{m r_0^2 \dot{a}^2}{2} - \frac{2\pi r_0^2 a^2 G m \ln(a)}{c^2} \varepsilon(t)\]

        Rearranging,

        \[\left(\frac{\dot{a}}{a}\right)^2 = \frac{4\pi G \ln(a)}{c^2} \varepsilon(t) + \frac{2E}{mr_0^2 a^2}\]

        From which we obtain

        \[\boxed{f(a) = \frac{4\pi G}{c^2} \ln(a)}\]

        \item The Friedmann equation for \(C = 0\) reduces to

        \[H(t)^2 = \frac{8\pi G}{3c^2} \varepsilon_c(t)\]

        Isolating \(\varepsilon_c(t)\),

        \[\boxed{\varepsilon_c(t) = \frac{3c^2}{8\pi G} H^2(t)}\]

        By definition,

        \[\varepsilon(t) = \Omega (t) \varepsilon_c(t)\]

        Substituting into the equation,

        \[H^2(t) = \frac{8\pi G}{3c^2} \Omega(t) \varepsilon_c(t) + \frac{C}{r_0^2 a^2}\]

        Substituting \(\varepsilon_c(t)\),

        \[H^2(t) = H^2(t) \Omega(t) + \frac{C}{r_0^2 a^2}\]

        Solving for \(H^2(t)\),

        \[\boxed{H^2(t) = \frac{C}{r_0^2 a^2 (1 - \Omega(t))}}\]
    \end{alternativas}
\end{pssolution*}
\end{pproblem}


\pts{3}
\begin{pproblem}(Adapted from NAO 2019) 
    Consider a flat universe in which the gravitational constant is no longer constant and is instead defined by
    \[G(a) = G_0 f(a)\]
    where \(f(a)\) is a function of the scale factor.
    \begin{alternativas}
        \item What would the Friedmann equation be in this universe? Assume that the universe is flat, \(C = 0\), and that it is composed only of baryonic (``bright'') matter. Leave your answer in terms of \(H_0\), \(f(a)\), \(a\), and \(\Omega_{m,0}\), where \(H_0\) is the value of the Hubble constant at the present time, and \(\Omega_{m,0}\) is the density parameter,
        
        In the case \(f(a) = e^{b(a-1)}\), where \(b=2\).
        
        \item Estimate the age of this universe, assuming that it consists only of baryonic matter (``bright'' matter).
                
        \item What is the behavior of the age when \(t \rightarrow \infty\)?
    \end{alternativas}

    You may find the following relations useful:
    \begin{paracol}{2}
        \[\int_0^\infty x^2 e^{-x^2} = \frac{\sqrt{\pi}}{4}\]

        \switchcolumn

        \[\int_0^1 x^2 e^{-x^2} \approx 0.189\]
    \end{paracol}

\begin{pssolution*}{}{ }
    \begin{alternativas}
        \item For a flat universe containing only matter, the Friedmann equation has the form
        \[H(t)^2 = H_0^2 \Omega_m\]

        where \(\Omega_m = \frac{\rho_m(t)}{\rho_{c, 0}}\) and

        \[\rho_c = \frac{3 H_0^2}{8\pi G}\]

        For baryonic matter,

        \[\rho_m(t) = \rho_0 a(t)^{-3}\]

        Then

        \[\Omega_m = \Omega_{m,0} f(a) a^{-3}\]

        For a universe composed only of baryonic matter, we have \(\Omega_{m,0} = 1\). Thus, the Friedmann equation becomes

        \[\boxed{H(a)^2 = H_0^2 \Omega_{m,0} f(a) a^{-3} \equiv H_0^2 f(a) a^{-3}}\]

        \item Using \(H = \frac{\dot{a}}{a}\), we have

        \[t = \int dt = \int da \frac{dt}{da} = \int \frac{da}{\dot{a}}\]

        Multiplying by \(a/a\),

        \[t = \int \frac{a}{\dot{a}a} da = \int \frac{da}{H(a) a}\]

        Substituting the value of \(H(a)\) found in the previous item,

        \[t = \frac{1}{H_0^2} \int \frac{da}{\sqrt{f(a) a^{-3}}} = \frac{1}{H_0^2} \int a^{3/2} e^{-b(a-1)/2} da = \frac{e^{b/2}}{H_0^2} \int a^{3/2} e^{-b a/2} da\]
        
        Using a substitution of the form \(x = \sqrt{ba/2}\), we have

        \[t = \frac{4\sqrt{2} e^{b/2}}{b^{3/2} H_0} \int x^2 e^{-x^2}\]

        The limits of integration run from \(a=0\) (the beginning of the universe) to \(a=1\) (the present). Since we substituted for \(x\),

        \[x = \sqrt{\frac{ab}{2}}\]

        \[t = \frac{4\sqrt{2} e^{b/2}}{b^{3/2} H_0} \int_0^{\sqrt{b/2}} x^2 e^{-x^2}\]

        Substituting \(b=2\),

        \[t = \frac{4\sqrt{2} e}{2^{3/2} H_0} \int_0^1 x^2 e^{-x^2}\]

        Substituting the values and using the integral provided, we obtain

        \[\boxed{t \approx 15 \text{ Gyr}}\]

        which is close to our universe!!

        \item Changing the limits of the integral to an infinite time, \(a(t) \rightarrow \infty\),

        \[\boxed{t_\infty = \frac{4\sqrt{2} e}{2^{3/2} H_0} \int_0^{\infty} x^2 e^{-x^2} = \frac{\sqrt{2\pi} e}{2^{3/2} H_0} \approx 34.1 \text{ Gyr}}\]

    \end{alternativas}
\end{pssolution*}
\end{pproblem}



\pts{5}
\begin{pproblem}(Problem set 8 - 2021)
    The Friedmann equation is given by
    \begin{equation}
        H^2 = \left( \frac{\dot{a}}{a} \right)^2 = \frac{8\pi G \varepsilon}{3c^2} - \frac{k c^2}{a^2},
    \end{equation}

    where \(a\) is the scale factor at time \(t\), \(\varepsilon\) is the energy density at time \(t\), and \(k\) is the parameter that characterizes the geometry of the universe and may take any real value. Considering a universe composed only of non-relativistic baryonic matter and solving this nonlinear differential equation for \(k > 0\), one obtains the following solutions in terms of the parameter \(\theta \in [0, 2\pi]\):

    \begin{align}
        a(\theta) &= \frac{4\pi G \varepsilon_0}{3k c^4} \left( 1 - \cos\theta \right), \\
        t(\theta) &= \frac{4\pi G \varepsilon_0}{3k^{3/2}c^5} \left( \theta - \sin\theta \right).
    \end{align}

    Consider a universe with \(\Omega_0 = 4\) and \(H_0 = 67.4 \ \text{km/s/Mpc}\).

    \begin{alternativas}
        \item Starting from the Friedmann equation, show that \(k c^2 = H^2 \left( 1 - \Omega \right) a^2\). Finally, rewrite the parametric equations for \(a\) and \(t\) in terms of the present density parameter \(\Omega_0\), the present Hubble constant \(H_0\), and the parameter \(\theta\). Do not substitute their respective numerical values.
        
        \item Find the age \(t_0\) of this universe in terms of \(H_0\), and then in billions of years.
        
        \item The so-called \textit{lookback time}, \(\Delta t_L\), is how far in the past the universe had a given scale factor \(a\). What is \(\Delta t_L\), in billions of years, for when the size of the universe was \(1/3\) of what it is at present?
        
        \item Determine \(\theta_n\), and then \(t_n\), for which \(H = 0\).
    \end{alternativas}

\begin{pssolution*}{}{}
    \begin{alternativas}
        \item 
    \end{alternativas}
    
\end{pssolution*}
\end{pproblem}

\pts{5} 
\begin{pproblem}(Problem set 6 - 2022)
aaa
\end{pproblem}

\newpage

\section{Miscellaneous}
\pts{4}
\begin{pproblem} \clock{1}{0}
    Dudu Leiteiro was observing the sky in the interior of his farm in Mato Grosso do Sul and observed the binary system formed by the stars Iaum and Sezenem. Dudu observed the stars and collected the following data, with an interval of 6 months between them:
    \\
    \begin{center}
        \begin{tabular}{|c|c|c|}
            \hline % Top horizontal line
            Measurement & Iaum & Sezenem   \\ 
            \hline
            Right ascension \(1\) & \(4^h19^m53.91^s\) & \(4^h19^m53.078^s\) \\
            Right ascension \(2\) & \(4^h19^m53.92^s\) & \(4^h19^m53.077^s\) \\
            Declination \(1\) & -\(13^\circ \ 33' \ 45.28''\) & -\(13^\circ \ 33' \ 47.78''\) \\
            Declination \(2\) & -\(13^\circ \ 33' \ 45.20''\) & -\(13^\circ \ 33' \ 48.03''\) \\
            \hline
        \end{tabular} 
\end{center}
    With these data, you and Dudu Leiteiro together will analyze properties of this binary system. The first important step is to find the coordinates of the center of mass of the system, denoted by the subscript \(_{CM}\). You remember that, in one of the classes on the study of binaries, your professor, LuCav, taught you that
    \[
    \delta_{CM} = \delta_A + \Delta\delta_A
    \]

    \[
    \alpha_{CM} = \alpha_A + \Delta\alpha_A
    \]

    where \(\delta_A\) and \(\alpha_A\) are the coordinates of one of the stars in the binary and \(\Delta\) represents the difference in coordinates between the star and the center of mass. We have the following relation:

    \[
    \frac{a_A}{a} = \frac{\Delta\delta_A}{\Delta\delta_{CM}} = \frac{\Delta\alpha_A}{\Delta\alpha_{CM}}
    \]

    where \(a_A\) is the distance from star \(A\) to the \(CM\), \(a\) is the semi-major axis of the orbit, and \(\Delta x_{CM}\) is the variation of the coordinates of the \(CM\). 
    
    \textbf{a)} Given that Iaum has \(3/2\) the mass of Sezenem, calculate \(\Delta\delta_{CM}\) and \(\Delta\alpha_{CM}\).
    
    \textbf{b)} With this, and considering that the angular variations are small enough for spherical triangles to be treated as plane, find the parallax of the system and its distance to Earth.

    \textbf{c)} Considering the mass of Iaum, \(M_I = 4.9 M_\odot\), and the period of the system equal to \(P = 29.01\) years, calculate the largest redshift coming from Sezenem, given that both orbits are circular and Sezenem has a tangential velocity of \(\mu = 1509''/\text{year}\).

\end{pproblem}

\pts{3}
\begin{pproblem} Marisso was tired of not being able to find the position of a star precisely with his telescope and decided to investigate the effects that could be causing this apparent error. After reading some papers, he discovered \(2\) main effects that make an object appear to lie at an angle \(\Delta \theta_i\), offset from its original position. They are \textit{parallax} and \textit{stellar aberration}. In this problem, your goal is to help Marrisso understand why these effects occur.
    \begin{alternativas}
        \item Parallax is the most basic of them, and it occurs because of the change in the position of the Earth over the course of the year. Consider a star located so that the line Sun--star is perpendicular to the plane of the Earth's orbit. Draw a diagram of the situation and, taking the radius of the Earth's orbit as \(r\) and the distance to the star as \(d\), obtain a formula for \(\Delta \theta_p\) caused by parallax.
        
        \item Now suppose that the line Sun--star makes an arbitrary angle \(\phi\) with the Earth's orbit. How does your answer change?
        
        Stellar aberration, in turn, comes from relativistic effects, which are explored below.

        \item Consider a frame \(S'\) moving with velocity \(v\hat{\mathbf{x}}\) relative to the frame \(S\). How are the coordinates \((x', t')\) related to the coordinates \((x,t)\)? Leave your answers in terms of \(\gamma\).
        
        \item Suppose there is a radiation emitter in the frame \(S'\) and that it emits light at an angle \(\alpha'\) relative to the \(x\) axis. In the frame \(S\), the device will appear to emit light at an angle \(\alpha\). Prove, using the Lorentz transformations, that the relation between \(\alpha\) and \(\alpha'\) is given by
        
        \[\cos\alpha = \frac{\cos\alpha' + v/c}{1+(\cos\alpha') v/c}\]

        \item Repeat the previous item, but prove it using relativistic velocity addition.

        \item Considering that the line Sun--star is perpendicular to the plane of the Earth's orbit, and that the Earth moves with velocity \(v\), find an expression for the shift \(\Delta\theta_A\) caused by stellar aberration.
    
        \item Which of these effects do you think is more significant, parallax or stellar aberration?
    \end{alternativas}

    \begin{pssolution*}{}{}
        \begin{alternativas}
            \item % too lazy to fetch GeoGebra
            
            \item % lazy GeoGebra
                
            \item The Lorentz transformations (using \(c=1\) ) are given by
            \[x = \gamma (x' + vt'), \ \ t = \gamma(t' + vx')\]

            \item Note that in the frame \(S\), we have \(\cos \alpha = \frac{x}{t}\). In the frame \(S'\) we have
            
            \[x' = t'\cos\alpha'\]

            Using the Lorentz transformations,

            \[t = \gamma t'(1+v\cos\alpha') \]

            \[x = \gamma t'(\cos\alpha'+ v)\]

            Using \(\cos\alpha = x/t\), we have

            \[\cos\alpha = \frac{\cos\alpha' + v}{1+ v\cos\alpha'}\]

            ``Putting the \(c\)'s back'', we have

            \[\boxed{\cos\alpha = \frac{\cos\alpha' + v/c}{1+(v/c)\cos\alpha'}}\]


            \item In the frame \(S'\), the velocity along \(x'\) is given by \(v_x' = c\cos\alpha'\). In the frame \(S\), we have
            
            \[v_x = "v_x'+v" = \frac{v_x'+v}{1+(v_x'v)/c^2}\]

            Substituting \(v_x'\) and using \(\cos\alpha = v_x/c\), we have

            \[\boxed{\cos\alpha =\frac{\cos\alpha' + v/c}{1+(v/c)\cos\alpha'}} \]

            \item In the Earth's frame (which is moving and therefore corresponds to the frame \(S'\)), we have \(\alpha' = \pi/2\). Plugging this into the answer, we obtain
            
            \[\cos\alpha = v/c \rightarrow \sin\alpha = \sqrt{1-(v/c)^2} = \frac{1}{\gamma}\]

            Using the small-angle approximation, \(\boxed{\alpha \approx 1/\gamma}\).

            \item While the Sun--Earth distance is on the order of a few light-minutes, the smallest distance between the Earth and another star is on the order of light-years, so \(\theta \sim 10^{-6}\). The speed of the Earth is approximately \(30\) km/s, so \(v/c \sim 10^{-4}\). Therefore the effects of aberration are more apparent than those of parallax.


        \end{alternativas}

        \end{pssolution*}

\end{pproblem}
\end{document}
